\par\textit{LoRA trains low-rank update matrices while freezing base weights for parameter-efficient adaptation. Tradeoff: Memory and storage fall, while rank and target-module choices limit expressiveness. Watch for: Merging incompatible adapters without evaluation creates interference.}
\paragraph{MCQ-1273 [hard]} Which failure mode should an interviewer expect you to identify for LoRA?
\begin{enumerate}[label=\Alph*.]
\item Omitting the square-root head-dimension scale saturates softmax gradients.
\item Training on conflicting instructions without provenance makes behavior unstable.
\item Merging incompatible adapters without evaluation creates interference.
\item Extending context without validating the position scheme can collapse quality.
\end{enumerate}
\noindent\textbf{Answer.} Merging incompatible adapters without evaluation creates interference.
\par\textit{LoRA trains low-rank update matrices while freezing base weights for parameter-efficient adaptation. Tradeoff: Memory and storage fall, while rank and target-module choices limit expressiveness. Watch for: Merging incompatible adapters without evaluation creates interference.}
\paragraph{MCQ-1274 [medium]} A design review is specifically concerned with this risk: Merging incompatible adapters without evaluation creates interference. Which topic should be reviewed first?
\begin{enumerate}[label=\Alph*.]
\item Self-attention math
\item Positional encoding
\item SFT
\item LoRA
\end{enumerate}
\noindent\textbf{Answer.} LoRA
\par\textit{LoRA trains low-rank update matrices while freezing base weights for parameter-efficient adaptation. Tradeoff: Memory and storage fall, while rank and target-module choices limit expressiveness. Watch for: Merging incompatible adapters without evaluation creates interference.}
\paragraph{MCQ-1275 [hard]} Which design note most accurately balances the benefits and costs of LoRA?
\begin{enumerate}[label=\Alph*.]
\item Relative methods extrapolate differently, while implementation and scaling choices affect long context.
\item Every token can interact directly, while standard sequence cost grows quadratically.
\item It is direct and stable but can overfit style and inherit annotation errors.
\item Memory and storage fall, while rank and target-module choices limit expressiveness.
\end{enumerate}
\noindent\textbf{Answer.} Memory and storage fall, while rank and target-module choices limit expressiveness.
\par\textit{LoRA trains low-rank update matrices while freezing base weights for parameter-efficient adaptation. Tradeoff: Memory and storage fall, while rank and target-module choices limit expressiveness. Watch for: Merging incompatible adapters without evaluation creates interference.}
\paragraph{MCQ-1276 [medium]} Which explanation would be strongest in a system-design interview about LoRA?
\begin{enumerate}[label=\Alph*.]
\item Attention computes scaled query-key scores, applies masking and softmax, then mixes value vectors.
\item Position signals break permutation invariance using absolute, relative, rotary, or bias-based schemes.
\item Supervised fine-tuning teaches instruction-response behavior from curated demonstrations.
\item LoRA trains low-rank update matrices while freezing base weights for parameter-efficient adaptation.
\end{enumerate}
\noindent\textbf{Answer.} LoRA trains low-rank update matrices while freezing base weights for parameter-efficient adaptation.
\par\textit{LoRA trains low-rank update matrices while freezing base weights for parameter-efficient adaptation. Tradeoff: Memory and storage fall, while rank and target-module choices limit expressiveness. Watch for: Merging incompatible adapters without evaluation creates interference.}
\paragraph{MCQ-1277 [hard]} Which production warning is most directly associated with LoRA?
\begin{enumerate}[label=\Alph*.]
\item Merging incompatible adapters without evaluation creates interference.
\item Extending context without validating the position scheme can collapse quality.
\item Training on conflicting instructions without provenance makes behavior unstable.
\item Omitting the square-root head-dimension scale saturates softmax gradients.
\end{enumerate}
\noindent\textbf{Answer.} Merging incompatible adapters without evaluation creates interference.
\par\textit{LoRA trains low-rank update matrices while freezing base weights for parameter-efficient adaptation. Tradeoff: Memory and storage fall, while rank and target-module choices limit expressiveness. Watch for: Merging incompatible adapters without evaluation creates interference.}
\paragraph{FC-1278 [medium]} Define LoRA in interview-ready language.
\noindent\textbf{Answer.} LoRA trains low-rank update matrices while freezing base weights for parameter-efficient adaptation.
\par\textit{LoRA trains low-rank update matrices while freezing base weights for parameter-efficient adaptation.}
\paragraph{FC-1279 [hard]} State the main tradeoff and production pitfall for LoRA.
\noindent\textbf{Answer.} Tradeoff: Memory and storage fall, while rank and target-module choices limit expressiveness. Pitfall: Merging incompatible adapters without evaluation creates interference.
\par\textit{Tradeoff: Memory and storage fall, while rank and target-module choices limit expressiveness. Pitfall: Merging incompatible adapters without evaluation creates interference.}
\paragraph{LA-1280 [hard]} Explain LoRA from first principles, then show how you would evaluate and operate it in production.
\noindent\textbf{Answer.} Start with the mechanism: LoRA trains low-rank update matrices while freezing base weights for parameter-efficient adaptation. Make the tradeoff explicit: Memory and storage fall, while rank and target-module choices limit expressiveness. Define workload-specific offline metrics and a latency/cost budget, test important slices, instrument the critical path, and create a rollback criterion. The production review must explicitly guard against this failure: Merging incompatible adapters without evaluation creates interference.
\par\textit{A strong answer connects mechanism, measurable tradeoffs, failure isolation, observability, and rollback rather than listing product features.}
\paragraph{MCQ-1281 [medium]} Which statement best describes QLoRA?
\begin{enumerate}[label=\Alph*.]
\item Rotary position embeddings rotate query and key coordinates so dot products encode relative offsets.
\item QLoRA backpropagates through a frozen quantized base model into LoRA adapters using memory-saving quantization techniques.
\item Attention computes scaled query-key scores, applies masking and softmax, then mixes value vectors.
\item Cross-entropy is negative log likelihood; perplexity exponentiates average token loss.
\end{enumerate}
\noindent\textbf{Answer.} QLoRA backpropagates through a frozen quantized base model into LoRA adapters using memory-saving quantization techniques.
\par\textit{QLoRA backpropagates through a frozen quantized base model into LoRA adapters using memory-saving quantization techniques. Tradeoff: Fine-tuning large models becomes cheaper, while kernels and quantization settings affect stability. Watch for: Confusing training quantization with final serving quantization produces wrong expectations.}
\paragraph{MCQ-1282 [hard]} What is the central engineering tradeoff of QLoRA?
\begin{enumerate}[label=\Alph*.]
\item Every token can interact directly, while standard sequence cost grows quadratically.
\item It is stable for model comparison on the same tokenization and data, not across arbitrary setups.
\item RoPE integrates cleanly with attention but long-context scaling needs careful frequency treatment.
\item Fine-tuning large models becomes cheaper, while kernels and quantization settings affect stability.
\end{enumerate}
\noindent\textbf{Answer.} Fine-tuning large models becomes cheaper, while kernels and quantization settings affect stability.
\par\textit{QLoRA backpropagates through a frozen quantized base model into LoRA adapters using memory-saving quantization techniques. Tradeoff: Fine-tuning large models becomes cheaper, while kernels and quantization settings affect stability. Watch for: Confusing training quantization with final serving quantization produces wrong expectations.}
\paragraph{MCQ-1283 [hard]} Which failure mode should an interviewer expect you to identify for QLoRA?
\begin{enumerate}[label=\Alph*.]
\item Naively increasing the maximum length does not preserve trained frequency behavior.
\item Confusing training quantization with final serving quantization produces wrong expectations.
\item Omitting the square-root head-dimension scale saturates softmax gradients.
\item Comparing perplexity across vocabularies or tokenizers is misleading.
\end{enumerate}
\noindent\textbf{Answer.} Confusing training quantization with final serving quantization produces wrong expectations.
\par\textit{QLoRA backpropagates through a frozen quantized base model into LoRA adapters using memory-saving quantization techniques. Tradeoff: Fine-tuning large models becomes cheaper, while kernels and quantization settings affect stability. Watch for: Confusing training quantization with final serving quantization produces wrong expectations.}
\paragraph{MCQ-1284 [medium]} A design review is specifically concerned with this risk: Confusing training quantization with final serving quantization produces wrong expectations. Which topic should be reviewed first?
\begin{enumerate}[label=\Alph*.]
\item RoPE
\item QLoRA
\item Self-attention math
\item Cross-entropy and perplexity
\end{enumerate}
\noindent\textbf{Answer.} QLoRA
\par\textit{QLoRA backpropagates through a frozen quantized base model into LoRA adapters using memory-saving quantization techniques. Tradeoff: Fine-tuning large models becomes cheaper, while kernels and quantization settings affect stability. Watch for: Confusing training quantization with final serving quantization produces wrong expectations.}
\paragraph{MCQ-1285 [hard]} Which design note most accurately balances the benefits and costs of QLoRA?
\begin{enumerate}[label=\Alph*.]
\item Fine-tuning large models becomes cheaper, while kernels and quantization settings affect stability.
\item It is stable for model comparison on the same tokenization and data, not across arbitrary setups.
\item Every token can interact directly, while standard sequence cost grows quadratically.
\item RoPE integrates cleanly with attention but long-context scaling needs careful frequency treatment.
\end{enumerate}
\noindent\textbf{Answer.} Fine-tuning large models becomes cheaper, while kernels and quantization settings affect stability.
\par\textit{QLoRA backpropagates through a frozen quantized base model into LoRA adapters using memory-saving quantization techniques. Tradeoff: Fine-tuning large models becomes cheaper, while kernels and quantization settings affect stability. Watch for: Confusing training quantization with final serving quantization produces wrong expectations.}
\paragraph{MCQ-1286 [medium]} Which explanation would be strongest in a system-design interview about QLoRA?
\begin{enumerate}[label=\Alph*.]
\item Cross-entropy is negative log likelihood; perplexity exponentiates average token loss.
\item Rotary position embeddings rotate query and key coordinates so dot products encode relative offsets.
\item Attention computes scaled query-key scores, applies masking and softmax, then mixes value vectors.
\item QLoRA backpropagates through a frozen quantized base model into LoRA adapters using memory-saving quantization techniques.
\end{enumerate}
\noindent\textbf{Answer.} QLoRA backpropagates through a frozen quantized base model into LoRA adapters using memory-saving quantization techniques.
\par\textit{QLoRA backpropagates through a frozen quantized base model into LoRA adapters using memory-saving quantization techniques. Tradeoff: Fine-tuning large models becomes cheaper, while kernels and quantization settings affect stability. Watch for: Confusing training quantization with final serving quantization produces wrong expectations.}
\paragraph{MCQ-1287 [hard]} Which production warning is most directly associated with QLoRA?
\begin{enumerate}[label=\Alph*.]
\item Naively increasing the maximum length does not preserve trained frequency behavior.
\item Omitting the square-root head-dimension scale saturates softmax gradients.
\item Confusing training quantization with final serving quantization produces wrong expectations.
\item Comparing perplexity across vocabularies or tokenizers is misleading.
\end{enumerate}
\noindent\textbf{Answer.} Confusing training quantization with final serving quantization produces wrong expectations.
\par\textit{QLoRA backpropagates through a frozen quantized base model into LoRA adapters using memory-saving quantization techniques. Tradeoff: Fine-tuning large models becomes cheaper, while kernels and quantization settings affect stability. Watch for: Confusing training quantization with final serving quantization produces wrong expectations.}
\paragraph{FC-1288 [medium]} Define QLoRA in interview-ready language.
\noindent\textbf{Answer.} QLoRA backpropagates through a frozen quantized base model into LoRA adapters using memory-saving quantization techniques.
\par\textit{QLoRA backpropagates through a frozen quantized base model into LoRA adapters using memory-saving quantization techniques.}
\paragraph{FC-1289 [hard]} State the main tradeoff and production pitfall for QLoRA.
\noindent\textbf{Answer.} Tradeoff: Fine-tuning large models becomes cheaper, while kernels and quantization settings affect stability. Pitfall: Confusing training quantization with final serving quantization produces wrong expectations.
\par\textit{Tradeoff: Fine-tuning large models becomes cheaper, while kernels and quantization settings affect stability. Pitfall: Confusing training quantization with final serving quantization produces wrong expectations.}
\paragraph{LA-1290 [hard]} Explain QLoRA from first principles, then show how you would evaluate and operate it in production.
\noindent\textbf{Answer.} Start with the mechanism: QLoRA backpropagates through a frozen quantized base model into LoRA adapters using memory-saving quantization techniques. Make the tradeoff explicit: Fine-tuning large models becomes cheaper, while kernels and quantization settings affect stability. Define workload-specific offline metrics and a latency/cost budget, test important slices, instrument the critical path, and create a rollback criterion. The production review must explicitly guard against this failure: Confusing training quantization with final serving quantization produces wrong expectations.
\par\textit{A strong answer connects mechanism, measurable tradeoffs, failure isolation, observability, and rollback rather than listing product features.}
\paragraph{MCQ-1291 [medium]} Which statement best describes Knowledge distillation?
\begin{enumerate}[label=\Alph*.]
\item Multiple heads project attention into separate subspaces before concatenation and output projection.
\item QLoRA backpropagates through a frozen quantized base model into LoRA adapters using memory-saving quantization techniques.
\item LoRA trains low-rank update matrices while freezing base weights for parameter-efficient adaptation.
\item A student learns from teacher probabilities, generated data, rationales, or intermediate representations.
\end{enumerate}
\noindent\textbf{Answer.} A student learns from teacher probabilities, generated data, rationales, or intermediate representations.
\par\textit{A student learns from teacher probabilities, generated data, rationales, or intermediate representations. Tradeoff: Serving cost falls but rare capabilities and calibration may be lost. Watch for: Distilling only easy teacher outputs creates a deceptively strong narrow model.}
\paragraph{MCQ-1292 [hard]} What is the central engineering tradeoff of Knowledge distillation?
\begin{enumerate}[label=\Alph*.]
\item Fine-tuning large models becomes cheaper, while kernels and quantization settings affect stability.
\item Memory and storage fall, while rank and target-module choices limit expressiveness.
\item Serving cost falls but rare capabilities and calibration may be lost.
\item Diverse interactions improve capacity but add projections and KV memory.
\end{enumerate}
\noindent\textbf{Answer.} Serving cost falls but rare capabilities and calibration may be lost.
\par\textit{A student learns from teacher probabilities, generated data, rationales, or intermediate representations. Tradeoff: Serving cost falls but rare capabilities and calibration may be lost. Watch for: Distilling only easy teacher outputs creates a deceptively strong narrow model.}
\paragraph{MCQ-1293 [hard]} Which failure mode should an interviewer expect you to identify for Knowledge distillation?
\begin{enumerate}[label=\Alph*.]
\item Merging incompatible adapters without evaluation creates interference.
\item Distilling only easy teacher outputs creates a deceptively strong narrow model.
\item Assuming heads are independently interpretable can overstate mechanistic conclusions.
\item Confusing training quantization with final serving quantization produces wrong expectations.
\end{enumerate}
\noindent\textbf{Answer.} Distilling only easy teacher outputs creates a deceptively strong narrow model.
\par\textit{A student learns from teacher probabilities, generated data, rationales, or intermediate representations. Tradeoff: Serving cost falls but rare capabilities and calibration may be lost. Watch for: Distilling only easy teacher outputs creates a deceptively strong narrow model.}
\paragraph{MCQ-1294 [medium]} A design review is specifically concerned with this risk: Distilling only easy teacher outputs creates a deceptively strong narrow model. Which topic should be reviewed first?
\begin{enumerate}[label=\Alph*.]
\item QLoRA
\item LoRA
\item Knowledge distillation
\item Multi-head attention
\end{enumerate}
\noindent\textbf{Answer.} Knowledge distillation
\par\textit{A student learns from teacher probabilities, generated data, rationales, or intermediate representations. Tradeoff: Serving cost falls but rare capabilities and calibration may be lost. Watch for: Distilling only easy teacher outputs creates a deceptively strong narrow model.}
\paragraph{MCQ-1295 [hard]} Which design note most accurately balances the benefits and costs of Knowledge distillation?
\begin{enumerate}[label=\Alph*.]
\item Fine-tuning large models becomes cheaper, while kernels and quantization settings affect stability.
\item Serving cost falls but rare capabilities and calibration may be lost.
\item Diverse interactions improve capacity but add projections and KV memory.
\item Memory and storage fall, while rank and target-module choices limit expressiveness.
\end{enumerate}
\noindent\textbf{Answer.} Serving cost falls but rare capabilities and calibration may be lost.
\par\textit{A student learns from teacher probabilities, generated data, rationales, or intermediate representations. Tradeoff: Serving cost falls but rare capabilities and calibration may be lost. Watch for: Distilling only easy teacher outputs creates a deceptively strong narrow model.}
\paragraph{MCQ-1296 [medium]} Which explanation would be strongest in a system-design interview about Knowledge distillation?
\begin{enumerate}[label=\Alph*.]
\item A student learns from teacher probabilities, generated data, rationales, or intermediate representations.
\item LoRA trains low-rank update matrices while freezing base weights for parameter-efficient adaptation.
\item QLoRA backpropagates through a frozen quantized base model into LoRA adapters using memory-saving quantization techniques.
\item Multiple heads project attention into separate subspaces before concatenation and output projection.
\end{enumerate}
\noindent\textbf{Answer.} A student learns from teacher probabilities, generated data, rationales, or intermediate representations.
\par\textit{A student learns from teacher probabilities, generated data, rationales, or intermediate representations. Tradeoff: Serving cost falls but rare capabilities and calibration may be lost. Watch for: Distilling only easy teacher outputs creates a deceptively strong narrow model.}
\paragraph{MCQ-1297 [hard]} Which production warning is most directly associated with Knowledge distillation?
\begin{enumerate}[label=\Alph*.]
\item Assuming heads are independently interpretable can overstate mechanistic conclusions.
\item Merging incompatible adapters without evaluation creates interference.
\item Distilling only easy teacher outputs creates a deceptively strong narrow model.
\item Confusing training quantization with final serving quantization produces wrong expectations.
\end{enumerate}
\noindent\textbf{Answer.} Distilling only easy teacher outputs creates a deceptively strong narrow model.
\par\textit{A student learns from teacher probabilities, generated data, rationales, or intermediate representations. Tradeoff: Serving cost falls but rare capabilities and calibration may be lost. Watch for: Distilling only easy teacher outputs creates a deceptively strong narrow model.}
\paragraph{FC-1298 [medium]} Define Knowledge distillation in interview-ready language.
\noindent\textbf{Answer.} A student learns from teacher probabilities, generated data, rationales, or intermediate representations.
\par\textit{A student learns from teacher probabilities, generated data, rationales, or intermediate representations.}
\paragraph{FC-1299 [hard]} State the main tradeoff and production pitfall for Knowledge distillation.
\noindent\textbf{Answer.} Tradeoff: Serving cost falls but rare capabilities and calibration may be lost. Pitfall: Distilling only easy teacher outputs creates a deceptively strong narrow model.
\par\textit{Tradeoff: Serving cost falls but rare capabilities and calibration may be lost. Pitfall: Distilling only easy teacher outputs creates a deceptively strong narrow model.}
\paragraph{LA-1300 [hard]} Explain Knowledge distillation from first principles, then show how you would evaluate and operate it in production.
\noindent\textbf{Answer.} Start with the mechanism: A student learns from teacher probabilities, generated data, rationales, or intermediate representations. Make the tradeoff explicit: Serving cost falls but rare capabilities and calibration may be lost. Define workload-specific offline metrics and a latency/cost budget, test important slices, instrument the critical path, and create a rollback criterion. The production review must explicitly guard against this failure: Distilling only easy teacher outputs creates a deceptively strong narrow model.
\par\textit{A strong answer connects mechanism, measurable tradeoffs, failure isolation, observability, and rollback rather than listing product features.}
\subsection{Inference Optimization}
\paragraph{MCQ-1301 [medium]} Which statement best describes KV caching?
\begin{enumerate}[label=\Alph*.]
\item Requests join and leave a running decode batch at iteration boundaries rather than waiting for a fixed batch.
\item Pipeline parallelism assigns layer stages to devices and streams microbatches through them.
\item MoE experts are distributed across devices and tokens are routed with all-to-all communication.
\item Autoregressive decoding stores prior keys and values so each new token avoids recomputing attention over the entire prefix.
\end{enumerate}
\noindent\textbf{Answer.} Autoregressive decoding stores prior keys and values so each new token avoids recomputing attention over the entire prefix.
\par\textit{Autoregressive decoding stores prior keys and values so each new token avoids recomputing attention over the entire prefix. Tradeoff: Decode speed improves while cache memory grows with layers, sequence, batch, and KV heads. Watch for: Capacity planning only model weights causes out-of-memory failures at long context.}
\paragraph{MCQ-1302 [hard]} What is the central engineering tradeoff of KV caching?
\begin{enumerate}[label=\Alph*.]
\item It scales model depth but creates bubbles and per-request latency.
\item Parameter capacity scales, while communication and load balance determine performance.
\item Throughput and utilization rise, while scheduling fairness and latency isolation become harder.
\item Decode speed improves while cache memory grows with layers, sequence, batch, and KV heads.
\end{enumerate}
\noindent\textbf{Answer.} Decode speed improves while cache memory grows with layers, sequence, batch, and KV heads.
\par\textit{Autoregressive decoding stores prior keys and values so each new token avoids recomputing attention over the entire prefix. Tradeoff: Decode speed improves while cache memory grows with layers, sequence, batch, and KV heads. Watch for: Capacity planning only model weights causes out-of-memory failures at long context.}
\paragraph{MCQ-1303 [hard]} Which failure mode should an interviewer expect you to identify for KV caching?
\begin{enumerate}[label=\Alph*.]
\item Capacity planning only model weights causes out-of-memory failures at long context.
\item Too few microbatches leave stages idle.
\item Optimizing total tokens per second alone can starve short requests.
\item Average expert load hides hot experts that gate step time.
\end{enumerate}
\noindent\textbf{Answer.} Capacity planning only model weights causes out-of-memory failures at long context.
\par\textit{Autoregressive decoding stores prior keys and values so each new token avoids recomputing attention over the entire prefix. Tradeoff: Decode speed improves while cache memory grows with layers, sequence, batch, and KV heads. Watch for: Capacity planning only model weights causes out-of-memory failures at long context.}
\paragraph{MCQ-1304 [medium]} A design review is specifically concerned with this risk: Capacity planning only model weights causes out-of-memory failures at long context. Which topic should be reviewed first?
\begin{enumerate}[label=\Alph*.]
\item Expert parallelism
\item Continuous batching
\item KV caching
\item Pipeline parallelism
\end{enumerate}
\noindent\textbf{Answer.} KV caching
\par\textit{Autoregressive decoding stores prior keys and values so each new token avoids recomputing attention over the entire prefix. Tradeoff: Decode speed improves while cache memory grows with layers, sequence, batch, and KV heads. Watch for: Capacity planning only model weights causes out-of-memory failures at long context.}
\paragraph{MCQ-1305 [hard]} Which design note most accurately balances the benefits and costs of KV caching?
\begin{enumerate}[label=\Alph*.]
\item It scales model depth but creates bubbles and per-request latency.
\item Throughput and utilization rise, while scheduling fairness and latency isolation become harder.
\item Parameter capacity scales, while communication and load balance determine performance.
\item Decode speed improves while cache memory grows with layers, sequence, batch, and KV heads.
\end{enumerate}
\noindent\textbf{Answer.} Decode speed improves while cache memory grows with layers, sequence, batch, and KV heads.
\par\textit{Autoregressive decoding stores prior keys and values so each new token avoids recomputing attention over the entire prefix. Tradeoff: Decode speed improves while cache memory grows with layers, sequence, batch, and KV heads. Watch for: Capacity planning only model weights causes out-of-memory failures at long context.}
\paragraph{MCQ-1306 [medium]} Which explanation would be strongest in a system-design interview about KV caching?
\begin{enumerate}[label=\Alph*.]
\item Autoregressive decoding stores prior keys and values so each new token avoids recomputing attention over the entire prefix.
\item MoE experts are distributed across devices and tokens are routed with all-to-all communication.
\item Requests join and leave a running decode batch at iteration boundaries rather than waiting for a fixed batch.
\item Pipeline parallelism assigns layer stages to devices and streams microbatches through them.
\end{enumerate}
\noindent\textbf{Answer.} Autoregressive decoding stores prior keys and values so each new token avoids recomputing attention over the entire prefix.
\par\textit{Autoregressive decoding stores prior keys and values so each new token avoids recomputing attention over the entire prefix. Tradeoff: Decode speed improves while cache memory grows with layers, sequence, batch, and KV heads. Watch for: Capacity planning only model weights causes out-of-memory failures at long context.}
\paragraph{MCQ-1307 [hard]} Which production warning is most directly associated with KV caching?
\begin{enumerate}[label=\Alph*.]
\item Capacity planning only model weights causes out-of-memory failures at long context.
\item Too few microbatches leave stages idle.
\item Optimizing total tokens per second alone can starve short requests.
\item Average expert load hides hot experts that gate step time.
\end{enumerate}
\noindent\textbf{Answer.} Capacity planning only model weights causes out-of-memory failures at long context.
\par\textit{Autoregressive decoding stores prior keys and values so each new token avoids recomputing attention over the entire prefix. Tradeoff: Decode speed improves while cache memory grows with layers, sequence, batch, and KV heads. Watch for: Capacity planning only model weights causes out-of-memory failures at long context.}
\paragraph{FC-1308 [medium]} Define KV caching in interview-ready language.
\noindent\textbf{Answer.} Autoregressive decoding stores prior keys and values so each new token avoids recomputing attention over the entire prefix.
\par\textit{Autoregressive decoding stores prior keys and values so each new token avoids recomputing attention over the entire prefix.}
\paragraph{FC-1309 [hard]} State the main tradeoff and production pitfall for KV caching.
\noindent\textbf{Answer.} Tradeoff: Decode speed improves while cache memory grows with layers, sequence, batch, and KV heads. Pitfall: Capacity planning only model weights causes out-of-memory failures at long context.
\par\textit{Tradeoff: Decode speed improves while cache memory grows with layers, sequence, batch, and KV heads. Pitfall: Capacity planning only model weights causes out-of-memory failures at long context.}
\paragraph{LA-1310 [hard]} Explain KV caching from first principles, then show how you would evaluate and operate it in production.
\noindent\textbf{Answer.} Start with the mechanism: Autoregressive decoding stores prior keys and values so each new token avoids recomputing attention over the entire prefix. Make the tradeoff explicit: Decode speed improves while cache memory grows with layers, sequence, batch, and KV heads. Define workload-specific offline metrics and a latency/cost budget, test important slices, instrument the critical path, and create a rollback criterion. The production review must explicitly guard against this failure: Capacity planning only model weights causes out-of-memory failures at long context.
\par\textit{A strong answer connects mechanism, measurable tradeoffs, failure isolation, observability, and rollback rather than listing product features.}
\paragraph{MCQ-1311 [medium]} Which statement best describes PagedAttention?
\begin{enumerate}[label=\Alph*.]
\item FlashAttention tiles exact attention to reduce high-bandwidth-memory traffic and materialization of the score matrix.
\item Tensor parallelism shards matrix operations across devices and communicates partial results within layers.
\item PagedAttention manages KV cache in non-contiguous blocks using paging-inspired allocation to reduce fragmentation and enable sharing.
\item Replicas serve independent requests and scale throughput with routing and load balancing.
\end{enumerate}
\noindent\textbf{Answer.} PagedAttention manages KV cache in non-contiguous blocks using paging-inspired allocation to reduce fragmentation and enable sharing.
\par\textit{PagedAttention manages KV cache in non-contiguous blocks using paging-inspired allocation to reduce fragmentation and enable sharing. Tradeoff: Utilization and batching improve, while block tables and kernels add runtime complexity. Watch for: Oversized blocks waste tail space; tiny blocks add metadata overhead.}
\paragraph{MCQ-1312 [hard]} What is the central engineering tradeoff of PagedAttention?
\begin{enumerate}[label=\Alph*.]
\item It is operationally simple but duplicates model weights on every replica.
\item It enables large models and lower per-request latency but requires high-bandwidth interconnects.
\item Utilization and batching improve, while block tables and kernels add runtime complexity.
\item Wall-clock speed and memory improve without approximate attention, but kernels depend on shapes and hardware.
\end{enumerate}
\noindent\textbf{Answer.} Utilization and batching improve, while block tables and kernels add runtime complexity.
\par\textit{PagedAttention manages KV cache in non-contiguous blocks using paging-inspired allocation to reduce fragmentation and enable sharing. Tradeoff: Utilization and batching improve, while block tables and kernels add runtime complexity. Watch for: Oversized blocks waste tail space; tiny blocks add metadata overhead.}
\paragraph{MCQ-1313 [hard]} Which failure mode should an interviewer expect you to identify for PagedAttention?
\begin{enumerate}[label=\Alph*.]
\item Random routing ignores prefix-cache locality and uneven sequence lengths.
\item Oversized blocks waste tail space; tiny blocks add metadata overhead.
\item Calling it subquadratic compute confuses IO optimization with arithmetic complexity.
\item Spanning slow network links can make communication dominate compute.
\end{enumerate}
\noindent\textbf{Answer.} Oversized blocks waste tail space; tiny blocks add metadata overhead.
\par\textit{PagedAttention manages KV cache in non-contiguous blocks using paging-inspired allocation to reduce fragmentation and enable sharing. Tradeoff: Utilization and batching improve, while block tables and kernels add runtime complexity. Watch for: Oversized blocks waste tail space; tiny blocks add metadata overhead.}
\paragraph{MCQ-1314 [medium]} A design review is specifically concerned with this risk: Oversized blocks waste tail space; tiny blocks add metadata overhead. Which topic should be reviewed first?
\begin{enumerate}[label=\Alph*.]
\item FlashAttention
\item Data parallel inference
\item PagedAttention
\item Tensor parallelism
\end{enumerate}
\noindent\textbf{Answer.} PagedAttention
\par\textit{PagedAttention manages KV cache in non-contiguous blocks using paging-inspired allocation to reduce fragmentation and enable sharing. Tradeoff: Utilization and batching improve, while block tables and kernels add runtime complexity. Watch for: Oversized blocks waste tail space; tiny blocks add metadata overhead.}
\paragraph{MCQ-1315 [hard]} Which design note most accurately balances the benefits and costs of PagedAttention?
\begin{enumerate}[label=\Alph*.]
\item Wall-clock speed and memory improve without approximate attention, but kernels depend on shapes and hardware.
\item It is operationally simple but duplicates model weights on every replica.
\item Utilization and batching improve, while block tables and kernels add runtime complexity.
\item It enables large models and lower per-request latency but requires high-bandwidth interconnects.
\end{enumerate}
\noindent\textbf{Answer.} Utilization and batching improve, while block tables and kernels add runtime complexity.
\par\textit{PagedAttention manages KV cache in non-contiguous blocks using paging-inspired allocation to reduce fragmentation and enable sharing. Tradeoff: Utilization and batching improve, while block tables and kernels add runtime complexity. Watch for: Oversized blocks waste tail space; tiny blocks add metadata overhead.}
\paragraph{MCQ-1316 [medium]} Which explanation would be strongest in a system-design interview about PagedAttention?
\begin{enumerate}[label=\Alph*.]
\item PagedAttention manages KV cache in non-contiguous blocks using paging-inspired allocation to reduce fragmentation and enable sharing.
\item Tensor parallelism shards matrix operations across devices and communicates partial results within layers.
\item Replicas serve independent requests and scale throughput with routing and load balancing.
\item FlashAttention tiles exact attention to reduce high-bandwidth-memory traffic and materialization of the score matrix.
\end{enumerate}
\noindent\textbf{Answer.} PagedAttention manages KV cache in non-contiguous blocks using paging-inspired allocation to reduce fragmentation and enable sharing.
\par\textit{PagedAttention manages KV cache in non-contiguous blocks using paging-inspired allocation to reduce fragmentation and enable sharing. Tradeoff: Utilization and batching improve, while block tables and kernels add runtime complexity. Watch for: Oversized blocks waste tail space; tiny blocks add metadata overhead.}
\paragraph{MCQ-1317 [hard]} Which production warning is most directly associated with PagedAttention?
\begin{enumerate}[label=\Alph*.]
\item Spanning slow network links can make communication dominate compute.
\item Calling it subquadratic compute confuses IO optimization with arithmetic complexity.
\item Random routing ignores prefix-cache locality and uneven sequence lengths.
\item Oversized blocks waste tail space; tiny blocks add metadata overhead.
\end{enumerate}
\noindent\textbf{Answer.} Oversized blocks waste tail space; tiny blocks add metadata overhead.
\par\textit{PagedAttention manages KV cache in non-contiguous blocks using paging-inspired allocation to reduce fragmentation and enable sharing. Tradeoff: Utilization and batching improve, while block tables and kernels add runtime complexity. Watch for: Oversized blocks waste tail space; tiny blocks add metadata overhead.}
\paragraph{FC-1318 [medium]} Define PagedAttention in interview-ready language.
\noindent\textbf{Answer.} PagedAttention manages KV cache in non-contiguous blocks using paging-inspired allocation to reduce fragmentation and enable sharing.
\par\textit{PagedAttention manages KV cache in non-contiguous blocks using paging-inspired allocation to reduce fragmentation and enable sharing.}
\paragraph{FC-1319 [hard]} State the main tradeoff and production pitfall for PagedAttention.
\noindent\textbf{Answer.} Tradeoff: Utilization and batching improve, while block tables and kernels add runtime complexity. Pitfall: Oversized blocks waste tail space; tiny blocks add metadata overhead.
\par\textit{Tradeoff: Utilization and batching improve, while block tables and kernels add runtime complexity. Pitfall: Oversized blocks waste tail space; tiny blocks add metadata overhead.}
\paragraph{LA-1320 [hard]} Explain PagedAttention from first principles, then show how you would evaluate and operate it in production.
\noindent\textbf{Answer.} Start with the mechanism: PagedAttention manages KV cache in non-contiguous blocks using paging-inspired allocation to reduce fragmentation and enable sharing. Make the tradeoff explicit: Utilization and batching improve, while block tables and kernels add runtime complexity. Define workload-specific offline metrics and a latency/cost budget, test important slices, instrument the critical path, and create a rollback criterion. The production review must explicitly guard against this failure: Oversized blocks waste tail space; tiny blocks add metadata overhead.
\par\textit{A strong answer connects mechanism, measurable tradeoffs, failure isolation, observability, and rollback rather than listing product features.}
\paragraph{MCQ-1321 [medium]} Which statement best describes Continuous batching?
\begin{enumerate}[label=\Alph*.]
\item Requests join and leave a running decode batch at iteration boundaries rather than waiting for a fixed batch.
\item PagedAttention manages KV cache in non-contiguous blocks using paging-inspired allocation to reduce fragmentation and enable sharing.
\item Autoregressive decoding stores prior keys and values so each new token avoids recomputing attention over the entire prefix.
\item FlashAttention tiles exact attention to reduce high-bandwidth-memory traffic and materialization of the score matrix.
\end{enumerate}
\noindent\textbf{Answer.} Requests join and leave a running decode batch at iteration boundaries rather than waiting for a fixed batch.
\par\textit{Requests join and leave a running decode batch at iteration boundaries rather than waiting for a fixed batch. Tradeoff: Throughput and utilization rise, while scheduling fairness and latency isolation become harder. Watch for: Optimizing total tokens per second alone can starve short requests.}
\paragraph{MCQ-1322 [hard]} What is the central engineering tradeoff of Continuous batching?
\begin{enumerate}[label=\Alph*.]
\item Utilization and batching improve, while block tables and kernels add runtime complexity.
\item Throughput and utilization rise, while scheduling fairness and latency isolation become harder.
\item Decode speed improves while cache memory grows with layers, sequence, batch, and KV heads.
\item Wall-clock speed and memory improve without approximate attention, but kernels depend on shapes and hardware.
\end{enumerate}
\noindent\textbf{Answer.} Throughput and utilization rise, while scheduling fairness and latency isolation become harder.
\par\textit{Requests join and leave a running decode batch at iteration boundaries rather than waiting for a fixed batch. Tradeoff: Throughput and utilization rise, while scheduling fairness and latency isolation become harder. Watch for: Optimizing total tokens per second alone can starve short requests.}
\paragraph{MCQ-1323 [hard]} Which failure mode should an interviewer expect you to identify for Continuous batching?
\begin{enumerate}[label=\Alph*.]
\item Oversized blocks waste tail space; tiny blocks add metadata overhead.
\item Calling it subquadratic compute confuses IO optimization with arithmetic complexity.
\item Capacity planning only model weights causes out-of-memory failures at long context.
\item Optimizing total tokens per second alone can starve short requests.
\end{enumerate}
\noindent\textbf{Answer.} Optimizing total tokens per second alone can starve short requests.
\par\textit{Requests join and leave a running decode batch at iteration boundaries rather than waiting for a fixed batch. Tradeoff: Throughput and utilization rise, while scheduling fairness and latency isolation become harder. Watch for: Optimizing total tokens per second alone can starve short requests.}
\paragraph{MCQ-1324 [medium]} A design review is specifically concerned with this risk: Optimizing total tokens per second alone can starve short requests. Which topic should be reviewed first?
\begin{enumerate}[label=\Alph*.]
\item PagedAttention
\item KV caching
\item Continuous batching
\item FlashAttention
\end{enumerate}
\noindent\textbf{Answer.} Continuous batching
\par\textit{Requests join and leave a running decode batch at iteration boundaries rather than waiting for a fixed batch. Tradeoff: Throughput and utilization rise, while scheduling fairness and latency isolation become harder. Watch for: Optimizing total tokens per second alone can starve short requests.}
\paragraph{MCQ-1325 [hard]} Which design note most accurately balances the benefits and costs of Continuous batching?
\begin{enumerate}[label=\Alph*.]
\item Wall-clock speed and memory improve without approximate attention, but kernels depend on shapes and hardware.
\item Throughput and utilization rise, while scheduling fairness and latency isolation become harder.
\item Decode speed improves while cache memory grows with layers, sequence, batch, and KV heads.
\item Utilization and batching improve, while block tables and kernels add runtime complexity.
\end{enumerate}
\noindent\textbf{Answer.} Throughput and utilization rise, while scheduling fairness and latency isolation become harder.
\par\textit{Requests join and leave a running decode batch at iteration boundaries rather than waiting for a fixed batch. Tradeoff: Throughput and utilization rise, while scheduling fairness and latency isolation become harder. Watch for: Optimizing total tokens per second alone can starve short requests.}
\paragraph{MCQ-1326 [medium]} Which explanation would be strongest in a system-design interview about Continuous batching?
\begin{enumerate}[label=\Alph*.]
\item Requests join and leave a running decode batch at iteration boundaries rather than waiting for a fixed batch.
\item Autoregressive decoding stores prior keys and values so each new token avoids recomputing attention over the entire prefix.
\item PagedAttention manages KV cache in non-contiguous blocks using paging-inspired allocation to reduce fragmentation and enable sharing.
\item FlashAttention tiles exact attention to reduce high-bandwidth-memory traffic and materialization of the score matrix.
\end{enumerate}
\noindent\textbf{Answer.} Requests join and leave a running decode batch at iteration boundaries rather than waiting for a fixed batch.
\par\textit{Requests join and leave a running decode batch at iteration boundaries rather than waiting for a fixed batch. Tradeoff: Throughput and utilization rise, while scheduling fairness and latency isolation become harder. Watch for: Optimizing total tokens per second alone can starve short requests.}
\paragraph{MCQ-1327 [hard]} Which production warning is most directly associated with Continuous batching?
\begin{enumerate}[label=\Alph*.]
\item Optimizing total tokens per second alone can starve short requests.
\item Oversized blocks waste tail space; tiny blocks add metadata overhead.
\item Calling it subquadratic compute confuses IO optimization with arithmetic complexity.
\item Capacity planning only model weights causes out-of-memory failures at long context.
\end{enumerate}
\noindent\textbf{Answer.} Optimizing total tokens per second alone can starve short requests.
\par\textit{Requests join and leave a running decode batch at iteration boundaries rather than waiting for a fixed batch. Tradeoff: Throughput and utilization rise, while scheduling fairness and latency isolation become harder. Watch for: Optimizing total tokens per second alone can starve short requests.}
\paragraph{FC-1328 [medium]} Define Continuous batching in interview-ready language.
\noindent\textbf{Answer.} Requests join and leave a running decode batch at iteration boundaries rather than waiting for a fixed batch.
\par\textit{Requests join and leave a running decode batch at iteration boundaries rather than waiting for a fixed batch.}
\paragraph{FC-1329 [hard]} State the main tradeoff and production pitfall for Continuous batching.
\noindent\textbf{Answer.} Tradeoff: Throughput and utilization rise, while scheduling fairness and latency isolation become harder. Pitfall: Optimizing total tokens per second alone can starve short requests.
\par\textit{Tradeoff: Throughput and utilization rise, while scheduling fairness and latency isolation become harder. Pitfall: Optimizing total tokens per second alone can starve short requests.}
\paragraph{LA-1330 [hard]} Explain Continuous batching from first principles, then show how you would evaluate and operate it in production.
\noindent\textbf{Answer.} Start with the mechanism: Requests join and leave a running decode batch at iteration boundaries rather than waiting for a fixed batch. Make the tradeoff explicit: Throughput and utilization rise, while scheduling fairness and latency isolation become harder. Define workload-specific offline metrics and a latency/cost budget, test important slices, instrument the critical path, and create a rollback criterion. The production review must explicitly guard against this failure: Optimizing total tokens per second alone can starve short requests.
\par\textit{A strong answer connects mechanism, measurable tradeoffs, failure isolation, observability, and rollback rather than listing product features.}
\paragraph{MCQ-1331 [medium]} Which statement best describes Prefix caching?
\begin{enumerate}[label=\Alph*.]
\item Grouped-query and multi-query attention share key-value heads across query heads to shrink KV cache and memory bandwidth.
\item Tensor parallelism shards matrix operations across devices and communicates partial results within layers.
\item Shared prompt prefixes reuse precomputed KV blocks across requests with matching tokens and model state.
\item A cheaper draft proposes tokens and a target model verifies them while preserving the target distribution.
\end{enumerate}
\noindent\textbf{Answer.} Shared prompt prefixes reuse precomputed KV blocks across requests with matching tokens and model state.
\par\textit{Shared prompt prefixes reuse precomputed KV blocks across requests with matching tokens and model state. Tradeoff: Repeated system prompts become cheaper, but hit rate, invalidation, and tenant isolation matter. Watch for: Cache keys that omit adapters or model revision return incorrect state.}
\paragraph{MCQ-1332 [hard]} What is the central engineering tradeoff of Prefix caching?
\begin{enumerate}[label=\Alph*.]
\item Decode efficiency improves, with possible quality loss relative to full multi-head KV.
\item Repeated system prompts become cheaper, but hit rate, invalidation, and tenant isolation matter.
\item It enables large models and lower per-request latency but requires high-bandwidth interconnects.
\item Latency falls when acceptance is high, but drafting and verification overhead can dominate.
\end{enumerate}
\noindent\textbf{Answer.} Repeated system prompts become cheaper, but hit rate, invalidation, and tenant isolation matter.
\par\textit{Shared prompt prefixes reuse precomputed KV blocks across requests with matching tokens and model state. Tradeoff: Repeated system prompts become cheaper, but hit rate, invalidation, and tenant isolation matter. Watch for: Cache keys that omit adapters or model revision return incorrect state.}
\paragraph{MCQ-1333 [hard]} Which failure mode should an interviewer expect you to identify for Prefix caching?
\begin{enumerate}[label=\Alph*.]
\item Converting weights naively after training does not reproduce a model trained for shared KV heads.
\item Spanning slow network links can make communication dominate compute.
\item Cache keys that omit adapters or model revision return incorrect state.
\item Measuring draft speed without acceptance rate gives meaningless projections.
\end{enumerate}
\noindent\textbf{Answer.} Cache keys that omit adapters or model revision return incorrect state.
\par\textit{Shared prompt prefixes reuse precomputed KV blocks across requests with matching tokens and model state. Tradeoff: Repeated system prompts become cheaper, but hit rate, invalidation, and tenant isolation matter. Watch for: Cache keys that omit adapters or model revision return incorrect state.}
\paragraph{MCQ-1334 [medium]} A design review is specifically concerned with this risk: Cache keys that omit adapters or model revision return incorrect state. Which topic should be reviewed first?
\begin{enumerate}[label=\Alph*.]
\item Prefix caching
\item Tensor parallelism
\item Speculative decoding
\item GQA and MQA
\end{enumerate}
\noindent\textbf{Answer.} Prefix caching
\par\textit{Shared prompt prefixes reuse precomputed KV blocks across requests with matching tokens and model state. Tradeoff: Repeated system prompts become cheaper, but hit rate, invalidation, and tenant isolation matter. Watch for: Cache keys that omit adapters or model revision return incorrect state.}
\paragraph{MCQ-1335 [hard]} Which design note most accurately balances the benefits and costs of Prefix caching?
\begin{enumerate}[label=\Alph*.]
\item It enables large models and lower per-request latency but requires high-bandwidth interconnects.
\item Decode efficiency improves, with possible quality loss relative to full multi-head KV.
\item Repeated system prompts become cheaper, but hit rate, invalidation, and tenant isolation matter.
\item Latency falls when acceptance is high, but drafting and verification overhead can dominate.
\end{enumerate}
\noindent\textbf{Answer.} Repeated system prompts become cheaper, but hit rate, invalidation, and tenant isolation matter.
\par\textit{Shared prompt prefixes reuse precomputed KV blocks across requests with matching tokens and model state. Tradeoff: Repeated system prompts become cheaper, but hit rate, invalidation, and tenant isolation matter. Watch for: Cache keys that omit adapters or model revision return incorrect state.}
\paragraph{MCQ-1336 [medium]} Which explanation would be strongest in a system-design interview about Prefix caching?
\begin{enumerate}[label=\Alph*.]
\item Shared prompt prefixes reuse precomputed KV blocks across requests with matching tokens and model state.
\item Tensor parallelism shards matrix operations across devices and communicates partial results within layers.
\item A cheaper draft proposes tokens and a target model verifies them while preserving the target distribution.
\item Grouped-query and multi-query attention share key-value heads across query heads to shrink KV cache and memory bandwidth.
\end{enumerate}
\noindent\textbf{Answer.} Shared prompt prefixes reuse precomputed KV blocks across requests with matching tokens and model state.
\par\textit{Shared prompt prefixes reuse precomputed KV blocks across requests with matching tokens and model state. Tradeoff: Repeated system prompts become cheaper, but hit rate, invalidation, and tenant isolation matter. Watch for: Cache keys that omit adapters or model revision return incorrect state.}
\paragraph{MCQ-1337 [hard]} Which production warning is most directly associated with Prefix caching?
\begin{enumerate}[label=\Alph*.]
\item Spanning slow network links can make communication dominate compute.
\item Cache keys that omit adapters or model revision return incorrect state.
\item Converting weights naively after training does not reproduce a model trained for shared KV heads.
\item Measuring draft speed without acceptance rate gives meaningless projections.
\end{enumerate}
\noindent\textbf{Answer.} Cache keys that omit adapters or model revision return incorrect state.
\par\textit{Shared prompt prefixes reuse precomputed KV blocks across requests with matching tokens and model state. Tradeoff: Repeated system prompts become cheaper, but hit rate, invalidation, and tenant isolation matter. Watch for: Cache keys that omit adapters or model revision return incorrect state.}
\paragraph{FC-1338 [medium]} Define Prefix caching in interview-ready language.
\noindent\textbf{Answer.} Shared prompt prefixes reuse precomputed KV blocks across requests with matching tokens and model state.
\par\textit{Shared prompt prefixes reuse precomputed KV blocks across requests with matching tokens and model state.}
\paragraph{FC-1339 [hard]} State the main tradeoff and production pitfall for Prefix caching.
\noindent\textbf{Answer.} Tradeoff: Repeated system prompts become cheaper, but hit rate, invalidation, and tenant isolation matter. Pitfall: Cache keys that omit adapters or model revision return incorrect state.
\par\textit{Tradeoff: Repeated system prompts become cheaper, but hit rate, invalidation, and tenant isolation matter. Pitfall: Cache keys that omit adapters or model revision return incorrect state.}
\paragraph{LA-1340 [hard]} Explain Prefix caching from first principles, then show how you would evaluate and operate it in production.
\noindent\textbf{Answer.} Start with the mechanism: Shared prompt prefixes reuse precomputed KV blocks across requests with matching tokens and model state. Make the tradeoff explicit: Repeated system prompts become cheaper, but hit rate, invalidation, and tenant isolation matter. Define workload-specific offline metrics and a latency/cost budget, test important slices, instrument the critical path, and create a rollback criterion. The production review must explicitly guard against this failure: Cache keys that omit adapters or model revision return incorrect state.
\par\textit{A strong answer connects mechanism, measurable tradeoffs, failure isolation, observability, and rollback rather than listing product features.}
\paragraph{MCQ-1341 [medium]} Which statement best describes FlashAttention?
\begin{enumerate}[label=\Alph*.]
\item PagedAttention manages KV cache in non-contiguous blocks using paging-inspired allocation to reduce fragmentation and enable sharing.
\item FlashAttention tiles exact attention to reduce high-bandwidth-memory traffic and materialization of the score matrix.
\item MoE experts are distributed across devices and tokens are routed with all-to-all communication.
\item Requests join and leave a running decode batch at iteration boundaries rather than waiting for a fixed batch.
\end{enumerate}
\noindent\textbf{Answer.} FlashAttention tiles exact attention to reduce high-bandwidth-memory traffic and materialization of the score matrix.
\par\textit{FlashAttention tiles exact attention to reduce high-bandwidth-memory traffic and materialization of the score matrix. Tradeoff: Wall-clock speed and memory improve without approximate attention, but kernels depend on shapes and hardware. Watch for: Calling it subquadratic compute confuses IO optimization with arithmetic complexity.}
\paragraph{MCQ-1342 [hard]} What is the central engineering tradeoff of FlashAttention?
\begin{enumerate}[label=\Alph*.]
\item Throughput and utilization rise, while scheduling fairness and latency isolation become harder.
\item Wall-clock speed and memory improve without approximate attention, but kernels depend on shapes and hardware.
\item Parameter capacity scales, while communication and load balance determine performance.
\item Utilization and batching improve, while block tables and kernels add runtime complexity.
\end{enumerate}
\noindent\textbf{Answer.} Wall-clock speed and memory improve without approximate attention, but kernels depend on shapes and hardware.
\par\textit{FlashAttention tiles exact attention to reduce high-bandwidth-memory traffic and materialization of the score matrix. Tradeoff: Wall-clock speed and memory improve without approximate attention, but kernels depend on shapes and hardware. Watch for: Calling it subquadratic compute confuses IO optimization with arithmetic complexity.}
\paragraph{MCQ-1343 [hard]} Which failure mode should an interviewer expect you to identify for FlashAttention?
\begin{enumerate}[label=\Alph*.]
\item Oversized blocks waste tail space; tiny blocks add metadata overhead.
\item Average expert load hides hot experts that gate step time.
\item Optimizing total tokens per second alone can starve short requests.
\item Calling it subquadratic compute confuses IO optimization with arithmetic complexity.
\end{enumerate}
\noindent\textbf{Answer.} Calling it subquadratic compute confuses IO optimization with arithmetic complexity.
\par\textit{FlashAttention tiles exact attention to reduce high-bandwidth-memory traffic and materialization of the score matrix. Tradeoff: Wall-clock speed and memory improve without approximate attention, but kernels depend on shapes and hardware. Watch for: Calling it subquadratic compute confuses IO optimization with arithmetic complexity.}
\paragraph{MCQ-1344 [medium]} A design review is specifically concerned with this risk: Calling it subquadratic compute confuses IO optimization with arithmetic complexity. Which topic should be reviewed first?
\begin{enumerate}[label=\Alph*.]
\item Expert parallelism
\item FlashAttention
\item PagedAttention
\item Continuous batching
\end{enumerate}
\noindent\textbf{Answer.} FlashAttention
\par\textit{FlashAttention tiles exact attention to reduce high-bandwidth-memory traffic and materialization of the score matrix. Tradeoff: Wall-clock speed and memory improve without approximate attention, but kernels depend on shapes and hardware. Watch for: Calling it subquadratic compute confuses IO optimization with arithmetic complexity.}
\paragraph{MCQ-1345 [hard]} Which design note most accurately balances the benefits and costs of FlashAttention?
\begin{enumerate}[label=\Alph*.]
\item Utilization and batching improve, while block tables and kernels add runtime complexity.
\item Wall-clock speed and memory improve without approximate attention, but kernels depend on shapes and hardware.
\item Parameter capacity scales, while communication and load balance determine performance.
\item Throughput and utilization rise, while scheduling fairness and latency isolation become harder.
\end{enumerate}
\noindent\textbf{Answer.} Wall-clock speed and memory improve without approximate attention, but kernels depend on shapes and hardware.
\par\textit{FlashAttention tiles exact attention to reduce high-bandwidth-memory traffic and materialization of the score matrix. Tradeoff: Wall-clock speed and memory improve without approximate attention, but kernels depend on shapes and hardware. Watch for: Calling it subquadratic compute confuses IO optimization with arithmetic complexity.}
\paragraph{MCQ-1346 [medium]} Which explanation would be strongest in a system-design interview about FlashAttention?
\begin{enumerate}[label=\Alph*.]
\item PagedAttention manages KV cache in non-contiguous blocks using paging-inspired allocation to reduce fragmentation and enable sharing.
\item Requests join and leave a running decode batch at iteration boundaries rather than waiting for a fixed batch.
\item MoE experts are distributed across devices and tokens are routed with all-to-all communication.
\item FlashAttention tiles exact attention to reduce high-bandwidth-memory traffic and materialization of the score matrix.
\end{enumerate}
\noindent\textbf{Answer.} FlashAttention tiles exact attention to reduce high-bandwidth-memory traffic and materialization of the score matrix.
\par\textit{FlashAttention tiles exact attention to reduce high-bandwidth-memory traffic and materialization of the score matrix. Tradeoff: Wall-clock speed and memory improve without approximate attention, but kernels depend on shapes and hardware. Watch for: Calling it subquadratic compute confuses IO optimization with arithmetic complexity.}
\paragraph{MCQ-1347 [hard]} Which production warning is most directly associated with FlashAttention?
\begin{enumerate}[label=\Alph*.]
\item Optimizing total tokens per second alone can starve short requests.
\item Oversized blocks waste tail space; tiny blocks add metadata overhead.
\item Calling it subquadratic compute confuses IO optimization with arithmetic complexity.
\item Average expert load hides hot experts that gate step time.
\end{enumerate}
\noindent\textbf{Answer.} Calling it subquadratic compute confuses IO optimization with arithmetic complexity.
\par\textit{FlashAttention tiles exact attention to reduce high-bandwidth-memory traffic and materialization of the score matrix. Tradeoff: Wall-clock speed and memory improve without approximate attention, but kernels depend on shapes and hardware. Watch for: Calling it subquadratic compute confuses IO optimization with arithmetic complexity.}
\paragraph{FC-1348 [medium]} Define FlashAttention in interview-ready language.
\noindent\textbf{Answer.} FlashAttention tiles exact attention to reduce high-bandwidth-memory traffic and materialization of the score matrix.
\par\textit{FlashAttention tiles exact attention to reduce high-bandwidth-memory traffic and materialization of the score matrix.}
\paragraph{FC-1349 [hard]} State the main tradeoff and production pitfall for FlashAttention.
\noindent\textbf{Answer.} Tradeoff: Wall-clock speed and memory improve without approximate attention, but kernels depend on shapes and hardware. Pitfall: Calling it subquadratic compute confuses IO optimization with arithmetic complexity.
\par\textit{Tradeoff: Wall-clock speed and memory improve without approximate attention, but kernels depend on shapes and hardware. Pitfall: Calling it subquadratic compute confuses IO optimization with arithmetic complexity.}
\paragraph{LA-1350 [hard]} Explain FlashAttention from first principles, then show how you would evaluate and operate it in production.
\noindent\textbf{Answer.} Start with the mechanism: FlashAttention tiles exact attention to reduce high-bandwidth-memory traffic and materialization of the score matrix. Make the tradeoff explicit: Wall-clock speed and memory improve without approximate attention, but kernels depend on shapes and hardware. Define workload-specific offline metrics and a latency/cost budget, test important slices, instrument the critical path, and create a rollback criterion. The production review must explicitly guard against this failure: Calling it subquadratic compute confuses IO optimization with arithmetic complexity.
\par\textit{A strong answer connects mechanism, measurable tradeoffs, failure isolation, observability, and rollback rather than listing product features.}
\paragraph{MCQ-1351 [medium]} Which statement best describes GQA and MQA?
\begin{enumerate}[label=\Alph*.]
\item PagedAttention manages KV cache in non-contiguous blocks using paging-inspired allocation to reduce fragmentation and enable sharing.
\item Requests join and leave a running decode batch at iteration boundaries rather than waiting for a fixed batch.
\item Servers emit tokens incrementally with backpressure, disconnect handling, and final usage accounting.
\item Grouped-query and multi-query attention share key-value heads across query heads to shrink KV cache and memory bandwidth.
\end{enumerate}
\noindent\textbf{Answer.} Grouped-query and multi-query attention share key-value heads across query heads to shrink KV cache and memory bandwidth.
\par\textit{Grouped-query and multi-query attention share key-value heads across query heads to shrink KV cache and memory bandwidth. Tradeoff: Decode efficiency improves, with possible quality loss relative to full multi-head KV. Watch for: Converting weights naively after training does not reproduce a model trained for shared KV heads.}
\paragraph{MCQ-1352 [hard]} What is the central engineering tradeoff of GQA and MQA?
\begin{enumerate}[label=\Alph*.]
\item Throughput and utilization rise, while scheduling fairness and latency isolation become harder.
\item Perceived latency improves, but partial outputs complicate moderation and retries.
\item Decode efficiency improves, with possible quality loss relative to full multi-head KV.
\item Utilization and batching improve, while block tables and kernels add runtime complexity.
\end{enumerate}
\noindent\textbf{Answer.} Decode efficiency improves, with possible quality loss relative to full multi-head KV.
\par\textit{Grouped-query and multi-query attention share key-value heads across query heads to shrink KV cache and memory bandwidth. Tradeoff: Decode efficiency improves, with possible quality loss relative to full multi-head KV. Watch for: Converting weights naively after training does not reproduce a model trained for shared KV heads.}
\paragraph{MCQ-1353 [hard]} Which failure mode should an interviewer expect you to identify for GQA and MQA?
\begin{enumerate}[label=\Alph*.]
\item Oversized blocks waste tail space; tiny blocks add metadata overhead.
\item Converting weights naively after training does not reproduce a model trained for shared KV heads.
\item Retrying after visible partial output can duplicate text or side effects.
\item Optimizing total tokens per second alone can starve short requests.
\end{enumerate}
\noindent\textbf{Answer.} Converting weights naively after training does not reproduce a model trained for shared KV heads.
\par\textit{Grouped-query and multi-query attention share key-value heads across query heads to shrink KV cache and memory bandwidth. Tradeoff: Decode efficiency improves, with possible quality loss relative to full multi-head KV. Watch for: Converting weights naively after training does not reproduce a model trained for shared KV heads.}
\paragraph{MCQ-1354 [medium]} A design review is specifically concerned with this risk: Converting weights naively after training does not reproduce a model trained for shared KV heads. Which topic should be reviewed first?
\begin{enumerate}[label=\Alph*.]
\item PagedAttention
\item Streaming generation
\item GQA and MQA
\item Continuous batching
\end{enumerate}
\noindent\textbf{Answer.} GQA and MQA
\par\textit{Grouped-query and multi-query attention share key-value heads across query heads to shrink KV cache and memory bandwidth. Tradeoff: Decode efficiency improves, with possible quality loss relative to full multi-head KV. Watch for: Converting weights naively after training does not reproduce a model trained for shared KV heads.}
\paragraph{MCQ-1355 [hard]} Which design note most accurately balances the benefits and costs of GQA and MQA?
\begin{enumerate}[label=\Alph*.]
\item Throughput and utilization rise, while scheduling fairness and latency isolation become harder.
\item Perceived latency improves, but partial outputs complicate moderation and retries.
\item Utilization and batching improve, while block tables and kernels add runtime complexity.
\item Decode efficiency improves, with possible quality loss relative to full multi-head KV.
\end{enumerate}
\noindent\textbf{Answer.} Decode efficiency improves, with possible quality loss relative to full multi-head KV.
\par\textit{Grouped-query and multi-query attention share key-value heads across query heads to shrink KV cache and memory bandwidth. Tradeoff: Decode efficiency improves, with possible quality loss relative to full multi-head KV. Watch for: Converting weights naively after training does not reproduce a model trained for shared KV heads.}
\paragraph{MCQ-1356 [medium]} Which explanation would be strongest in a system-design interview about GQA and MQA?
\begin{enumerate}[label=\Alph*.]
\item Requests join and leave a running decode batch at iteration boundaries rather than waiting for a fixed batch.
\item PagedAttention manages KV cache in non-contiguous blocks using paging-inspired allocation to reduce fragmentation and enable sharing.
\item Grouped-query and multi-query attention share key-value heads across query heads to shrink KV cache and memory bandwidth.
\item Servers emit tokens incrementally with backpressure, disconnect handling, and final usage accounting.
\end{enumerate}
\noindent\textbf{Answer.} Grouped-query and multi-query attention share key-value heads across query heads to shrink KV cache and memory bandwidth.
\par\textit{Grouped-query and multi-query attention share key-value heads across query heads to shrink KV cache and memory bandwidth. Tradeoff: Decode efficiency improves, with possible quality loss relative to full multi-head KV. Watch for: Converting weights naively after training does not reproduce a model trained for shared KV heads.}
\paragraph{MCQ-1357 [hard]} Which production warning is most directly associated with GQA and MQA?
\begin{enumerate}[label=\Alph*.]
\item Optimizing total tokens per second alone can starve short requests.
\item Retrying after visible partial output can duplicate text or side effects.
\item Oversized blocks waste tail space; tiny blocks add metadata overhead.
\item Converting weights naively after training does not reproduce a model trained for shared KV heads.
\end{enumerate}
\noindent\textbf{Answer.} Converting weights naively after training does not reproduce a model trained for shared KV heads.
\par\textit{Grouped-query and multi-query attention share key-value heads across query heads to shrink KV cache and memory bandwidth. Tradeoff: Decode efficiency improves, with possible quality loss relative to full multi-head KV. Watch for: Converting weights naively after training does not reproduce a model trained for shared KV heads.}
\paragraph{FC-1358 [medium]} Define GQA and MQA in interview-ready language.
\noindent\textbf{Answer.} Grouped-query and multi-query attention share key-value heads across query heads to shrink KV cache and memory bandwidth.
\par\textit{Grouped-query and multi-query attention share key-value heads across query heads to shrink KV cache and memory bandwidth.}
\paragraph{FC-1359 [hard]} State the main tradeoff and production pitfall for GQA and MQA.
\noindent\textbf{Answer.} Tradeoff: Decode efficiency improves, with possible quality loss relative to full multi-head KV. Pitfall: Converting weights naively after training does not reproduce a model trained for shared KV heads.
\par\textit{Tradeoff: Decode efficiency improves, with possible quality loss relative to full multi-head KV. Pitfall: Converting weights naively after training does not reproduce a model trained for shared KV heads.}
\paragraph{LA-1360 [hard]} Explain GQA and MQA from first principles, then show how you would evaluate and operate it in production.
\noindent\textbf{Answer.} Start with the mechanism: Grouped-query and multi-query attention share key-value heads across query heads to shrink KV cache and memory bandwidth. Make the tradeoff explicit: Decode efficiency improves, with possible quality loss relative to full multi-head KV. Define workload-specific offline metrics and a latency/cost budget, test important slices, instrument the critical path, and create a rollback criterion. The production review must explicitly guard against this failure: Converting weights naively after training does not reproduce a model trained for shared KV heads.
\par\textit{A strong answer connects mechanism, measurable tradeoffs, failure isolation, observability, and rollback rather than listing product features.}
\paragraph{MCQ-1361 [medium]} Which statement best describes Speculative decoding?
\begin{enumerate}[label=\Alph*.]
\item A cheaper draft proposes tokens and a target model verifies them while preserving the target distribution.
\item Autoregressive decoding stores prior keys and values so each new token avoids recomputing attention over the entire prefix.
\item Requests join and leave a running decode batch at iteration boundaries rather than waiting for a fixed batch.
\item Replicas serve independent requests and scale throughput with routing and load balancing.
\end{enumerate}
\noindent\textbf{Answer.} A cheaper draft proposes tokens and a target model verifies them while preserving the target distribution.
\par\textit{A cheaper draft proposes tokens and a target model verifies them while preserving the target distribution. Tradeoff: Latency falls when acceptance is high, but drafting and verification overhead can dominate. Watch for: Measuring draft speed without acceptance rate gives meaningless projections.}
\paragraph{MCQ-1362 [hard]} What is the central engineering tradeoff of Speculative decoding?
\begin{enumerate}[label=\Alph*.]
\item Throughput and utilization rise, while scheduling fairness and latency isolation become harder.
\item Decode speed improves while cache memory grows with layers, sequence, batch, and KV heads.
\item Latency falls when acceptance is high, but drafting and verification overhead can dominate.
\item It is operationally simple but duplicates model weights on every replica.
\end{enumerate}
\noindent\textbf{Answer.} Latency falls when acceptance is high, but drafting and verification overhead can dominate.
\par\textit{A cheaper draft proposes tokens and a target model verifies them while preserving the target distribution. Tradeoff: Latency falls when acceptance is high, but drafting and verification overhead can dominate. Watch for: Measuring draft speed without acceptance rate gives meaningless projections.}
\paragraph{MCQ-1363 [hard]} Which failure mode should an interviewer expect you to identify for Speculative decoding?
\begin{enumerate}[label=\Alph*.]
\item Optimizing total tokens per second alone can starve short requests.
\item Capacity planning only model weights causes out-of-memory failures at long context.
\item Measuring draft speed without acceptance rate gives meaningless projections.
\item Random routing ignores prefix-cache locality and uneven sequence lengths.
\end{enumerate}
\noindent\textbf{Answer.} Measuring draft speed without acceptance rate gives meaningless projections.
\par\textit{A cheaper draft proposes tokens and a target model verifies them while preserving the target distribution. Tradeoff: Latency falls when acceptance is high, but drafting and verification overhead can dominate. Watch for: Measuring draft speed without acceptance rate gives meaningless projections.}
\paragraph{MCQ-1364 [medium]} A design review is specifically concerned with this risk: Measuring draft speed without acceptance rate gives meaningless projections. Which topic should be reviewed first?
\begin{enumerate}[label=\Alph*.]
\item Continuous batching
\item Data parallel inference
\item Speculative decoding
\item KV caching
\end{enumerate}
\noindent\textbf{Answer.} Speculative decoding
\par\textit{A cheaper draft proposes tokens and a target model verifies them while preserving the target distribution. Tradeoff: Latency falls when acceptance is high, but drafting and verification overhead can dominate. Watch for: Measuring draft speed without acceptance rate gives meaningless projections.}
\paragraph{MCQ-1365 [hard]} Which design note most accurately balances the benefits and costs of Speculative decoding?
\begin{enumerate}[label=\Alph*.]
\item Decode speed improves while cache memory grows with layers, sequence, batch, and KV heads.
\item Latency falls when acceptance is high, but drafting and verification overhead can dominate.
\item It is operationally simple but duplicates model weights on every replica.
\item Throughput and utilization rise, while scheduling fairness and latency isolation become harder.
\end{enumerate}
\noindent\textbf{Answer.} Latency falls when acceptance is high, but drafting and verification overhead can dominate.
\par\textit{A cheaper draft proposes tokens and a target model verifies them while preserving the target distribution. Tradeoff: Latency falls when acceptance is high, but drafting and verification overhead can dominate. Watch for: Measuring draft speed without acceptance rate gives meaningless projections.}
\paragraph{MCQ-1366 [medium]} Which explanation would be strongest in a system-design interview about Speculative decoding?
\begin{enumerate}[label=\Alph*.]
\item Replicas serve independent requests and scale throughput with routing and load balancing.
\item Autoregressive decoding stores prior keys and values so each new token avoids recomputing attention over the entire prefix.
\item Requests join and leave a running decode batch at iteration boundaries rather than waiting for a fixed batch.
\item A cheaper draft proposes tokens and a target model verifies them while preserving the target distribution.
\end{enumerate}
\noindent\textbf{Answer.} A cheaper draft proposes tokens and a target model verifies them while preserving the target distribution.
\par\textit{A cheaper draft proposes tokens and a target model verifies them while preserving the target distribution. Tradeoff: Latency falls when acceptance is high, but drafting and verification overhead can dominate. Watch for: Measuring draft speed without acceptance rate gives meaningless projections.}
\paragraph{MCQ-1367 [hard]} Which production warning is most directly associated with Speculative decoding?
\begin{enumerate}[label=\Alph*.]
\item Measuring draft speed without acceptance rate gives meaningless projections.
\item Capacity planning only model weights causes out-of-memory failures at long context.
\item Random routing ignores prefix-cache locality and uneven sequence lengths.
\item Optimizing total tokens per second alone can starve short requests.
\end{enumerate}
\noindent\textbf{Answer.} Measuring draft speed without acceptance rate gives meaningless projections.
\par\textit{A cheaper draft proposes tokens and a target model verifies them while preserving the target distribution. Tradeoff: Latency falls when acceptance is high, but drafting and verification overhead can dominate. Watch for: Measuring draft speed without acceptance rate gives meaningless projections.}
\paragraph{FC-1368 [medium]} Define Speculative decoding in interview-ready language.
\noindent\textbf{Answer.} A cheaper draft proposes tokens and a target model verifies them while preserving the target distribution.
\par\textit{A cheaper draft proposes tokens and a target model verifies them while preserving the target distribution.}
\paragraph{FC-1369 [hard]} State the main tradeoff and production pitfall for Speculative decoding.
\noindent\textbf{Answer.} Tradeoff: Latency falls when acceptance is high, but drafting and verification overhead can dominate. Pitfall: Measuring draft speed without acceptance rate gives meaningless projections.
\par\textit{Tradeoff: Latency falls when acceptance is high, but drafting and verification overhead can dominate. Pitfall: Measuring draft speed without acceptance rate gives meaningless projections.}
\paragraph{LA-1370 [hard]} Explain Speculative decoding from first principles, then show how you would evaluate and operate it in production.
\noindent\textbf{Answer.} Start with the mechanism: A cheaper draft proposes tokens and a target model verifies them while preserving the target distribution. Make the tradeoff explicit: Latency falls when acceptance is high, but drafting and verification overhead can dominate. Define workload-specific offline metrics and a latency/cost budget, test important slices, instrument the critical path, and create a rollback criterion. The production review must explicitly guard against this failure: Measuring draft speed without acceptance rate gives meaningless projections.
\par\textit{A strong answer connects mechanism, measurable tradeoffs, failure isolation, observability, and rollback rather than listing product features.}
\paragraph{MCQ-1371 [medium]} Which statement best describes Quantization?
\begin{enumerate}[label=\Alph*.]
\item Autoregressive decoding stores prior keys and values so each new token avoids recomputing attention over the entire prefix.
\item Tensor parallelism shards matrix operations across devices and communicates partial results within layers.
\item Pipeline parallelism assigns layer stages to devices and streams microbatches through them.
\item Quantization maps weights or activations to lower precision using scales, groups, calibration, or quantization-aware training.
\end{enumerate}
\noindent\textbf{Answer.} Quantization maps weights or activations to lower precision using scales, groups, calibration, or quantization-aware training.
\par\textit{Quantization maps weights or activations to lower precision using scales, groups, calibration, or quantization-aware training. Tradeoff: Memory and bandwidth fall, while accuracy and kernel support can degrade. Watch for: A smaller model file may run slower if hardware lacks optimized kernels.}
\paragraph{MCQ-1372 [hard]} What is the central engineering tradeoff of Quantization?
\begin{enumerate}[label=\Alph*.]
\item It enables large models and lower per-request latency but requires high-bandwidth interconnects.
\item Memory and bandwidth fall, while accuracy and kernel support can degrade.
\item Decode speed improves while cache memory grows with layers, sequence, batch, and KV heads.
\item It scales model depth but creates bubbles and per-request latency.
\end{enumerate}
\noindent\textbf{Answer.} Memory and bandwidth fall, while accuracy and kernel support can degrade.
\par\textit{Quantization maps weights or activations to lower precision using scales, groups, calibration, or quantization-aware training. Tradeoff: Memory and bandwidth fall, while accuracy and kernel support can degrade. Watch for: A smaller model file may run slower if hardware lacks optimized kernels.}
\paragraph{MCQ-1373 [hard]} Which failure mode should an interviewer expect you to identify for Quantization?
\begin{enumerate}[label=\Alph*.]
\item Capacity planning only model weights causes out-of-memory failures at long context.
\item A smaller model file may run slower if hardware lacks optimized kernels.
\item Spanning slow network links can make communication dominate compute.
\item Too few microbatches leave stages idle.
\end{enumerate}
\noindent\textbf{Answer.} A smaller model file may run slower if hardware lacks optimized kernels.
\par\textit{Quantization maps weights or activations to lower precision using scales, groups, calibration, or quantization-aware training. Tradeoff: Memory and bandwidth fall, while accuracy and kernel support can degrade. Watch for: A smaller model file may run slower if hardware lacks optimized kernels.}
\paragraph{MCQ-1374 [medium]} A design review is specifically concerned with this risk: A smaller model file may run slower if hardware lacks optimized kernels. Which topic should be reviewed first?
\begin{enumerate}[label=\Alph*.]
\item Quantization
\item Pipeline parallelism
\item Tensor parallelism
\item KV caching
\end{enumerate}
\noindent\textbf{Answer.} Quantization
\par\textit{Quantization maps weights or activations to lower precision using scales, groups, calibration, or quantization-aware training. Tradeoff: Memory and bandwidth fall, while accuracy and kernel support can degrade. Watch for: A smaller model file may run slower if hardware lacks optimized kernels.}
\paragraph{MCQ-1375 [hard]} Which design note most accurately balances the benefits and costs of Quantization?
\begin{enumerate}[label=\Alph*.]
\item Memory and bandwidth fall, while accuracy and kernel support can degrade.
\item Decode speed improves while cache memory grows with layers, sequence, batch, and KV heads.
\item It enables large models and lower per-request latency but requires high-bandwidth interconnects.
\item It scales model depth but creates bubbles and per-request latency.
\end{enumerate}
\noindent\textbf{Answer.} Memory and bandwidth fall, while accuracy and kernel support can degrade.
\par\textit{Quantization maps weights or activations to lower precision using scales, groups, calibration, or quantization-aware training. Tradeoff: Memory and bandwidth fall, while accuracy and kernel support can degrade. Watch for: A smaller model file may run slower if hardware lacks optimized kernels.}
\paragraph{MCQ-1376 [medium]} Which explanation would be strongest in a system-design interview about Quantization?
\begin{enumerate}[label=\Alph*.]
\item Quantization maps weights or activations to lower precision using scales, groups, calibration, or quantization-aware training.
\item Autoregressive decoding stores prior keys and values so each new token avoids recomputing attention over the entire prefix.
\item Pipeline parallelism assigns layer stages to devices and streams microbatches through them.
\item Tensor parallelism shards matrix operations across devices and communicates partial results within layers.
\end{enumerate}
\noindent\textbf{Answer.} Quantization maps weights or activations to lower precision using scales, groups, calibration, or quantization-aware training.
\par\textit{Quantization maps weights or activations to lower precision using scales, groups, calibration, or quantization-aware training. Tradeoff: Memory and bandwidth fall, while accuracy and kernel support can degrade. Watch for: A smaller model file may run slower if hardware lacks optimized kernels.}
\paragraph{MCQ-1377 [hard]} Which production warning is most directly associated with Quantization?
\begin{enumerate}[label=\Alph*.]
\item Spanning slow network links can make communication dominate compute.
\item A smaller model file may run slower if hardware lacks optimized kernels.
\item Capacity planning only model weights causes out-of-memory failures at long context.
\item Too few microbatches leave stages idle.
\end{enumerate}
\noindent\textbf{Answer.} A smaller model file may run slower if hardware lacks optimized kernels.
\par\textit{Quantization maps weights or activations to lower precision using scales, groups, calibration, or quantization-aware training. Tradeoff: Memory and bandwidth fall, while accuracy and kernel support can degrade. Watch for: A smaller model file may run slower if hardware lacks optimized kernels.}
\paragraph{FC-1378 [medium]} Define Quantization in interview-ready language.
\noindent\textbf{Answer.} Quantization maps weights or activations to lower precision using scales, groups, calibration, or quantization-aware training.
\par\textit{Quantization maps weights or activations to lower precision using scales, groups, calibration, or quantization-aware training.}
\paragraph{FC-1379 [hard]} State the main tradeoff and production pitfall for Quantization.
\noindent\textbf{Answer.} Tradeoff: Memory and bandwidth fall, while accuracy and kernel support can degrade. Pitfall: A smaller model file may run slower if hardware lacks optimized kernels.
\par\textit{Tradeoff: Memory and bandwidth fall, while accuracy and kernel support can degrade. Pitfall: A smaller model file may run slower if hardware lacks optimized kernels.}
\paragraph{LA-1380 [hard]} Explain Quantization from first principles, then show how you would evaluate and operate it in production.
\noindent\textbf{Answer.} Start with the mechanism: Quantization maps weights or activations to lower precision using scales, groups, calibration, or quantization-aware training. Make the tradeoff explicit: Memory and bandwidth fall, while accuracy and kernel support can degrade. Define workload-specific offline metrics and a latency/cost budget, test important slices, instrument the critical path, and create a rollback criterion. The production review must explicitly guard against this failure: A smaller model file may run slower if hardware lacks optimized kernels.
\par\textit{A strong answer connects mechanism, measurable tradeoffs, failure isolation, observability, and rollback rather than listing product features.}
\paragraph{MCQ-1381 [medium]} Which statement best describes Tensor parallelism?
\begin{enumerate}[label=\Alph*.]
\item Grouped-query and multi-query attention share key-value heads across query heads to shrink KV cache and memory bandwidth.
\item Replicas serve independent requests and scale throughput with routing and load balancing.
\item Tensor parallelism shards matrix operations across devices and communicates partial results within layers.
\item Autoregressive decoding stores prior keys and values so each new token avoids recomputing attention over the entire prefix.
\end{enumerate}
\noindent\textbf{Answer.} Tensor parallelism shards matrix operations across devices and communicates partial results within layers.
\par\textit{Tensor parallelism shards matrix operations across devices and communicates partial results within layers. Tradeoff: It enables large models and lower per-request latency but requires high-bandwidth interconnects. Watch for: Spanning slow network links can make communication dominate compute.}
\paragraph{MCQ-1382 [hard]} What is the central engineering tradeoff of Tensor parallelism?
\begin{enumerate}[label=\Alph*.]
\item It is operationally simple but duplicates model weights on every replica.
\item Decode efficiency improves, with possible quality loss relative to full multi-head KV.
\item Decode speed improves while cache memory grows with layers, sequence, batch, and KV heads.
\item It enables large models and lower per-request latency but requires high-bandwidth interconnects.
\end{enumerate}
\noindent\textbf{Answer.} It enables large models and lower per-request latency but requires high-bandwidth interconnects.
\par\textit{Tensor parallelism shards matrix operations across devices and communicates partial results within layers. Tradeoff: It enables large models and lower per-request latency but requires high-bandwidth interconnects. Watch for: Spanning slow network links can make communication dominate compute.}
\paragraph{MCQ-1383 [hard]} Which failure mode should an interviewer expect you to identify for Tensor parallelism?
\begin{enumerate}[label=\Alph*.]
\item Random routing ignores prefix-cache locality and uneven sequence lengths.
\item Converting weights naively after training does not reproduce a model trained for shared KV heads.
\item Spanning slow network links can make communication dominate compute.
\item Capacity planning only model weights causes out-of-memory failures at long context.
\end{enumerate}
\noindent\textbf{Answer.} Spanning slow network links can make communication dominate compute.
\par\textit{Tensor parallelism shards matrix operations across devices and communicates partial results within layers. Tradeoff: It enables large models and lower per-request latency but requires high-bandwidth interconnects. Watch for: Spanning slow network links can make communication dominate compute.}
\paragraph{MCQ-1384 [medium]} A design review is specifically concerned with this risk: Spanning slow network links can make communication dominate compute. Which topic should be reviewed first?
\begin{enumerate}[label=\Alph*.]
\item Data parallel inference
\item Tensor parallelism
\item GQA and MQA
\item KV caching
\end{enumerate}
\noindent\textbf{Answer.} Tensor parallelism
\par\textit{Tensor parallelism shards matrix operations across devices and communicates partial results within layers. Tradeoff: It enables large models and lower per-request latency but requires high-bandwidth interconnects. Watch for: Spanning slow network links can make communication dominate compute.}
\paragraph{MCQ-1385 [hard]} Which design note most accurately balances the benefits and costs of Tensor parallelism?
\begin{enumerate}[label=\Alph*.]
\item It is operationally simple but duplicates model weights on every replica.
\item It enables large models and lower per-request latency but requires high-bandwidth interconnects.
\item Decode speed improves while cache memory grows with layers, sequence, batch, and KV heads.
\item Decode efficiency improves, with possible quality loss relative to full multi-head KV.
\end{enumerate}
\noindent\textbf{Answer.} It enables large models and lower per-request latency but requires high-bandwidth interconnects.
\par\textit{Tensor parallelism shards matrix operations across devices and communicates partial results within layers. Tradeoff: It enables large models and lower per-request latency but requires high-bandwidth interconnects. Watch for: Spanning slow network links can make communication dominate compute.}
\paragraph{MCQ-1386 [medium]} Which explanation would be strongest in a system-design interview about Tensor parallelism?
\begin{enumerate}[label=\Alph*.]
\item Tensor parallelism shards matrix operations across devices and communicates partial results within layers.
\item Autoregressive decoding stores prior keys and values so each new token avoids recomputing attention over the entire prefix.
\item Grouped-query and multi-query attention share key-value heads across query heads to shrink KV cache and memory bandwidth.
\item Replicas serve independent requests and scale throughput with routing and load balancing.
\end{enumerate}
\noindent\textbf{Answer.} Tensor parallelism shards matrix operations across devices and communicates partial results within layers.
\par\textit{Tensor parallelism shards matrix operations across devices and communicates partial results within layers. Tradeoff: It enables large models and lower per-request latency but requires high-bandwidth interconnects. Watch for: Spanning slow network links can make communication dominate compute.}
\paragraph{MCQ-1387 [hard]} Which production warning is most directly associated with Tensor parallelism?
\begin{enumerate}[label=\Alph*.]
\item Random routing ignores prefix-cache locality and uneven sequence lengths.
\item Capacity planning only model weights causes out-of-memory failures at long context.
\item Spanning slow network links can make communication dominate compute.
\item Converting weights naively after training does not reproduce a model trained for shared KV heads.
\end{enumerate}
\noindent\textbf{Answer.} Spanning slow network links can make communication dominate compute.
\par\textit{Tensor parallelism shards matrix operations across devices and communicates partial results within layers. Tradeoff: It enables large models and lower per-request latency but requires high-bandwidth interconnects. Watch for: Spanning slow network links can make communication dominate compute.}
\paragraph{FC-1388 [medium]} Define Tensor parallelism in interview-ready language.
\noindent\textbf{Answer.} Tensor parallelism shards matrix operations across devices and communicates partial results within layers.
\par\textit{Tensor parallelism shards matrix operations across devices and communicates partial results within layers.}
\paragraph{FC-1389 [hard]} State the main tradeoff and production pitfall for Tensor parallelism.
\noindent\textbf{Answer.} Tradeoff: It enables large models and lower per-request latency but requires high-bandwidth interconnects. Pitfall: Spanning slow network links can make communication dominate compute.
\par\textit{Tradeoff: It enables large models and lower per-request latency but requires high-bandwidth interconnects. Pitfall: Spanning slow network links can make communication dominate compute.}
\paragraph{LA-1390 [hard]} Explain Tensor parallelism from first principles, then show how you would evaluate and operate it in production.
\noindent\textbf{Answer.} Start with the mechanism: Tensor parallelism shards matrix operations across devices and communicates partial results within layers. Make the tradeoff explicit: It enables large models and lower per-request latency but requires high-bandwidth interconnects. Define workload-specific offline metrics and a latency/cost budget, test important slices, instrument the critical path, and create a rollback criterion. The production review must explicitly guard against this failure: Spanning slow network links can make communication dominate compute.
\par\textit{A strong answer connects mechanism, measurable tradeoffs, failure isolation, observability, and rollback rather than listing product features.}
\paragraph{MCQ-1391 [medium]} Which statement best describes Pipeline parallelism?
\begin{enumerate}[label=\Alph*.]
\item Shared prompt prefixes reuse precomputed KV blocks across requests with matching tokens and model state.
\item Pipeline parallelism assigns layer stages to devices and streams microbatches through them.
\item Grouped-query and multi-query attention share key-value heads across query heads to shrink KV cache and memory bandwidth.
\item Replicas serve independent requests and scale throughput with routing and load balancing.
\end{enumerate}
\noindent\textbf{Answer.} Pipeline parallelism assigns layer stages to devices and streams microbatches through them.
\par\textit{Pipeline parallelism assigns layer stages to devices and streams microbatches through them. Tradeoff: It scales model depth but creates bubbles and per-request latency. Watch for: Too few microbatches leave stages idle.}
\paragraph{MCQ-1392 [hard]} What is the central engineering tradeoff of Pipeline parallelism?
\begin{enumerate}[label=\Alph*.]
\item Repeated system prompts become cheaper, but hit rate, invalidation, and tenant isolation matter.
\item It scales model depth but creates bubbles and per-request latency.
\item Decode efficiency improves, with possible quality loss relative to full multi-head KV.
\item It is operationally simple but duplicates model weights on every replica.
\end{enumerate}
\noindent\textbf{Answer.} It scales model depth but creates bubbles and per-request latency.
\par\textit{Pipeline parallelism assigns layer stages to devices and streams microbatches through them. Tradeoff: It scales model depth but creates bubbles and per-request latency. Watch for: Too few microbatches leave stages idle.}
\paragraph{MCQ-1393 [hard]} Which failure mode should an interviewer expect you to identify for Pipeline parallelism?
\begin{enumerate}[label=\Alph*.]
\item Too few microbatches leave stages idle.
\item Random routing ignores prefix-cache locality and uneven sequence lengths.
\item Converting weights naively after training does not reproduce a model trained for shared KV heads.
\item Cache keys that omit adapters or model revision return incorrect state.
\end{enumerate}
\noindent\textbf{Answer.} Too few microbatches leave stages idle.
\par\textit{Pipeline parallelism assigns layer stages to devices and streams microbatches through them. Tradeoff: It scales model depth but creates bubbles and per-request latency. Watch for: Too few microbatches leave stages idle.}
\paragraph{MCQ-1394 [medium]} A design review is specifically concerned with this risk: Too few microbatches leave stages idle. Which topic should be reviewed first?
\begin{enumerate}[label=\Alph*.]
\item Data parallel inference
\item Pipeline parallelism
\item GQA and MQA
\item Prefix caching
\end{enumerate}
\noindent\textbf{Answer.} Pipeline parallelism
\par\textit{Pipeline parallelism assigns layer stages to devices and streams microbatches through them. Tradeoff: It scales model depth but creates bubbles and per-request latency. Watch for: Too few microbatches leave stages idle.}
\paragraph{MCQ-1395 [hard]} Which design note most accurately balances the benefits and costs of Pipeline parallelism?
\begin{enumerate}[label=\Alph*.]
\item Decode efficiency improves, with possible quality loss relative to full multi-head KV.
\item It is operationally simple but duplicates model weights on every replica.
\item Repeated system prompts become cheaper, but hit rate, invalidation, and tenant isolation matter.
\item It scales model depth but creates bubbles and per-request latency.
\end{enumerate}
\noindent\textbf{Answer.} It scales model depth but creates bubbles and per-request latency.
\par\textit{Pipeline parallelism assigns layer stages to devices and streams microbatches through them. Tradeoff: It scales model depth but creates bubbles and per-request latency. Watch for: Too few microbatches leave stages idle.}
\paragraph{MCQ-1396 [medium]} Which explanation would be strongest in a system-design interview about Pipeline parallelism?
\begin{enumerate}[label=\Alph*.]
\item Replicas serve independent requests and scale throughput with routing and load balancing.
\item Shared prompt prefixes reuse precomputed KV blocks across requests with matching tokens and model state.
\item Pipeline parallelism assigns layer stages to devices and streams microbatches through them.
\item Grouped-query and multi-query attention share key-value heads across query heads to shrink KV cache and memory bandwidth.
\end{enumerate}
\noindent\textbf{Answer.} Pipeline parallelism assigns layer stages to devices and streams microbatches through them.
\par\textit{Pipeline parallelism assigns layer stages to devices and streams microbatches through them. Tradeoff: It scales model depth but creates bubbles and per-request latency. Watch for: Too few microbatches leave stages idle.}
\paragraph{MCQ-1397 [hard]} Which production warning is most directly associated with Pipeline parallelism?
\begin{enumerate}[label=\Alph*.]
\item Random routing ignores prefix-cache locality and uneven sequence lengths.
\item Cache keys that omit adapters or model revision return incorrect state.
\item Converting weights naively after training does not reproduce a model trained for shared KV heads.
\item Too few microbatches leave stages idle.
\end{enumerate}
\noindent\textbf{Answer.} Too few microbatches leave stages idle.
\par\textit{Pipeline parallelism assigns layer stages to devices and streams microbatches through them. Tradeoff: It scales model depth but creates bubbles and per-request latency. Watch for: Too few microbatches leave stages idle.}
\paragraph{FC-1398 [medium]} Define Pipeline parallelism in interview-ready language.
\noindent\textbf{Answer.} Pipeline parallelism assigns layer stages to devices and streams microbatches through them.
\par\textit{Pipeline parallelism assigns layer stages to devices and streams microbatches through them.}
\paragraph{FC-1399 [hard]} State the main tradeoff and production pitfall for Pipeline parallelism.
\noindent\textbf{Answer.} Tradeoff: It scales model depth but creates bubbles and per-request latency. Pitfall: Too few microbatches leave stages idle.
\par\textit{Tradeoff: It scales model depth but creates bubbles and per-request latency. Pitfall: Too few microbatches leave stages idle.}
\paragraph{LA-1400 [hard]} Explain Pipeline parallelism from first principles, then show how you would evaluate and operate it in production.
\noindent\textbf{Answer.} Start with the mechanism: Pipeline parallelism assigns layer stages to devices and streams microbatches through them. Make the tradeoff explicit: It scales model depth but creates bubbles and per-request latency. Define workload-specific offline metrics and a latency/cost budget, test important slices, instrument the critical path, and create a rollback criterion. The production review must explicitly guard against this failure: Too few microbatches leave stages idle.
\par\textit{A strong answer connects mechanism, measurable tradeoffs, failure isolation, observability, and rollback rather than listing product features.}
\paragraph{MCQ-1401 [medium]} Which statement best describes Data parallel inference?
\begin{enumerate}[label=\Alph*.]
\item Requests join and leave a running decode batch at iteration boundaries rather than waiting for a fixed batch.
\item Replicas serve independent requests and scale throughput with routing and load balancing.
\item Grouped-query and multi-query attention share key-value heads across query heads to shrink KV cache and memory bandwidth.
\item Servers emit tokens incrementally with backpressure, disconnect handling, and final usage accounting.
\end{enumerate}
\noindent\textbf{Answer.} Replicas serve independent requests and scale throughput with routing and load balancing.
\par\textit{Replicas serve independent requests and scale throughput with routing and load balancing. Tradeoff: It is operationally simple but duplicates model weights on every replica. Watch for: Random routing ignores prefix-cache locality and uneven sequence lengths.}
\paragraph{MCQ-1402 [hard]} What is the central engineering tradeoff of Data parallel inference?
\begin{enumerate}[label=\Alph*.]
\item Perceived latency improves, but partial outputs complicate moderation and retries.
\item Throughput and utilization rise, while scheduling fairness and latency isolation become harder.
\item Decode efficiency improves, with possible quality loss relative to full multi-head KV.
\item It is operationally simple but duplicates model weights on every replica.
\end{enumerate}
\noindent\textbf{Answer.} It is operationally simple but duplicates model weights on every replica.
\par\textit{Replicas serve independent requests and scale throughput with routing and load balancing. Tradeoff: It is operationally simple but duplicates model weights on every replica. Watch for: Random routing ignores prefix-cache locality and uneven sequence lengths.}
\paragraph{MCQ-1403 [hard]} Which failure mode should an interviewer expect you to identify for Data parallel inference?
\begin{enumerate}[label=\Alph*.]
\item Converting weights naively after training does not reproduce a model trained for shared KV heads.
\item Retrying after visible partial output can duplicate text or side effects.
\item Optimizing total tokens per second alone can starve short requests.
\item Random routing ignores prefix-cache locality and uneven sequence lengths.
\end{enumerate}
\noindent\textbf{Answer.} Random routing ignores prefix-cache locality and uneven sequence lengths.
\par\textit{Replicas serve independent requests and scale throughput with routing and load balancing. Tradeoff: It is operationally simple but duplicates model weights on every replica. Watch for: Random routing ignores prefix-cache locality and uneven sequence lengths.}
\paragraph{MCQ-1404 [medium]} A design review is specifically concerned with this risk: Random routing ignores prefix-cache locality and uneven sequence lengths. Which topic should be reviewed first?
\begin{enumerate}[label=\Alph*.]
\item GQA and MQA
\item Continuous batching
\item Streaming generation
\item Data parallel inference
\end{enumerate}
\noindent\textbf{Answer.} Data parallel inference
\par\textit{Replicas serve independent requests and scale throughput with routing and load balancing. Tradeoff: It is operationally simple but duplicates model weights on every replica. Watch for: Random routing ignores prefix-cache locality and uneven sequence lengths.}
\paragraph{MCQ-1405 [hard]} Which design note most accurately balances the benefits and costs of Data parallel inference?
\begin{enumerate}[label=\Alph*.]
\item Perceived latency improves, but partial outputs complicate moderation and retries.
\item Throughput and utilization rise, while scheduling fairness and latency isolation become harder.
\item Decode efficiency improves, with possible quality loss relative to full multi-head KV.
\item It is operationally simple but duplicates model weights on every replica.
\end{enumerate}
\noindent\textbf{Answer.} It is operationally simple but duplicates model weights on every replica.
\par\textit{Replicas serve independent requests and scale throughput with routing and load balancing. Tradeoff: It is operationally simple but duplicates model weights on every replica. Watch for: Random routing ignores prefix-cache locality and uneven sequence lengths.}
\paragraph{MCQ-1406 [medium]} Which explanation would be strongest in a system-design interview about Data parallel inference?
\begin{enumerate}[label=\Alph*.]
\item Grouped-query and multi-query attention share key-value heads across query heads to shrink KV cache and memory bandwidth.
\item Servers emit tokens incrementally with backpressure, disconnect handling, and final usage accounting.
\item Replicas serve independent requests and scale throughput with routing and load balancing.
\item Requests join and leave a running decode batch at iteration boundaries rather than waiting for a fixed batch.
\end{enumerate}
\noindent\textbf{Answer.} Replicas serve independent requests and scale throughput with routing and load balancing.
\par\textit{Replicas serve independent requests and scale throughput with routing and load balancing. Tradeoff: It is operationally simple but duplicates model weights on every replica. Watch for: Random routing ignores prefix-cache locality and uneven sequence lengths.}
\paragraph{MCQ-1407 [hard]} Which production warning is most directly associated with Data parallel inference?
\begin{enumerate}[label=\Alph*.]
\item Random routing ignores prefix-cache locality and uneven sequence lengths.
\item Converting weights naively after training does not reproduce a model trained for shared KV heads.
\item Optimizing total tokens per second alone can starve short requests.
\item Retrying after visible partial output can duplicate text or side effects.
\end{enumerate}
\noindent\textbf{Answer.} Random routing ignores prefix-cache locality and uneven sequence lengths.
\par\textit{Replicas serve independent requests and scale throughput with routing and load balancing. Tradeoff: It is operationally simple but duplicates model weights on every replica. Watch for: Random routing ignores prefix-cache locality and uneven sequence lengths.}
\paragraph{FC-1408 [medium]} Define Data parallel inference in interview-ready language.
\noindent\textbf{Answer.} Replicas serve independent requests and scale throughput with routing and load balancing.
\par\textit{Replicas serve independent requests and scale throughput with routing and load balancing.}
\paragraph{FC-1409 [hard]} State the main tradeoff and production pitfall for Data parallel inference.
\noindent\textbf{Answer.} Tradeoff: It is operationally simple but duplicates model weights on every replica. Pitfall: Random routing ignores prefix-cache locality and uneven sequence lengths.
\par\textit{Tradeoff: It is operationally simple but duplicates model weights on every replica. Pitfall: Random routing ignores prefix-cache locality and uneven sequence lengths.}
\paragraph{LA-1410 [hard]} Explain Data parallel inference from first principles, then show how you would evaluate and operate it in production.
\noindent\textbf{Answer.} Start with the mechanism: Replicas serve independent requests and scale throughput with routing and load balancing. Make the tradeoff explicit: It is operationally simple but duplicates model weights on every replica. Define workload-specific offline metrics and a latency/cost budget, test important slices, instrument the critical path, and create a rollback criterion. The production review must explicitly guard against this failure: Random routing ignores prefix-cache locality and uneven sequence lengths.
\par\textit{A strong answer connects mechanism, measurable tradeoffs, failure isolation, observability, and rollback rather than listing product features.}
\paragraph{MCQ-1411 [medium]} Which statement best describes Expert parallelism?
\begin{enumerate}[label=\Alph*.]
\item Replicas serve independent requests and scale throughput with routing and load balancing.
\item Requests join and leave a running decode batch at iteration boundaries rather than waiting for a fixed batch.
\item MoE experts are distributed across devices and tokens are routed with all-to-all communication.
\item Grouped-query and multi-query attention share key-value heads across query heads to shrink KV cache and memory bandwidth.
\end{enumerate}
\noindent\textbf{Answer.} MoE experts are distributed across devices and tokens are routed with all-to-all communication.
\par\textit{MoE experts are distributed across devices and tokens are routed with all-to-all communication. Tradeoff: Parameter capacity scales, while communication and load balance determine performance. Watch for: Average expert load hides hot experts that gate step time.}
\paragraph{MCQ-1412 [hard]} What is the central engineering tradeoff of Expert parallelism?
\begin{enumerate}[label=\Alph*.]
\item Parameter capacity scales, while communication and load balance determine performance.
\item Decode efficiency improves, with possible quality loss relative to full multi-head KV.
\item It is operationally simple but duplicates model weights on every replica.
\item Throughput and utilization rise, while scheduling fairness and latency isolation become harder.
\end{enumerate}
\noindent\textbf{Answer.} Parameter capacity scales, while communication and load balance determine performance.
\par\textit{MoE experts are distributed across devices and tokens are routed with all-to-all communication. Tradeoff: Parameter capacity scales, while communication and load balance determine performance. Watch for: Average expert load hides hot experts that gate step time.}
\paragraph{MCQ-1413 [hard]} Which failure mode should an interviewer expect you to identify for Expert parallelism?
\begin{enumerate}[label=\Alph*.]
\item Optimizing total tokens per second alone can starve short requests.
\item Converting weights naively after training does not reproduce a model trained for shared KV heads.
\item Average expert load hides hot experts that gate step time.
\item Random routing ignores prefix-cache locality and uneven sequence lengths.
\end{enumerate}
\noindent\textbf{Answer.} Average expert load hides hot experts that gate step time.
\par\textit{MoE experts are distributed across devices and tokens are routed with all-to-all communication. Tradeoff: Parameter capacity scales, while communication and load balance determine performance. Watch for: Average expert load hides hot experts that gate step time.}
\paragraph{MCQ-1414 [medium]} A design review is specifically concerned with this risk: Average expert load hides hot experts that gate step time. Which topic should be reviewed first?
\begin{enumerate}[label=\Alph*.]
\item Data parallel inference
\item Expert parallelism
\item GQA and MQA
\item Continuous batching
\end{enumerate}
\noindent\textbf{Answer.} Expert parallelism
\par\textit{MoE experts are distributed across devices and tokens are routed with all-to-all communication. Tradeoff: Parameter capacity scales, while communication and load balance determine performance. Watch for: Average expert load hides hot experts that gate step time.}
\paragraph{MCQ-1415 [hard]} Which design note most accurately balances the benefits and costs of Expert parallelism?
\begin{enumerate}[label=\Alph*.]
\item Parameter capacity scales, while communication and load balance determine performance.
\item It is operationally simple but duplicates model weights on every replica.
\item Decode efficiency improves, with possible quality loss relative to full multi-head KV.
\item Throughput and utilization rise, while scheduling fairness and latency isolation become harder.
\end{enumerate}
\noindent\textbf{Answer.} Parameter capacity scales, while communication and load balance determine performance.
\par\textit{MoE experts are distributed across devices and tokens are routed with all-to-all communication. Tradeoff: Parameter capacity scales, while communication and load balance determine performance. Watch for: Average expert load hides hot experts that gate step time.}
\paragraph{MCQ-1416 [medium]} Which explanation would be strongest in a system-design interview about Expert parallelism?
\begin{enumerate}[label=\Alph*.]
\item Replicas serve independent requests and scale throughput with routing and load balancing.
\item MoE experts are distributed across devices and tokens are routed with all-to-all communication.
\item Grouped-query and multi-query attention share key-value heads across query heads to shrink KV cache and memory bandwidth.
\item Requests join and leave a running decode batch at iteration boundaries rather than waiting for a fixed batch.
\end{enumerate}
\noindent\textbf{Answer.} MoE experts are distributed across devices and tokens are routed with all-to-all communication.
\par\textit{MoE experts are distributed across devices and tokens are routed with all-to-all communication. Tradeoff: Parameter capacity scales, while communication and load balance determine performance. Watch for: Average expert load hides hot experts that gate step time.}
\paragraph{MCQ-1417 [hard]} Which production warning is most directly associated with Expert parallelism?
\begin{enumerate}[label=\Alph*.]
\item Converting weights naively after training does not reproduce a model trained for shared KV heads.
\item Optimizing total tokens per second alone can starve short requests.
\item Average expert load hides hot experts that gate step time.
\item Random routing ignores prefix-cache locality and uneven sequence lengths.
\end{enumerate}
\noindent\textbf{Answer.} Average expert load hides hot experts that gate step time.
\par\textit{MoE experts are distributed across devices and tokens are routed with all-to-all communication. Tradeoff: Parameter capacity scales, while communication and load balance determine performance. Watch for: Average expert load hides hot experts that gate step time.}
\paragraph{FC-1418 [medium]} Define Expert parallelism in interview-ready language.
\noindent\textbf{Answer.} MoE experts are distributed across devices and tokens are routed with all-to-all communication.
\par\textit{MoE experts are distributed across devices and tokens are routed with all-to-all communication.}
\paragraph{FC-1419 [hard]} State the main tradeoff and production pitfall for Expert parallelism.
\noindent\textbf{Answer.} Tradeoff: Parameter capacity scales, while communication and load balance determine performance. Pitfall: Average expert load hides hot experts that gate step time.
\par\textit{Tradeoff: Parameter capacity scales, while communication and load balance determine performance. Pitfall: Average expert load hides hot experts that gate step time.}
\paragraph{LA-1420 [hard]} Explain Expert parallelism from first principles, then show how you would evaluate and operate it in production.
\noindent\textbf{Answer.} Start with the mechanism: MoE experts are distributed across devices and tokens are routed with all-to-all communication. Make the tradeoff explicit: Parameter capacity scales, while communication and load balance determine performance. Define workload-specific offline metrics and a latency/cost budget, test important slices, instrument the critical path, and create a rollback criterion. The production review must explicitly guard against this failure: Average expert load hides hot experts that gate step time.
\par\textit{A strong answer connects mechanism, measurable tradeoffs, failure isolation, observability, and rollback rather than listing product features.}
\paragraph{MCQ-1421 [medium]} Which statement best describes Dynamic batching?
\begin{enumerate}[label=\Alph*.]
\item Shared prompt prefixes reuse precomputed KV blocks across requests with matching tokens and model state.
\item A server waits briefly to combine compatible requests into a batch.
\item MoE experts are distributed across devices and tokens are routed with all-to-all communication.
\item PagedAttention manages KV cache in non-contiguous blocks using paging-inspired allocation to reduce fragmentation and enable sharing.
\end{enumerate}
\noindent\textbf{Answer.} A server waits briefly to combine compatible requests into a batch.
\par\textit{A server waits briefly to combine compatible requests into a batch. Tradeoff: GPU efficiency improves while queue delay harms time-to-first-token. Watch for: One global maximum delay cannot serve both interactive and batch workloads well.}
\paragraph{MCQ-1422 [hard]} What is the central engineering tradeoff of Dynamic batching?
\begin{enumerate}[label=\Alph*.]
\item Repeated system prompts become cheaper, but hit rate, invalidation, and tenant isolation matter.
\item Parameter capacity scales, while communication and load balance determine performance.
\item Utilization and batching improve, while block tables and kernels add runtime complexity.
\item GPU efficiency improves while queue delay harms time-to-first-token.
\end{enumerate}
\noindent\textbf{Answer.} GPU efficiency improves while queue delay harms time-to-first-token.
\par\textit{A server waits briefly to combine compatible requests into a batch. Tradeoff: GPU efficiency improves while queue delay harms time-to-first-token. Watch for: One global maximum delay cannot serve both interactive and batch workloads well.}
\paragraph{MCQ-1423 [hard]} Which failure mode should an interviewer expect you to identify for Dynamic batching?
\begin{enumerate}[label=\Alph*.]
\item Oversized blocks waste tail space; tiny blocks add metadata overhead.
\item Average expert load hides hot experts that gate step time.
\item One global maximum delay cannot serve both interactive and batch workloads well.
\item Cache keys that omit adapters or model revision return incorrect state.
\end{enumerate}
\noindent\textbf{Answer.} One global maximum delay cannot serve both interactive and batch workloads well.
\par\textit{A server waits briefly to combine compatible requests into a batch. Tradeoff: GPU efficiency improves while queue delay harms time-to-first-token. Watch for: One global maximum delay cannot serve both interactive and batch workloads well.}
\paragraph{MCQ-1424 [medium]} A design review is specifically concerned with this risk: One global maximum delay cannot serve both interactive and batch workloads well. Which topic should be reviewed first?
\begin{enumerate}[label=\Alph*.]
\item Prefix caching
\item PagedAttention
\item Dynamic batching
\item Expert parallelism
\end{enumerate}
\noindent\textbf{Answer.} Dynamic batching
\par\textit{A server waits briefly to combine compatible requests into a batch. Tradeoff: GPU efficiency improves while queue delay harms time-to-first-token. Watch for: One global maximum delay cannot serve both interactive and batch workloads well.}
\paragraph{MCQ-1425 [hard]} Which design note most accurately balances the benefits and costs of Dynamic batching?
\begin{enumerate}[label=\Alph*.]
\item Parameter capacity scales, while communication and load balance determine performance.
\item GPU efficiency improves while queue delay harms time-to-first-token.
\item Repeated system prompts become cheaper, but hit rate, invalidation, and tenant isolation matter.
\item Utilization and batching improve, while block tables and kernels add runtime complexity.
\end{enumerate}
\noindent\textbf{Answer.} GPU efficiency improves while queue delay harms time-to-first-token.
\par\textit{A server waits briefly to combine compatible requests into a batch. Tradeoff: GPU efficiency improves while queue delay harms time-to-first-token. Watch for: One global maximum delay cannot serve both interactive and batch workloads well.}
\paragraph{MCQ-1426 [medium]} Which explanation would be strongest in a system-design interview about Dynamic batching?
\begin{enumerate}[label=\Alph*.]
\item PagedAttention manages KV cache in non-contiguous blocks using paging-inspired allocation to reduce fragmentation and enable sharing.
\item Shared prompt prefixes reuse precomputed KV blocks across requests with matching tokens and model state.
\item MoE experts are distributed across devices and tokens are routed with all-to-all communication.
\item A server waits briefly to combine compatible requests into a batch.
\end{enumerate}
\noindent\textbf{Answer.} A server waits briefly to combine compatible requests into a batch.
\par\textit{A server waits briefly to combine compatible requests into a batch. Tradeoff: GPU efficiency improves while queue delay harms time-to-first-token. Watch for: One global maximum delay cannot serve both interactive and batch workloads well.}
\paragraph{MCQ-1427 [hard]} Which production warning is most directly associated with Dynamic batching?
\begin{enumerate}[label=\Alph*.]
\item One global maximum delay cannot serve both interactive and batch workloads well.
\item Cache keys that omit adapters or model revision return incorrect state.
\item Oversized blocks waste tail space; tiny blocks add metadata overhead.
\item Average expert load hides hot experts that gate step time.
\end{enumerate}
\noindent\textbf{Answer.} One global maximum delay cannot serve both interactive and batch workloads well.
\par\textit{A server waits briefly to combine compatible requests into a batch. Tradeoff: GPU efficiency improves while queue delay harms time-to-first-token. Watch for: One global maximum delay cannot serve both interactive and batch workloads well.}
\paragraph{FC-1428 [medium]} Define Dynamic batching in interview-ready language.
\noindent\textbf{Answer.} A server waits briefly to combine compatible requests into a batch.
\par\textit{A server waits briefly to combine compatible requests into a batch.}
\paragraph{FC-1429 [hard]} State the main tradeoff and production pitfall for Dynamic batching.
\noindent\textbf{Answer.} Tradeoff: GPU efficiency improves while queue delay harms time-to-first-token. Pitfall: One global maximum delay cannot serve both interactive and batch workloads well.
\par\textit{Tradeoff: GPU efficiency improves while queue delay harms time-to-first-token. Pitfall: One global maximum delay cannot serve both interactive and batch workloads well.}
\paragraph{LA-1430 [hard]} Explain Dynamic batching from first principles, then show how you would evaluate and operate it in production.
\noindent\textbf{Answer.} Start with the mechanism: A server waits briefly to combine compatible requests into a batch. Make the tradeoff explicit: GPU efficiency improves while queue delay harms time-to-first-token. Define workload-specific offline metrics and a latency/cost budget, test important slices, instrument the critical path, and create a rollback criterion. The production review must explicitly guard against this failure: One global maximum delay cannot serve both interactive and batch workloads well.
\par\textit{A strong answer connects mechanism, measurable tradeoffs, failure isolation, observability, and rollback rather than listing product features.}
\paragraph{MCQ-1431 [medium]} Which statement best describes Streaming generation?
\begin{enumerate}[label=\Alph*.]
\item Servers emit tokens incrementally with backpressure, disconnect handling, and final usage accounting.
\item A cheaper draft proposes tokens and a target model verifies them while preserving the target distribution.
\item Requests join and leave a running decode batch at iteration boundaries rather than waiting for a fixed batch.
\item Quantization maps weights or activations to lower precision using scales, groups, calibration, or quantization-aware training.
\end{enumerate}
\noindent\textbf{Answer.} Servers emit tokens incrementally with backpressure, disconnect handling, and final usage accounting.
\par\textit{Servers emit tokens incrementally with backpressure, disconnect handling, and final usage accounting. Tradeoff: Perceived latency improves, but partial outputs complicate moderation and retries. Watch for: Retrying after visible partial output can duplicate text or side effects.}
\paragraph{MCQ-1432 [hard]} What is the central engineering tradeoff of Streaming generation?
\begin{enumerate}[label=\Alph*.]
\item Latency falls when acceptance is high, but drafting and verification overhead can dominate.
\item Memory and bandwidth fall, while accuracy and kernel support can degrade.
\item Throughput and utilization rise, while scheduling fairness and latency isolation become harder.
\item Perceived latency improves, but partial outputs complicate moderation and retries.
\end{enumerate}
\noindent\textbf{Answer.} Perceived latency improves, but partial outputs complicate moderation and retries.
\par\textit{Servers emit tokens incrementally with backpressure, disconnect handling, and final usage accounting. Tradeoff: Perceived latency improves, but partial outputs complicate moderation and retries. Watch for: Retrying after visible partial output can duplicate text or side effects.}
\paragraph{MCQ-1433 [hard]} Which failure mode should an interviewer expect you to identify for Streaming generation?
\begin{enumerate}[label=\Alph*.]
\item Measuring draft speed without acceptance rate gives meaningless projections.
\item Retrying after visible partial output can duplicate text or side effects.
\item A smaller model file may run slower if hardware lacks optimized kernels.
\item Optimizing total tokens per second alone can starve short requests.
\end{enumerate}
\noindent\textbf{Answer.} Retrying after visible partial output can duplicate text or side effects.
\par\textit{Servers emit tokens incrementally with backpressure, disconnect handling, and final usage accounting. Tradeoff: Perceived latency improves, but partial outputs complicate moderation and retries. Watch for: Retrying after visible partial output can duplicate text or side effects.}
\paragraph{MCQ-1434 [medium]} A design review is specifically concerned with this risk: Retrying after visible partial output can duplicate text or side effects. Which topic should be reviewed first?
\begin{enumerate}[label=\Alph*.]
\item Speculative decoding
\item Quantization
\item Continuous batching
\item Streaming generation
\end{enumerate}
\noindent\textbf{Answer.} Streaming generation
\par\textit{Servers emit tokens incrementally with backpressure, disconnect handling, and final usage accounting. Tradeoff: Perceived latency improves, but partial outputs complicate moderation and retries. Watch for: Retrying after visible partial output can duplicate text or side effects.}
\paragraph{MCQ-1435 [hard]} Which design note most accurately balances the benefits and costs of Streaming generation?
\begin{enumerate}[label=\Alph*.]
\item Throughput and utilization rise, while scheduling fairness and latency isolation become harder.
\item Perceived latency improves, but partial outputs complicate moderation and retries.
\item Memory and bandwidth fall, while accuracy and kernel support can degrade.
\item Latency falls when acceptance is high, but drafting and verification overhead can dominate.
\end{enumerate}
\noindent\textbf{Answer.} Perceived latency improves, but partial outputs complicate moderation and retries.
\par\textit{Servers emit tokens incrementally with backpressure, disconnect handling, and final usage accounting. Tradeoff: Perceived latency improves, but partial outputs complicate moderation and retries. Watch for: Retrying after visible partial output can duplicate text or side effects.}
\paragraph{MCQ-1436 [medium]} Which explanation would be strongest in a system-design interview about Streaming generation?
\begin{enumerate}[label=\Alph*.]
\item A cheaper draft proposes tokens and a target model verifies them while preserving the target distribution.
\item Requests join and leave a running decode batch at iteration boundaries rather than waiting for a fixed batch.
\item Quantization maps weights or activations to lower precision using scales, groups, calibration, or quantization-aware training.
\item Servers emit tokens incrementally with backpressure, disconnect handling, and final usage accounting.
\end{enumerate}
\noindent\textbf{Answer.} Servers emit tokens incrementally with backpressure, disconnect handling, and final usage accounting.
\par\textit{Servers emit tokens incrementally with backpressure, disconnect handling, and final usage accounting. Tradeoff: Perceived latency improves, but partial outputs complicate moderation and retries. Watch for: Retrying after visible partial output can duplicate text or side effects.}
\paragraph{MCQ-1437 [hard]} Which production warning is most directly associated with Streaming generation?
\begin{enumerate}[label=\Alph*.]
\item A smaller model file may run slower if hardware lacks optimized kernels.
\item Measuring draft speed without acceptance rate gives meaningless projections.
\item Retrying after visible partial output can duplicate text or side effects.
\item Optimizing total tokens per second alone can starve short requests.
\end{enumerate}
\noindent\textbf{Answer.} Retrying after visible partial output can duplicate text or side effects.
\par\textit{Servers emit tokens incrementally with backpressure, disconnect handling, and final usage accounting. Tradeoff: Perceived latency improves, but partial outputs complicate moderation and retries. Watch for: Retrying after visible partial output can duplicate text or side effects.}
\paragraph{FC-1438 [medium]} Define Streaming generation in interview-ready language.
\noindent\textbf{Answer.} Servers emit tokens incrementally with backpressure, disconnect handling, and final usage accounting.
\par\textit{Servers emit tokens incrementally with backpressure, disconnect handling, and final usage accounting.}
\paragraph{FC-1439 [hard]} State the main tradeoff and production pitfall for Streaming generation.
\noindent\textbf{Answer.} Tradeoff: Perceived latency improves, but partial outputs complicate moderation and retries. Pitfall: Retrying after visible partial output can duplicate text or side effects.
\par\textit{Tradeoff: Perceived latency improves, but partial outputs complicate moderation and retries. Pitfall: Retrying after visible partial output can duplicate text or side effects.}
\paragraph{LA-1440 [hard]} Explain Streaming generation from first principles, then show how you would evaluate and operate it in production.
\noindent\textbf{Answer.} Start with the mechanism: Servers emit tokens incrementally with backpressure, disconnect handling, and final usage accounting. Make the tradeoff explicit: Perceived latency improves, but partial outputs complicate moderation and retries. Define workload-specific offline metrics and a latency/cost budget, test important slices, instrument the critical path, and create a rollback criterion. The production review must explicitly guard against this failure: Retrying after visible partial output can duplicate text or side effects.
\par\textit{A strong answer connects mechanism, measurable tradeoffs, failure isolation, observability, and rollback rather than listing product features.}
\subsection{Backend, APIs, and Microservices}
\paragraph{MCQ-1441 [medium]} Which statement best describes FastAPI?
\begin{enumerate}[label=\Alph*.]
\item Token buckets, leaky buckets, and concurrency limits protect shared resources by identity and cost.
\item FastAPI builds typed ASGI APIs from Python hints with validation, dependency injection, and OpenAPI generation.
\item Litestar is an ASGI framework with typing, dependency injection, OpenAPI, plugins, and data-transfer objects.
\item gRPC uses protobuf contracts and HTTP/2 for efficient unary and streaming RPCs with generated clients.
\end{enumerate}
\noindent\textbf{Answer.} FastAPI builds typed ASGI APIs from Python hints with validation, dependency injection, and OpenAPI generation.
\par\textit{FastAPI builds typed ASGI APIs from Python hints with validation, dependency injection, and OpenAPI generation. Tradeoff: Development is fast, while blocking work must be isolated from the event loop. Watch for: Declaring a blocking database or model call inside async code stalls all requests.}
\paragraph{MCQ-1442 [hard]} What is the central engineering tradeoff of FastAPI?
\begin{enumerate}[label=\Alph*.]
\item Stability and fairness improve but strict limits can reject valuable bursts.
\item Development is fast, while blocking work must be isolated from the event loop.
\item Rich framework features reduce boilerplate but introduce conventions and migration cost.
\item Typed internal services perform well, while browser and human debugging are less direct.
\end{enumerate}
\noindent\textbf{Answer.} Development is fast, while blocking work must be isolated from the event loop.
\par\textit{FastAPI builds typed ASGI APIs from Python hints with validation, dependency injection, and OpenAPI generation. Tradeoff: Development is fast, while blocking work must be isolated from the event loop. Watch for: Declaring a blocking database or model call inside async code stalls all requests.}
\paragraph{MCQ-1443 [hard]} Which failure mode should an interviewer expect you to identify for FastAPI?
\begin{enumerate}[label=\Alph*.]
\item Limiting requests instead of tokens lets a few long prompts monopolize capacity.
\item Choosing by benchmark alone ignores ecosystem and team familiarity.
\item Declaring a blocking database or model call inside async code stalls all requests.
\item Changing field numbers or incompatible meanings breaks wire compatibility.
\end{enumerate}
\noindent\textbf{Answer.} Declaring a blocking database or model call inside async code stalls all requests.
\par\textit{FastAPI builds typed ASGI APIs from Python hints with validation, dependency injection, and OpenAPI generation. Tradeoff: Development is fast, while blocking work must be isolated from the event loop. Watch for: Declaring a blocking database or model call inside async code stalls all requests.}
\paragraph{MCQ-1444 [medium]} A design review is specifically concerned with this risk: Declaring a blocking database or model call inside async code stalls all requests. Which topic should be reviewed first?
\begin{enumerate}[label=\Alph*.]
\item Litestar
\item gRPC
\item Rate limiting
\item FastAPI
\end{enumerate}
\noindent\textbf{Answer.} FastAPI
\par\textit{FastAPI builds typed ASGI APIs from Python hints with validation, dependency injection, and OpenAPI generation. Tradeoff: Development is fast, while blocking work must be isolated from the event loop. Watch for: Declaring a blocking database or model call inside async code stalls all requests.}
\paragraph{MCQ-1445 [hard]} Which design note most accurately balances the benefits and costs of FastAPI?
\begin{enumerate}[label=\Alph*.]
\item Development is fast, while blocking work must be isolated from the event loop.
\item Rich framework features reduce boilerplate but introduce conventions and migration cost.
\item Stability and fairness improve but strict limits can reject valuable bursts.
\item Typed internal services perform well, while browser and human debugging are less direct.
\end{enumerate}
\noindent\textbf{Answer.} Development is fast, while blocking work must be isolated from the event loop.
\par\textit{FastAPI builds typed ASGI APIs from Python hints with validation, dependency injection, and OpenAPI generation. Tradeoff: Development is fast, while blocking work must be isolated from the event loop. Watch for: Declaring a blocking database or model call inside async code stalls all requests.}
\paragraph{MCQ-1446 [medium]} Which explanation would be strongest in a system-design interview about FastAPI?
\begin{enumerate}[label=\Alph*.]
\item Token buckets, leaky buckets, and concurrency limits protect shared resources by identity and cost.
\item FastAPI builds typed ASGI APIs from Python hints with validation, dependency injection, and OpenAPI generation.
\item gRPC uses protobuf contracts and HTTP/2 for efficient unary and streaming RPCs with generated clients.
\item Litestar is an ASGI framework with typing, dependency injection, OpenAPI, plugins, and data-transfer objects.
\end{enumerate}
\noindent\textbf{Answer.} FastAPI builds typed ASGI APIs from Python hints with validation, dependency injection, and OpenAPI generation.
\par\textit{FastAPI builds typed ASGI APIs from Python hints with validation, dependency injection, and OpenAPI generation. Tradeoff: Development is fast, while blocking work must be isolated from the event loop. Watch for: Declaring a blocking database or model call inside async code stalls all requests.}
\paragraph{MCQ-1447 [hard]} Which production warning is most directly associated with FastAPI?
\begin{enumerate}[label=\Alph*.]
\item Choosing by benchmark alone ignores ecosystem and team familiarity.
\item Changing field numbers or incompatible meanings breaks wire compatibility.
\item Limiting requests instead of tokens lets a few long prompts monopolize capacity.
\item Declaring a blocking database or model call inside async code stalls all requests.
\end{enumerate}
\noindent\textbf{Answer.} Declaring a blocking database or model call inside async code stalls all requests.
\par\textit{FastAPI builds typed ASGI APIs from Python hints with validation, dependency injection, and OpenAPI generation. Tradeoff: Development is fast, while blocking work must be isolated from the event loop. Watch for: Declaring a blocking database or model call inside async code stalls all requests.}
\paragraph{FC-1448 [medium]} Define FastAPI in interview-ready language.
\noindent\textbf{Answer.} FastAPI builds typed ASGI APIs from Python hints with validation, dependency injection, and OpenAPI generation.
\par\textit{FastAPI builds typed ASGI APIs from Python hints with validation, dependency injection, and OpenAPI generation.}
\paragraph{FC-1449 [hard]} State the main tradeoff and production pitfall for FastAPI.
\noindent\textbf{Answer.} Tradeoff: Development is fast, while blocking work must be isolated from the event loop. Pitfall: Declaring a blocking database or model call inside async code stalls all requests.
\par\textit{Tradeoff: Development is fast, while blocking work must be isolated from the event loop. Pitfall: Declaring a blocking database or model call inside async code stalls all requests.}
\paragraph{LA-1450 [hard]} Explain FastAPI from first principles, then show how you would evaluate and operate it in production.
\noindent\textbf{Answer.} Start with the mechanism: FastAPI builds typed ASGI APIs from Python hints with validation, dependency injection, and OpenAPI generation. Make the tradeoff explicit: Development is fast, while blocking work must be isolated from the event loop. Define workload-specific offline metrics and a latency/cost budget, test important slices, instrument the critical path, and create a rollback criterion. The production review must explicitly guard against this failure: Declaring a blocking database or model call inside async code stalls all requests.
\par\textit{A strong answer connects mechanism, measurable tradeoffs, failure isolation, observability, and rollback rather than listing product features.}
\paragraph{MCQ-1451 [medium]} Which statement best describes Starlette?
\begin{enumerate}[label=\Alph*.]
\item Token buckets, leaky buckets, and concurrency limits protect shared resources by identity and cost.
\item Idempotency keys let retried mutation requests resolve to the same business operation and stored outcome.
\item Starlette is a lightweight ASGI toolkit providing routing, middleware, WebSockets, background tasks, and test support.
\item Circuit breakers stop calls to a failing dependency, probe recovery, and combine with timeouts and fallbacks.
\end{enumerate}
\noindent\textbf{Answer.} Starlette is a lightweight ASGI toolkit providing routing, middleware, WebSockets, background tasks, and test support.
\par\textit{Starlette is a lightweight ASGI toolkit providing routing, middleware, WebSockets, background tasks, and test support. Tradeoff: It offers control with fewer batteries than a full API framework. Watch for: Reimplementing validation and dependency patterns inconsistently increases defects.}
\paragraph{MCQ-1452 [hard]} What is the central engineering tradeoff of Starlette?
\begin{enumerate}[label=\Alph*.]
\item Stability and fairness improve but strict limits can reject valuable bursts.
\item Duplicate effects are prevented at the cost of durable deduplication state.
\item Cascading failure is reduced, while bad thresholds can create synchronized outages.
\item It offers control with fewer batteries than a full API framework.
\end{enumerate}
\noindent\textbf{Answer.} It offers control with fewer batteries than a full API framework.
\par\textit{Starlette is a lightweight ASGI toolkit providing routing, middleware, WebSockets, background tasks, and test support. Tradeoff: It offers control with fewer batteries than a full API framework. Watch for: Reimplementing validation and dependency patterns inconsistently increases defects.}
\paragraph{MCQ-1453 [hard]} Which failure mode should an interviewer expect you to identify for Starlette?
\begin{enumerate}[label=\Alph*.]
\item Reimplementing validation and dependency patterns inconsistently increases defects.
\item A breaker without bounded timeouts reacts only after resources are already exhausted.
\item Expiring keys before upstream retry windows can reintroduce duplicates.
\item Limiting requests instead of tokens lets a few long prompts monopolize capacity.
\end{enumerate}
\noindent\textbf{Answer.} Reimplementing validation and dependency patterns inconsistently increases defects.
\par\textit{Starlette is a lightweight ASGI toolkit providing routing, middleware, WebSockets, background tasks, and test support. Tradeoff: It offers control with fewer batteries than a full API framework. Watch for: Reimplementing validation and dependency patterns inconsistently increases defects.}
\paragraph{MCQ-1454 [medium]} A design review is specifically concerned with this risk: Reimplementing validation and dependency patterns inconsistently increases defects. Which topic should be reviewed first?
\begin{enumerate}[label=\Alph*.]
\item API idempotency
\item Rate limiting
\item Starlette
\item Circuit breakers
\end{enumerate}
\noindent\textbf{Answer.} Starlette
\par\textit{Starlette is a lightweight ASGI toolkit providing routing, middleware, WebSockets, background tasks, and test support. Tradeoff: It offers control with fewer batteries than a full API framework. Watch for: Reimplementing validation and dependency patterns inconsistently increases defects.}
\paragraph{MCQ-1455 [hard]} Which design note most accurately balances the benefits and costs of Starlette?
\begin{enumerate}[label=\Alph*.]
\item Stability and fairness improve but strict limits can reject valuable bursts.
\item Cascading failure is reduced, while bad thresholds can create synchronized outages.
\item It offers control with fewer batteries than a full API framework.
\item Duplicate effects are prevented at the cost of durable deduplication state.
\end{enumerate}
\noindent\textbf{Answer.} It offers control with fewer batteries than a full API framework.
\par\textit{Starlette is a lightweight ASGI toolkit providing routing, middleware, WebSockets, background tasks, and test support. Tradeoff: It offers control with fewer batteries than a full API framework. Watch for: Reimplementing validation and dependency patterns inconsistently increases defects.}
\paragraph{MCQ-1456 [medium]} Which explanation would be strongest in a system-design interview about Starlette?
\begin{enumerate}[label=\Alph*.]
\item Token buckets, leaky buckets, and concurrency limits protect shared resources by identity and cost.
\item Circuit breakers stop calls to a failing dependency, probe recovery, and combine with timeouts and fallbacks.
\item Idempotency keys let retried mutation requests resolve to the same business operation and stored outcome.
\item Starlette is a lightweight ASGI toolkit providing routing, middleware, WebSockets, background tasks, and test support.
\end{enumerate}
\noindent\textbf{Answer.} Starlette is a lightweight ASGI toolkit providing routing, middleware, WebSockets, background tasks, and test support.
\par\textit{Starlette is a lightweight ASGI toolkit providing routing, middleware, WebSockets, background tasks, and test support. Tradeoff: It offers control with fewer batteries than a full API framework. Watch for: Reimplementing validation and dependency patterns inconsistently increases defects.}
\paragraph{MCQ-1457 [hard]} Which production warning is most directly associated with Starlette?
\begin{enumerate}[label=\Alph*.]
\item Reimplementing validation and dependency patterns inconsistently increases defects.
\item A breaker without bounded timeouts reacts only after resources are already exhausted.
\item Expiring keys before upstream retry windows can reintroduce duplicates.
\item Limiting requests instead of tokens lets a few long prompts monopolize capacity.
\end{enumerate}
\noindent\textbf{Answer.} Reimplementing validation and dependency patterns inconsistently increases defects.
\par\textit{Starlette is a lightweight ASGI toolkit providing routing, middleware, WebSockets, background tasks, and test support. Tradeoff: It offers control with fewer batteries than a full API framework. Watch for: Reimplementing validation and dependency patterns inconsistently increases defects.}
\paragraph{FC-1458 [medium]} Define Starlette in interview-ready language.
\noindent\textbf{Answer.} Starlette is a lightweight ASGI toolkit providing routing, middleware, WebSockets, background tasks, and test support.
\par\textit{Starlette is a lightweight ASGI toolkit providing routing, middleware, WebSockets, background tasks, and test support.}
\paragraph{FC-1459 [hard]} State the main tradeoff and production pitfall for Starlette.
\noindent\textbf{Answer.} Tradeoff: It offers control with fewer batteries than a full API framework. Pitfall: Reimplementing validation and dependency patterns inconsistently increases defects.
\par\textit{Tradeoff: It offers control with fewer batteries than a full API framework. Pitfall: Reimplementing validation and dependency patterns inconsistently increases defects.}
\paragraph{LA-1460 [hard]} Explain Starlette from first principles, then show how you would evaluate and operate it in production.
\noindent\textbf{Answer.} Start with the mechanism: Starlette is a lightweight ASGI toolkit providing routing, middleware, WebSockets, background tasks, and test support. Make the tradeoff explicit: It offers control with fewer batteries than a full API framework. Define workload-specific offline metrics and a latency/cost budget, test important slices, instrument the critical path, and create a rollback criterion. The production review must explicitly guard against this failure: Reimplementing validation and dependency patterns inconsistently increases defects.
\par\textit{A strong answer connects mechanism, measurable tradeoffs, failure isolation, observability, and rollback rather than listing product features.}
\paragraph{MCQ-1461 [medium]} Which statement best describes Litestar?
\begin{enumerate}[label=\Alph*.]
\item REST uses resource-oriented URLs, HTTP methods, status codes, caching, content negotiation, and idempotency semantics.
\item Idempotency keys let retried mutation requests resolve to the same business operation and stored outcome.
\item Boundaries should follow ownership, data, scaling, failure isolation, and change cadence rather than nouns alone.
\item Litestar is an ASGI framework with typing, dependency injection, OpenAPI, plugins, and data-transfer objects.
\end{enumerate}
\noindent\textbf{Answer.} Litestar is an ASGI framework with typing, dependency injection, OpenAPI, plugins, and data-transfer objects.
\par\textit{Litestar is an ASGI framework with typing, dependency injection, OpenAPI, plugins, and data-transfer objects. Tradeoff: Rich framework features reduce boilerplate but introduce conventions and migration cost. Watch for: Choosing by benchmark alone ignores ecosystem and team familiarity.}
\paragraph{MCQ-1462 [hard]} What is the central engineering tradeoff of Litestar?
\begin{enumerate}[label=\Alph*.]
\item Rich framework features reduce boilerplate but introduce conventions and migration cost.
\item Independent deployment helps teams but adds networks, consistency, and operations.
\item Broad interoperability is strong, while chatty workflows and schema evolution need care.
\item Duplicate effects are prevented at the cost of durable deduplication state.
\end{enumerate}
\noindent\textbf{Answer.} Rich framework features reduce boilerplate but introduce conventions and migration cost.
\par\textit{Litestar is an ASGI framework with typing, dependency injection, OpenAPI, plugins, and data-transfer objects. Tradeoff: Rich framework features reduce boilerplate but introduce conventions and migration cost. Watch for: Choosing by benchmark alone ignores ecosystem and team familiarity.}
\paragraph{MCQ-1463 [hard]} Which failure mode should an interviewer expect you to identify for Litestar?
\begin{enumerate}[label=\Alph*.]
\item Returning 200 for every failure breaks clients and observability.
\item Expiring keys before upstream retry windows can reintroduce duplicates.
\item Choosing by benchmark alone ignores ecosystem and team familiarity.
\item Splitting too early creates a distributed monolith.
\end{enumerate}
\noindent\textbf{Answer.} Choosing by benchmark alone ignores ecosystem and team familiarity.
\par\textit{Litestar is an ASGI framework with typing, dependency injection, OpenAPI, plugins, and data-transfer objects. Tradeoff: Rich framework features reduce boilerplate but introduce conventions and migration cost. Watch for: Choosing by benchmark alone ignores ecosystem and team familiarity.}
\paragraph{MCQ-1464 [medium]} A design review is specifically concerned with this risk: Choosing by benchmark alone ignores ecosystem and team familiarity. Which topic should be reviewed first?
\begin{enumerate}[label=\Alph*.]
\item REST API design
\item Microservice boundaries
\item API idempotency
\item Litestar
\end{enumerate}
\noindent\textbf{Answer.} Litestar
\par\textit{Litestar is an ASGI framework with typing, dependency injection, OpenAPI, plugins, and data-transfer objects. Tradeoff: Rich framework features reduce boilerplate but introduce conventions and migration cost. Watch for: Choosing by benchmark alone ignores ecosystem and team familiarity.}
\paragraph{MCQ-1465 [hard]} Which design note most accurately balances the benefits and costs of Litestar?
\begin{enumerate}[label=\Alph*.]
\item Independent deployment helps teams but adds networks, consistency, and operations.
\item Rich framework features reduce boilerplate but introduce conventions and migration cost.
\item Broad interoperability is strong, while chatty workflows and schema evolution need care.
\item Duplicate effects are prevented at the cost of durable deduplication state.
\end{enumerate}
\noindent\textbf{Answer.} Rich framework features reduce boilerplate but introduce conventions and migration cost.
\par\textit{Litestar is an ASGI framework with typing, dependency injection, OpenAPI, plugins, and data-transfer objects. Tradeoff: Rich framework features reduce boilerplate but introduce conventions and migration cost. Watch for: Choosing by benchmark alone ignores ecosystem and team familiarity.}
\paragraph{MCQ-1466 [medium]} Which explanation would be strongest in a system-design interview about Litestar?
\begin{enumerate}[label=\Alph*.]
\item REST uses resource-oriented URLs, HTTP methods, status codes, caching, content negotiation, and idempotency semantics.
\item Litestar is an ASGI framework with typing, dependency injection, OpenAPI, plugins, and data-transfer objects.
\item Idempotency keys let retried mutation requests resolve to the same business operation and stored outcome.
\item Boundaries should follow ownership, data, scaling, failure isolation, and change cadence rather than nouns alone.
\end{enumerate}
\noindent\textbf{Answer.} Litestar is an ASGI framework with typing, dependency injection, OpenAPI, plugins, and data-transfer objects.
\par\textit{Litestar is an ASGI framework with typing, dependency injection, OpenAPI, plugins, and data-transfer objects. Tradeoff: Rich framework features reduce boilerplate but introduce conventions and migration cost. Watch for: Choosing by benchmark alone ignores ecosystem and team familiarity.}
\paragraph{MCQ-1467 [hard]} Which production warning is most directly associated with Litestar?
\begin{enumerate}[label=\Alph*.]
\item Expiring keys before upstream retry windows can reintroduce duplicates.
\item Choosing by benchmark alone ignores ecosystem and team familiarity.
\item Splitting too early creates a distributed monolith.
\item Returning 200 for every failure breaks clients and observability.
\end{enumerate}
\noindent\textbf{Answer.} Choosing by benchmark alone ignores ecosystem and team familiarity.
\par\textit{Litestar is an ASGI framework with typing, dependency injection, OpenAPI, plugins, and data-transfer objects. Tradeoff: Rich framework features reduce boilerplate but introduce conventions and migration cost. Watch for: Choosing by benchmark alone ignores ecosystem and team familiarity.}
\paragraph{FC-1468 [medium]} Define Litestar in interview-ready language.
\noindent\textbf{Answer.} Litestar is an ASGI framework with typing, dependency injection, OpenAPI, plugins, and data-transfer objects.
\par\textit{Litestar is an ASGI framework with typing, dependency injection, OpenAPI, plugins, and data-transfer objects.}
\paragraph{FC-1469 [hard]} State the main tradeoff and production pitfall for Litestar.
\noindent\textbf{Answer.} Tradeoff: Rich framework features reduce boilerplate but introduce conventions and migration cost. Pitfall: Choosing by benchmark alone ignores ecosystem and team familiarity.
\par\textit{Tradeoff: Rich framework features reduce boilerplate but introduce conventions and migration cost. Pitfall: Choosing by benchmark alone ignores ecosystem and team familiarity.}
\paragraph{LA-1470 [hard]} Explain Litestar from first principles, then show how you would evaluate and operate it in production.
\noindent\textbf{Answer.} Start with the mechanism: Litestar is an ASGI framework with typing, dependency injection, OpenAPI, plugins, and data-transfer objects. Make the tradeoff explicit: Rich framework features reduce boilerplate but introduce conventions and migration cost. Define workload-specific offline metrics and a latency/cost budget, test important slices, instrument the critical path, and create a rollback criterion. The production review must explicitly guard against this failure: Choosing by benchmark alone ignores ecosystem and team familiarity.
\par\textit{A strong answer connects mechanism, measurable tradeoffs, failure isolation, observability, and rollback rather than listing product features.}
\paragraph{MCQ-1471 [medium]} Which statement best describes Flask?
\begin{enumerate}[label=\Alph*.]
\item Starlette is a lightweight ASGI toolkit providing routing, middleware, WebSockets, background tasks, and test support.
\item Flask is a minimal WSGI web framework with a large extension ecosystem.
\item Boundaries should follow ownership, data, scaling, failure isolation, and change cadence rather than nouns alone.
\item gRPC uses protobuf contracts and HTTP/2 for efficient unary and streaming RPCs with generated clients.
\end{enumerate}
\noindent\textbf{Answer.} Flask is a minimal WSGI web framework with a large extension ecosystem.
\par\textit{Flask is a minimal WSGI web framework with a large extension ecosystem. Tradeoff: Simplicity suits synchronous services, while async-first streaming needs different architecture. Watch for: Assuming an async view converts a WSGI deployment into high-concurrency ASGI is wrong.}
\paragraph{MCQ-1472 [hard]} What is the central engineering tradeoff of Flask?
\begin{enumerate}[label=\Alph*.]
\item Typed internal services perform well, while browser and human debugging are less direct.
\item Independent deployment helps teams but adds networks, consistency, and operations.
\item It offers control with fewer batteries than a full API framework.
\item Simplicity suits synchronous services, while async-first streaming needs different architecture.
\end{enumerate}
\noindent\textbf{Answer.} Simplicity suits synchronous services, while async-first streaming needs different architecture.
\par\textit{Flask is a minimal WSGI web framework with a large extension ecosystem. Tradeoff: Simplicity suits synchronous services, while async-first streaming needs different architecture. Watch for: Assuming an async view converts a WSGI deployment into high-concurrency ASGI is wrong.}
\paragraph{MCQ-1473 [hard]} Which failure mode should an interviewer expect you to identify for Flask?
\begin{enumerate}[label=\Alph*.]
\item Reimplementing validation and dependency patterns inconsistently increases defects.
\item Assuming an async view converts a WSGI deployment into high-concurrency ASGI is wrong.
\item Splitting too early creates a distributed monolith.
\item Changing field numbers or incompatible meanings breaks wire compatibility.
\end{enumerate}
\noindent\textbf{Answer.} Assuming an async view converts a WSGI deployment into high-concurrency ASGI is wrong.
\par\textit{Flask is a minimal WSGI web framework with a large extension ecosystem. Tradeoff: Simplicity suits synchronous services, while async-first streaming needs different architecture. Watch for: Assuming an async view converts a WSGI deployment into high-concurrency ASGI is wrong.}
\paragraph{MCQ-1474 [medium]} A design review is specifically concerned with this risk: Assuming an async view converts a WSGI deployment into high-concurrency ASGI is wrong. Which topic should be reviewed first?
\begin{enumerate}[label=\Alph*.]
\item Microservice boundaries
\item Flask
\item Starlette
\item gRPC
\end{enumerate}
\noindent\textbf{Answer.} Flask
\par\textit{Flask is a minimal WSGI web framework with a large extension ecosystem. Tradeoff: Simplicity suits synchronous services, while async-first streaming needs different architecture. Watch for: Assuming an async view converts a WSGI deployment into high-concurrency ASGI is wrong.}
\paragraph{MCQ-1475 [hard]} Which design note most accurately balances the benefits and costs of Flask?
\begin{enumerate}[label=\Alph*.]
\item Typed internal services perform well, while browser and human debugging are less direct.
\item Simplicity suits synchronous services, while async-first streaming needs different architecture.
\item Independent deployment helps teams but adds networks, consistency, and operations.
\item It offers control with fewer batteries than a full API framework.
\end{enumerate}
\noindent\textbf{Answer.} Simplicity suits synchronous services, while async-first streaming needs different architecture.
\par\textit{Flask is a minimal WSGI web framework with a large extension ecosystem. Tradeoff: Simplicity suits synchronous services, while async-first streaming needs different architecture. Watch for: Assuming an async view converts a WSGI deployment into high-concurrency ASGI is wrong.}
\paragraph{MCQ-1476 [medium]} Which explanation would be strongest in a system-design interview about Flask?
\begin{enumerate}[label=\Alph*.]
\item Boundaries should follow ownership, data, scaling, failure isolation, and change cadence rather than nouns alone.
\item Starlette is a lightweight ASGI toolkit providing routing, middleware, WebSockets, background tasks, and test support.
\item Flask is a minimal WSGI web framework with a large extension ecosystem.
\item gRPC uses protobuf contracts and HTTP/2 for efficient unary and streaming RPCs with generated clients.
\end{enumerate}
\noindent\textbf{Answer.} Flask is a minimal WSGI web framework with a large extension ecosystem.
\par\textit{Flask is a minimal WSGI web framework with a large extension ecosystem. Tradeoff: Simplicity suits synchronous services, while async-first streaming needs different architecture. Watch for: Assuming an async view converts a WSGI deployment into high-concurrency ASGI is wrong.}
\paragraph{MCQ-1477 [hard]} Which production warning is most directly associated with Flask?
\begin{enumerate}[label=\Alph*.]
\item Reimplementing validation and dependency patterns inconsistently increases defects.
\item Splitting too early creates a distributed monolith.
\item Assuming an async view converts a WSGI deployment into high-concurrency ASGI is wrong.
\item Changing field numbers or incompatible meanings breaks wire compatibility.
\end{enumerate}
\noindent\textbf{Answer.} Assuming an async view converts a WSGI deployment into high-concurrency ASGI is wrong.
\par\textit{Flask is a minimal WSGI web framework with a large extension ecosystem. Tradeoff: Simplicity suits synchronous services, while async-first streaming needs different architecture. Watch for: Assuming an async view converts a WSGI deployment into high-concurrency ASGI is wrong.}
\paragraph{FC-1478 [medium]} Define Flask in interview-ready language.
\noindent\textbf{Answer.} Flask is a minimal WSGI web framework with a large extension ecosystem.
\par\textit{Flask is a minimal WSGI web framework with a large extension ecosystem.}
\paragraph{FC-1479 [hard]} State the main tradeoff and production pitfall for Flask.
\noindent\textbf{Answer.} Tradeoff: Simplicity suits synchronous services, while async-first streaming needs different architecture. Pitfall: Assuming an async view converts a WSGI deployment into high-concurrency ASGI is wrong.
\par\textit{Tradeoff: Simplicity suits synchronous services, while async-first streaming needs different architecture. Pitfall: Assuming an async view converts a WSGI deployment into high-concurrency ASGI is wrong.}
\paragraph{LA-1480 [hard]} Explain Flask from first principles, then show how you would evaluate and operate it in production.
\noindent\textbf{Answer.} Start with the mechanism: Flask is a minimal WSGI web framework with a large extension ecosystem. Make the tradeoff explicit: Simplicity suits synchronous services, while async-first streaming needs different architecture. Define workload-specific offline metrics and a latency/cost budget, test important slices, instrument the critical path, and create a rollback criterion. The production review must explicitly guard against this failure: Assuming an async view converts a WSGI deployment into high-concurrency ASGI is wrong.
\par\textit{A strong answer connects mechanism, measurable tradeoffs, failure isolation, observability, and rollback rather than listing product features.}
\paragraph{MCQ-1481 [medium]} Which statement best describes Django and DRF?
\begin{enumerate}[label=\Alph*.]
\item Flask is a minimal WSGI web framework with a large extension ecosystem.
\item Litestar is an ASGI framework with typing, dependency injection, OpenAPI, plugins, and data-transfer objects.
\item Encapsulation, composition, interfaces, and explicit invariants organize stateful behavior.
\item Django supplies an ORM, admin, authentication, migrations, and conventions; DRF adds APIs.
\end{enumerate}
\noindent\textbf{Answer.} Django supplies an ORM, admin, authentication, migrations, and conventions; DRF adds APIs.
\par\textit{Django supplies an ORM, admin, authentication, migrations, and conventions; DRF adds APIs. Tradeoff: Batteries accelerate business systems but can be heavy for a narrow inference gateway. Watch for: Performing N+1 ORM queries around an expensive model magnifies tail latency.}
\paragraph{MCQ-1482 [hard]} What is the central engineering tradeoff of Django and DRF?
\begin{enumerate}[label=\Alph*.]
\item Abstraction improves changeability, while deep inheritance hides coupling.
\item Simplicity suits synchronous services, while async-first streaming needs different architecture.
\item Batteries accelerate business systems but can be heavy for a narrow inference gateway.
\item Rich framework features reduce boilerplate but introduce conventions and migration cost.
\end{enumerate}
\noindent\textbf{Answer.} Batteries accelerate business systems but can be heavy for a narrow inference gateway.
\par\textit{Django supplies an ORM, admin, authentication, migrations, and conventions; DRF adds APIs. Tradeoff: Batteries accelerate business systems but can be heavy for a narrow inference gateway. Watch for: Performing N+1 ORM queries around an expensive model magnifies tail latency.}
\paragraph{MCQ-1483 [hard]} Which failure mode should an interviewer expect you to identify for Django and DRF?
\begin{enumerate}[label=\Alph*.]
\item Choosing by benchmark alone ignores ecosystem and team familiarity.
\item Performing N+1 ORM queries around an expensive model magnifies tail latency.
\item Modeling every data record as a behavior-heavy class adds ceremony without invariants.
\item Assuming an async view converts a WSGI deployment into high-concurrency ASGI is wrong.
\end{enumerate}
\noindent\textbf{Answer.} Performing N+1 ORM queries around an expensive model magnifies tail latency.
\par\textit{Django supplies an ORM, admin, authentication, migrations, and conventions; DRF adds APIs. Tradeoff: Batteries accelerate business systems but can be heavy for a narrow inference gateway. Watch for: Performing N+1 ORM queries around an expensive model magnifies tail latency.}
\paragraph{MCQ-1484 [medium]} A design review is specifically concerned with this risk: Performing N+1 ORM queries around an expensive model magnifies tail latency. Which topic should be reviewed first?
\begin{enumerate}[label=\Alph*.]
\item Litestar
\item Flask
\item Object-oriented design
\item Django and DRF
\end{enumerate}
\noindent\textbf{Answer.} Django and DRF
\par\textit{Django supplies an ORM, admin, authentication, migrations, and conventions; DRF adds APIs. Tradeoff: Batteries accelerate business systems but can be heavy for a narrow inference gateway. Watch for: Performing N+1 ORM queries around an expensive model magnifies tail latency.}
\paragraph{MCQ-1485 [hard]} Which design note most accurately balances the benefits and costs of Django and DRF?
\begin{enumerate}[label=\Alph*.]
\item Batteries accelerate business systems but can be heavy for a narrow inference gateway.
\item Abstraction improves changeability, while deep inheritance hides coupling.
\item Rich framework features reduce boilerplate but introduce conventions and migration cost.
\item Simplicity suits synchronous services, while async-first streaming needs different architecture.
\end{enumerate}
\noindent\textbf{Answer.} Batteries accelerate business systems but can be heavy for a narrow inference gateway.
\par\textit{Django supplies an ORM, admin, authentication, migrations, and conventions; DRF adds APIs. Tradeoff: Batteries accelerate business systems but can be heavy for a narrow inference gateway. Watch for: Performing N+1 ORM queries around an expensive model magnifies tail latency.}
\paragraph{MCQ-1486 [medium]} Which explanation would be strongest in a system-design interview about Django and DRF?
\begin{enumerate}[label=\Alph*.]
\item Encapsulation, composition, interfaces, and explicit invariants organize stateful behavior.
\item Litestar is an ASGI framework with typing, dependency injection, OpenAPI, plugins, and data-transfer objects.
\item Flask is a minimal WSGI web framework with a large extension ecosystem.
\item Django supplies an ORM, admin, authentication, migrations, and conventions; DRF adds APIs.
\end{enumerate}
\noindent\textbf{Answer.} Django supplies an ORM, admin, authentication, migrations, and conventions; DRF adds APIs.
\par\textit{Django supplies an ORM, admin, authentication, migrations, and conventions; DRF adds APIs. Tradeoff: Batteries accelerate business systems but can be heavy for a narrow inference gateway. Watch for: Performing N+1 ORM queries around an expensive model magnifies tail latency.}
\paragraph{MCQ-1487 [hard]} Which production warning is most directly associated with Django and DRF?
\begin{enumerate}[label=\Alph*.]
\item Assuming an async view converts a WSGI deployment into high-concurrency ASGI is wrong.
\item Modeling every data record as a behavior-heavy class adds ceremony without invariants.
\item Choosing by benchmark alone ignores ecosystem and team familiarity.
\item Performing N+1 ORM queries around an expensive model magnifies tail latency.
\end{enumerate}
\noindent\textbf{Answer.} Performing N+1 ORM queries around an expensive model magnifies tail latency.
\par\textit{Django supplies an ORM, admin, authentication, migrations, and conventions; DRF adds APIs. Tradeoff: Batteries accelerate business systems but can be heavy for a narrow inference gateway. Watch for: Performing N+1 ORM queries around an expensive model magnifies tail latency.}
\paragraph{FC-1488 [medium]} Define Django and DRF in interview-ready language.
\noindent\textbf{Answer.} Django supplies an ORM, admin, authentication, migrations, and conventions; DRF adds APIs.
\par\textit{Django supplies an ORM, admin, authentication, migrations, and conventions; DRF adds APIs.}
\paragraph{FC-1489 [hard]} State the main tradeoff and production pitfall for Django and DRF.
\noindent\textbf{Answer.} Tradeoff: Batteries accelerate business systems but can be heavy for a narrow inference gateway. Pitfall: Performing N+1 ORM queries around an expensive model magnifies tail latency.
\par\textit{Tradeoff: Batteries accelerate business systems but can be heavy for a narrow inference gateway. Pitfall: Performing N+1 ORM queries around an expensive model magnifies tail latency.}
\paragraph{LA-1490 [hard]} Explain Django and DRF from first principles, then show how you would evaluate and operate it in production.
\noindent\textbf{Answer.} Start with the mechanism: Django supplies an ORM, admin, authentication, migrations, and conventions; DRF adds APIs. Make the tradeoff explicit: Batteries accelerate business systems but can be heavy for a narrow inference gateway. Define workload-specific offline metrics and a latency/cost budget, test important slices, instrument the critical path, and create a rollback criterion. The production review must explicitly guard against this failure: Performing N+1 ORM queries around an expensive model magnifies tail latency.
\par\textit{A strong answer connects mechanism, measurable tradeoffs, failure isolation, observability, and rollback rather than listing product features.}
\paragraph{MCQ-1491 [medium]} Which statement best describes REST API design?
\begin{enumerate}[label=\Alph*.]
\item Litestar is an ASGI framework with typing, dependency injection, OpenAPI, plugins, and data-transfer objects.
\item Circuit breakers stop calls to a failing dependency, probe recovery, and combine with timeouts and fallbacks.
\item Producers publish immutable events to brokers while consumers process independently with delivery semantics.
\item REST uses resource-oriented URLs, HTTP methods, status codes, caching, content negotiation, and idempotency semantics.
\end{enumerate}
\noindent\textbf{Answer.} REST uses resource-oriented URLs, HTTP methods, status codes, caching, content negotiation, and idempotency semantics.
\par\textit{REST uses resource-oriented URLs, HTTP methods, status codes, caching, content negotiation, and idempotency semantics. Tradeoff: Broad interoperability is strong, while chatty workflows and schema evolution need care. Watch for: Returning 200 for every failure breaks clients and observability.}
\paragraph{MCQ-1492 [hard]} What is the central engineering tradeoff of REST API design?
\begin{enumerate}[label=\Alph*.]
\item Broad interoperability is strong, while chatty workflows and schema evolution need care.
\item Rich framework features reduce boilerplate but introduce conventions and migration cost.
\item Decoupling and replay improve, while ordering, duplication, schemas, and lag must be managed.
\item Cascading failure is reduced, while bad thresholds can create synchronized outages.
\end{enumerate}
\noindent\textbf{Answer.} Broad interoperability is strong, while chatty workflows and schema evolution need care.
\par\textit{REST uses resource-oriented URLs, HTTP methods, status codes, caching, content negotiation, and idempotency semantics. Tradeoff: Broad interoperability is strong, while chatty workflows and schema evolution need care. Watch for: Returning 200 for every failure breaks clients and observability.}
\paragraph{MCQ-1493 [hard]} Which failure mode should an interviewer expect you to identify for REST API design?
\begin{enumerate}[label=\Alph*.]
\item Returning 200 for every failure breaks clients and observability.
\item A breaker without bounded timeouts reacts only after resources are already exhausted.
\item Choosing by benchmark alone ignores ecosystem and team familiarity.
\item Assuming exactly-once business effects from at-least-once delivery causes duplicates.
\end{enumerate}
\noindent\textbf{Answer.} Returning 200 for every failure breaks clients and observability.
\par\textit{REST uses resource-oriented URLs, HTTP methods, status codes, caching, content negotiation, and idempotency semantics. Tradeoff: Broad interoperability is strong, while chatty workflows and schema evolution need care. Watch for: Returning 200 for every failure breaks clients and observability.}
\paragraph{MCQ-1494 [medium]} A design review is specifically concerned with this risk: Returning 200 for every failure breaks clients and observability. Which topic should be reviewed first?
\begin{enumerate}[label=\Alph*.]
\item REST API design
\item Litestar
\item Circuit breakers
\item Event-driven architecture
\end{enumerate}
\noindent\textbf{Answer.} REST API design
\par\textit{REST uses resource-oriented URLs, HTTP methods, status codes, caching, content negotiation, and idempotency semantics. Tradeoff: Broad interoperability is strong, while chatty workflows and schema evolution need care. Watch for: Returning 200 for every failure breaks clients and observability.}
\paragraph{MCQ-1495 [hard]} Which design note most accurately balances the benefits and costs of REST API design?
\begin{enumerate}[label=\Alph*.]
\item Cascading failure is reduced, while bad thresholds can create synchronized outages.
\item Decoupling and replay improve, while ordering, duplication, schemas, and lag must be managed.
\item Rich framework features reduce boilerplate but introduce conventions and migration cost.
\item Broad interoperability is strong, while chatty workflows and schema evolution need care.
\end{enumerate}
\noindent\textbf{Answer.} Broad interoperability is strong, while chatty workflows and schema evolution need care.
\par\textit{REST uses resource-oriented URLs, HTTP methods, status codes, caching, content negotiation, and idempotency semantics. Tradeoff: Broad interoperability is strong, while chatty workflows and schema evolution need care. Watch for: Returning 200 for every failure breaks clients and observability.}
\paragraph{MCQ-1496 [medium]} Which explanation would be strongest in a system-design interview about REST API design?
\begin{enumerate}[label=\Alph*.]
\item Circuit breakers stop calls to a failing dependency, probe recovery, and combine with timeouts and fallbacks.
\item Producers publish immutable events to brokers while consumers process independently with delivery semantics.
\item REST uses resource-oriented URLs, HTTP methods, status codes, caching, content negotiation, and idempotency semantics.
\item Litestar is an ASGI framework with typing, dependency injection, OpenAPI, plugins, and data-transfer objects.
\end{enumerate}
\noindent\textbf{Answer.} REST uses resource-oriented URLs, HTTP methods, status codes, caching, content negotiation, and idempotency semantics.
\par\textit{REST uses resource-oriented URLs, HTTP methods, status codes, caching, content negotiation, and idempotency semantics. Tradeoff: Broad interoperability is strong, while chatty workflows and schema evolution need care. Watch for: Returning 200 for every failure breaks clients and observability.}
\paragraph{MCQ-1497 [hard]} Which production warning is most directly associated with REST API design?
\begin{enumerate}[label=\Alph*.]
\item Choosing by benchmark alone ignores ecosystem and team familiarity.
\item A breaker without bounded timeouts reacts only after resources are already exhausted.
\item Assuming exactly-once business effects from at-least-once delivery causes duplicates.
\item Returning 200 for every failure breaks clients and observability.
\end{enumerate}
\noindent\textbf{Answer.} Returning 200 for every failure breaks clients and observability.
\par\textit{REST uses resource-oriented URLs, HTTP methods, status codes, caching, content negotiation, and idempotency semantics. Tradeoff: Broad interoperability is strong, while chatty workflows and schema evolution need care. Watch for: Returning 200 for every failure breaks clients and observability.}
\paragraph{FC-1498 [medium]} Define REST API design in interview-ready language.
\noindent\textbf{Answer.} REST uses resource-oriented URLs, HTTP methods, status codes, caching, content negotiation, and idempotency semantics.
\par\textit{REST uses resource-oriented URLs, HTTP methods, status codes, caching, content negotiation, and idempotency semantics.}
\paragraph{FC-1499 [hard]} State the main tradeoff and production pitfall for REST API design.
\noindent\textbf{Answer.} Tradeoff: Broad interoperability is strong, while chatty workflows and schema evolution need care. Pitfall: Returning 200 for every failure breaks clients and observability.
\par\textit{Tradeoff: Broad interoperability is strong, while chatty workflows and schema evolution need care. Pitfall: Returning 200 for every failure breaks clients and observability.}
\paragraph{LA-1500 [hard]} Explain REST API design from first principles, then show how you would evaluate and operate it in production.
\noindent\textbf{Answer.} Start with the mechanism: REST uses resource-oriented URLs, HTTP methods, status codes, caching, content negotiation, and idempotency semantics. Make the tradeoff explicit: Broad interoperability is strong, while chatty workflows and schema evolution need care. Define workload-specific offline metrics and a latency/cost budget, test important slices, instrument the critical path, and create a rollback criterion. The production review must explicitly guard against this failure: Returning 200 for every failure breaks clients and observability.
\par\textit{A strong answer connects mechanism, measurable tradeoffs, failure isolation, observability, and rollback rather than listing product features.}
\paragraph{MCQ-1501 [medium]} Which statement best describes gRPC?
\begin{enumerate}[label=\Alph*.]
\item gRPC uses protobuf contracts and HTTP/2 for efficient unary and streaming RPCs with generated clients.
\item Flask is a minimal WSGI web framework with a large extension ecosystem.
\item Producers publish immutable events to brokers while consumers process independently with delivery semantics.
\item FastAPI builds typed ASGI APIs from Python hints with validation, dependency injection, and OpenAPI generation.
\end{enumerate}
\noindent\textbf{Answer.} gRPC uses protobuf contracts and HTTP/2 for efficient unary and streaming RPCs with generated clients.
\par\textit{gRPC uses protobuf contracts and HTTP/2 for efficient unary and streaming RPCs with generated clients. Tradeoff: Typed internal services perform well, while browser and human debugging are less direct. Watch for: Changing field numbers or incompatible meanings breaks wire compatibility.}
\paragraph{MCQ-1502 [hard]} What is the central engineering tradeoff of gRPC?
\begin{enumerate}[label=\Alph*.]
\item Development is fast, while blocking work must be isolated from the event loop.
\item Simplicity suits synchronous services, while async-first streaming needs different architecture.
\item Decoupling and replay improve, while ordering, duplication, schemas, and lag must be managed.
\item Typed internal services perform well, while browser and human debugging are less direct.
\end{enumerate}
\noindent\textbf{Answer.} Typed internal services perform well, while browser and human debugging are less direct.
\par\textit{gRPC uses protobuf contracts and HTTP/2 for efficient unary and streaming RPCs with generated clients. Tradeoff: Typed internal services perform well, while browser and human debugging are less direct. Watch for: Changing field numbers or incompatible meanings breaks wire compatibility.}
\paragraph{MCQ-1503 [hard]} Which failure mode should an interviewer expect you to identify for gRPC?
\begin{enumerate}[label=\Alph*.]
\item Changing field numbers or incompatible meanings breaks wire compatibility.
\item Declaring a blocking database or model call inside async code stalls all requests.
\item Assuming exactly-once business effects from at-least-once delivery causes duplicates.
\item Assuming an async view converts a WSGI deployment into high-concurrency ASGI is wrong.
\end{enumerate}
\noindent\textbf{Answer.} Changing field numbers or incompatible meanings breaks wire compatibility.
\par\textit{gRPC uses protobuf contracts and HTTP/2 for efficient unary and streaming RPCs with generated clients. Tradeoff: Typed internal services perform well, while browser and human debugging are less direct. Watch for: Changing field numbers or incompatible meanings breaks wire compatibility.}
\paragraph{MCQ-1504 [medium]} A design review is specifically concerned with this risk: Changing field numbers or incompatible meanings breaks wire compatibility. Which topic should be reviewed first?
\begin{enumerate}[label=\Alph*.]
\item Event-driven architecture
\item Flask
\item FastAPI
\item gRPC
\end{enumerate}
\noindent\textbf{Answer.} gRPC
\par\textit{gRPC uses protobuf contracts and HTTP/2 for efficient unary and streaming RPCs with generated clients. Tradeoff: Typed internal services perform well, while browser and human debugging are less direct. Watch for: Changing field numbers or incompatible meanings breaks wire compatibility.}
\paragraph{MCQ-1505 [hard]} Which design note most accurately balances the benefits and costs of gRPC?
\begin{enumerate}[label=\Alph*.]
\item Typed internal services perform well, while browser and human debugging are less direct.
\item Development is fast, while blocking work must be isolated from the event loop.
\item Decoupling and replay improve, while ordering, duplication, schemas, and lag must be managed.
\item Simplicity suits synchronous services, while async-first streaming needs different architecture.
\end{enumerate}
\noindent\textbf{Answer.} Typed internal services perform well, while browser and human debugging are less direct.
\par\textit{gRPC uses protobuf contracts and HTTP/2 for efficient unary and streaming RPCs with generated clients. Tradeoff: Typed internal services perform well, while browser and human debugging are less direct. Watch for: Changing field numbers or incompatible meanings breaks wire compatibility.}
\paragraph{MCQ-1506 [medium]} Which explanation would be strongest in a system-design interview about gRPC?
\begin{enumerate}[label=\Alph*.]
\item gRPC uses protobuf contracts and HTTP/2 for efficient unary and streaming RPCs with generated clients.
\item Producers publish immutable events to brokers while consumers process independently with delivery semantics.
\item FastAPI builds typed ASGI APIs from Python hints with validation, dependency injection, and OpenAPI generation.
\item Flask is a minimal WSGI web framework with a large extension ecosystem.
\end{enumerate}
\noindent\textbf{Answer.} gRPC uses protobuf contracts and HTTP/2 for efficient unary and streaming RPCs with generated clients.
\par\textit{gRPC uses protobuf contracts and HTTP/2 for efficient unary and streaming RPCs with generated clients. Tradeoff: Typed internal services perform well, while browser and human debugging are less direct. Watch for: Changing field numbers or incompatible meanings breaks wire compatibility.}
\paragraph{MCQ-1507 [hard]} Which production warning is most directly associated with gRPC?
\begin{enumerate}[label=\Alph*.]
\item Assuming an async view converts a WSGI deployment into high-concurrency ASGI is wrong.
\item Declaring a blocking database or model call inside async code stalls all requests.
\item Assuming exactly-once business effects from at-least-once delivery causes duplicates.
\item Changing field numbers or incompatible meanings breaks wire compatibility.
\end{enumerate}
\noindent\textbf{Answer.} Changing field numbers or incompatible meanings breaks wire compatibility.
\par\textit{gRPC uses protobuf contracts and HTTP/2 for efficient unary and streaming RPCs with generated clients. Tradeoff: Typed internal services perform well, while browser and human debugging are less direct. Watch for: Changing field numbers or incompatible meanings breaks wire compatibility.}
\paragraph{FC-1508 [medium]} Define gRPC in interview-ready language.
\noindent\textbf{Answer.} gRPC uses protobuf contracts and HTTP/2 for efficient unary and streaming RPCs with generated clients.
\par\textit{gRPC uses protobuf contracts and HTTP/2 for efficient unary and streaming RPCs with generated clients.}
\paragraph{FC-1509 [hard]} State the main tradeoff and production pitfall for gRPC.
\noindent\textbf{Answer.} Tradeoff: Typed internal services perform well, while browser and human debugging are less direct. Pitfall: Changing field numbers or incompatible meanings breaks wire compatibility.
\par\textit{Tradeoff: Typed internal services perform well, while browser and human debugging are less direct. Pitfall: Changing field numbers or incompatible meanings breaks wire compatibility.}
\paragraph{LA-1510 [hard]} Explain gRPC from first principles, then show how you would evaluate and operate it in production.
\noindent\textbf{Answer.} Start with the mechanism: gRPC uses protobuf contracts and HTTP/2 for efficient unary and streaming RPCs with generated clients. Make the tradeoff explicit: Typed internal services perform well, while browser and human debugging are less direct. Define workload-specific offline metrics and a latency/cost budget, test important slices, instrument the critical path, and create a rollback criterion. The production review must explicitly guard against this failure: Changing field numbers or incompatible meanings breaks wire compatibility.
\par\textit{A strong answer connects mechanism, measurable tradeoffs, failure isolation, observability, and rollback rather than listing product features.}
\paragraph{MCQ-1511 [medium]} Which statement best describes WebSockets and SSE?
\begin{enumerate}[label=\Alph*.]
\item Flask is a minimal WSGI web framework with a large extension ecosystem.
\item Idempotency keys let retried mutation requests resolve to the same business operation and stored outcome.
\item WebSockets provide bidirectional frames; server-sent events provide simpler one-way HTTP streams.
\item gRPC uses protobuf contracts and HTTP/2 for efficient unary and streaming RPCs with generated clients.
\end{enumerate}
\noindent\textbf{Answer.} WebSockets provide bidirectional frames; server-sent events provide simpler one-way HTTP streams.
\par\textit{WebSockets provide bidirectional frames; server-sent events provide simpler one-way HTTP streams. Tradeoff: SSE fits token streaming and reconnect semantics; WebSockets fit interactive duplex protocols. Watch for: Using WebSockets when one-way streaming is enough adds state and proxy complexity.}
\paragraph{MCQ-1512 [hard]} What is the central engineering tradeoff of WebSockets and SSE?
\begin{enumerate}[label=\Alph*.]
\item Duplicate effects are prevented at the cost of durable deduplication state.
\item Typed internal services perform well, while browser and human debugging are less direct.
\item Simplicity suits synchronous services, while async-first streaming needs different architecture.
\item SSE fits token streaming and reconnect semantics; WebSockets fit interactive duplex protocols.
\end{enumerate}
\noindent\textbf{Answer.} SSE fits token streaming and reconnect semantics; WebSockets fit interactive duplex protocols.
\par\textit{WebSockets provide bidirectional frames; server-sent events provide simpler one-way HTTP streams. Tradeoff: SSE fits token streaming and reconnect semantics; WebSockets fit interactive duplex protocols. Watch for: Using WebSockets when one-way streaming is enough adds state and proxy complexity.}
\paragraph{MCQ-1513 [hard]} Which failure mode should an interviewer expect you to identify for WebSockets and SSE?
\begin{enumerate}[label=\Alph*.]
\item Expiring keys before upstream retry windows can reintroduce duplicates.
\item Using WebSockets when one-way streaming is enough adds state and proxy complexity.
\item Changing field numbers or incompatible meanings breaks wire compatibility.
\item Assuming an async view converts a WSGI deployment into high-concurrency ASGI is wrong.
\end{enumerate}
\noindent\textbf{Answer.} Using WebSockets when one-way streaming is enough adds state and proxy complexity.
\par\textit{WebSockets provide bidirectional frames; server-sent events provide simpler one-way HTTP streams. Tradeoff: SSE fits token streaming and reconnect semantics; WebSockets fit interactive duplex protocols. Watch for: Using WebSockets when one-way streaming is enough adds state and proxy complexity.}
\paragraph{MCQ-1514 [medium]} A design review is specifically concerned with this risk: Using WebSockets when one-way streaming is enough adds state and proxy complexity. Which topic should be reviewed first?
\begin{enumerate}[label=\Alph*.]
\item API idempotency
\item gRPC
\item WebSockets and SSE
\item Flask
\end{enumerate}
\noindent\textbf{Answer.} WebSockets and SSE
\par\textit{WebSockets provide bidirectional frames; server-sent events provide simpler one-way HTTP streams. Tradeoff: SSE fits token streaming and reconnect semantics; WebSockets fit interactive duplex protocols. Watch for: Using WebSockets when one-way streaming is enough adds state and proxy complexity.}
\paragraph{MCQ-1515 [hard]} Which design note most accurately balances the benefits and costs of WebSockets and SSE?
\begin{enumerate}[label=\Alph*.]
\item Duplicate effects are prevented at the cost of durable deduplication state.
\item Typed internal services perform well, while browser and human debugging are less direct.
\item Simplicity suits synchronous services, while async-first streaming needs different architecture.
\item SSE fits token streaming and reconnect semantics; WebSockets fit interactive duplex protocols.
\end{enumerate}
\noindent\textbf{Answer.} SSE fits token streaming and reconnect semantics; WebSockets fit interactive duplex protocols.
\par\textit{WebSockets provide bidirectional frames; server-sent events provide simpler one-way HTTP streams. Tradeoff: SSE fits token streaming and reconnect semantics; WebSockets fit interactive duplex protocols. Watch for: Using WebSockets when one-way streaming is enough adds state and proxy complexity.}
\paragraph{MCQ-1516 [medium]} Which explanation would be strongest in a system-design interview about WebSockets and SSE?
\begin{enumerate}[label=\Alph*.]
\item gRPC uses protobuf contracts and HTTP/2 for efficient unary and streaming RPCs with generated clients.
\item Idempotency keys let retried mutation requests resolve to the same business operation and stored outcome.
\item WebSockets provide bidirectional frames; server-sent events provide simpler one-way HTTP streams.
\item Flask is a minimal WSGI web framework with a large extension ecosystem.
\end{enumerate}
\noindent\textbf{Answer.} WebSockets provide bidirectional frames; server-sent events provide simpler one-way HTTP streams.
\par\textit{WebSockets provide bidirectional frames; server-sent events provide simpler one-way HTTP streams. Tradeoff: SSE fits token streaming and reconnect semantics; WebSockets fit interactive duplex protocols. Watch for: Using WebSockets when one-way streaming is enough adds state and proxy complexity.}
\paragraph{MCQ-1517 [hard]} Which production warning is most directly associated with WebSockets and SSE?
\begin{enumerate}[label=\Alph*.]
\item Using WebSockets when one-way streaming is enough adds state and proxy complexity.
\item Changing field numbers or incompatible meanings breaks wire compatibility.
\item Assuming an async view converts a WSGI deployment into high-concurrency ASGI is wrong.
\item Expiring keys before upstream retry windows can reintroduce duplicates.
\end{enumerate}
\noindent\textbf{Answer.} Using WebSockets when one-way streaming is enough adds state and proxy complexity.
\par\textit{WebSockets provide bidirectional frames; server-sent events provide simpler one-way HTTP streams. Tradeoff: SSE fits token streaming and reconnect semantics; WebSockets fit interactive duplex protocols. Watch for: Using WebSockets when one-way streaming is enough adds state and proxy complexity.}
\paragraph{FC-1518 [medium]} Define WebSockets and SSE in interview-ready language.
\noindent\textbf{Answer.} WebSockets provide bidirectional frames; server-sent events provide simpler one-way HTTP streams.
\par\textit{WebSockets provide bidirectional frames; server-sent events provide simpler one-way HTTP streams.}
\paragraph{FC-1519 [hard]} State the main tradeoff and production pitfall for WebSockets and SSE.
\noindent\textbf{Answer.} Tradeoff: SSE fits token streaming and reconnect semantics; WebSockets fit interactive duplex protocols. Pitfall: Using WebSockets when one-way streaming is enough adds state and proxy complexity.
\par\textit{Tradeoff: SSE fits token streaming and reconnect semantics; WebSockets fit interactive duplex protocols. Pitfall: Using WebSockets when one-way streaming is enough adds state and proxy complexity.}
\paragraph{LA-1520 [hard]} Explain WebSockets and SSE from first principles, then show how you would evaluate and operate it in production.
\noindent\textbf{Answer.} Start with the mechanism: WebSockets provide bidirectional frames; server-sent events provide simpler one-way HTTP streams. Make the tradeoff explicit: SSE fits token streaming and reconnect semantics; WebSockets fit interactive duplex protocols. Define workload-specific offline metrics and a latency/cost budget, test important slices, instrument the critical path, and create a rollback criterion. The production review must explicitly guard against this failure: Using WebSockets when one-way streaming is enough adds state and proxy complexity.
\par\textit{A strong answer connects mechanism, measurable tradeoffs, failure isolation, observability, and rollback rather than listing product features.}
\paragraph{MCQ-1521 [medium]} Which statement best describes API idempotency?
\begin{enumerate}[label=\Alph*.]
\item WebSockets provide bidirectional frames; server-sent events provide simpler one-way HTTP streams.
\item gRPC uses protobuf contracts and HTTP/2 for efficient unary and streaming RPCs with generated clients.
\item Token buckets, leaky buckets, and concurrency limits protect shared resources by identity and cost.
\item Idempotency keys let retried mutation requests resolve to the same business operation and stored outcome.
\end{enumerate}
\noindent\textbf{Answer.} Idempotency keys let retried mutation requests resolve to the same business operation and stored outcome.
\par\textit{Idempotency keys let retried mutation requests resolve to the same business operation and stored outcome. Tradeoff: Duplicate effects are prevented at the cost of durable deduplication state. Watch for: Expiring keys before upstream retry windows can reintroduce duplicates.}
\paragraph{MCQ-1522 [hard]} What is the central engineering tradeoff of API idempotency?
\begin{enumerate}[label=\Alph*.]
\item Duplicate effects are prevented at the cost of durable deduplication state.
\item Stability and fairness improve but strict limits can reject valuable bursts.
\item Typed internal services perform well, while browser and human debugging are less direct.
\item SSE fits token streaming and reconnect semantics; WebSockets fit interactive duplex protocols.
\end{enumerate}
\noindent\textbf{Answer.} Duplicate effects are prevented at the cost of durable deduplication state.
\par\textit{Idempotency keys let retried mutation requests resolve to the same business operation and stored outcome. Tradeoff: Duplicate effects are prevented at the cost of durable deduplication state. Watch for: Expiring keys before upstream retry windows can reintroduce duplicates.}
\paragraph{MCQ-1523 [hard]} Which failure mode should an interviewer expect you to identify for API idempotency?
\begin{enumerate}[label=\Alph*.]
\item Limiting requests instead of tokens lets a few long prompts monopolize capacity.
\item Using WebSockets when one-way streaming is enough adds state and proxy complexity.
\item Expiring keys before upstream retry windows can reintroduce duplicates.
\item Changing field numbers or incompatible meanings breaks wire compatibility.
\end{enumerate}
\noindent\textbf{Answer.} Expiring keys before upstream retry windows can reintroduce duplicates.
\par\textit{Idempotency keys let retried mutation requests resolve to the same business operation and stored outcome. Tradeoff: Duplicate effects are prevented at the cost of durable deduplication state. Watch for: Expiring keys before upstream retry windows can reintroduce duplicates.}
\paragraph{MCQ-1524 [medium]} A design review is specifically concerned with this risk: Expiring keys before upstream retry windows can reintroduce duplicates. Which topic should be reviewed first?
\begin{enumerate}[label=\Alph*.]
\item gRPC
\item API idempotency
\item WebSockets and SSE
\item Rate limiting
\end{enumerate}
\noindent\textbf{Answer.} API idempotency
\par\textit{Idempotency keys let retried mutation requests resolve to the same business operation and stored outcome. Tradeoff: Duplicate effects are prevented at the cost of durable deduplication state. Watch for: Expiring keys before upstream retry windows can reintroduce duplicates.}
\paragraph{MCQ-1525 [hard]} Which design note most accurately balances the benefits and costs of API idempotency?
\begin{enumerate}[label=\Alph*.]
\item Typed internal services perform well, while browser and human debugging are less direct.
\item Duplicate effects are prevented at the cost of durable deduplication state.
\item SSE fits token streaming and reconnect semantics; WebSockets fit interactive duplex protocols.
\item Stability and fairness improve but strict limits can reject valuable bursts.
\end{enumerate}
\noindent\textbf{Answer.} Duplicate effects are prevented at the cost of durable deduplication state.
\par\textit{Idempotency keys let retried mutation requests resolve to the same business operation and stored outcome. Tradeoff: Duplicate effects are prevented at the cost of durable deduplication state. Watch for: Expiring keys before upstream retry windows can reintroduce duplicates.}
\paragraph{MCQ-1526 [medium]} Which explanation would be strongest in a system-design interview about API idempotency?
\begin{enumerate}[label=\Alph*.]
\item WebSockets provide bidirectional frames; server-sent events provide simpler one-way HTTP streams.
\item gRPC uses protobuf contracts and HTTP/2 for efficient unary and streaming RPCs with generated clients.
\item Idempotency keys let retried mutation requests resolve to the same business operation and stored outcome.
\item Token buckets, leaky buckets, and concurrency limits protect shared resources by identity and cost.
\end{enumerate}
\noindent\textbf{Answer.} Idempotency keys let retried mutation requests resolve to the same business operation and stored outcome.
\par\textit{Idempotency keys let retried mutation requests resolve to the same business operation and stored outcome. Tradeoff: Duplicate effects are prevented at the cost of durable deduplication state. Watch for: Expiring keys before upstream retry windows can reintroduce duplicates.}
\paragraph{MCQ-1527 [hard]} Which production warning is most directly associated with API idempotency?
\begin{enumerate}[label=\Alph*.]
\item Changing field numbers or incompatible meanings breaks wire compatibility.
\item Limiting requests instead of tokens lets a few long prompts monopolize capacity.
\item Using WebSockets when one-way streaming is enough adds state and proxy complexity.
\item Expiring keys before upstream retry windows can reintroduce duplicates.
\end{enumerate}
\noindent\textbf{Answer.} Expiring keys before upstream retry windows can reintroduce duplicates.
\par\textit{Idempotency keys let retried mutation requests resolve to the same business operation and stored outcome. Tradeoff: Duplicate effects are prevented at the cost of durable deduplication state. Watch for: Expiring keys before upstream retry windows can reintroduce duplicates.}
\paragraph{FC-1528 [medium]} Define API idempotency in interview-ready language.
\noindent\textbf{Answer.} Idempotency keys let retried mutation requests resolve to the same business operation and stored outcome.
\par\textit{Idempotency keys let retried mutation requests resolve to the same business operation and stored outcome.}
\paragraph{FC-1529 [hard]} State the main tradeoff and production pitfall for API idempotency.
\noindent\textbf{Answer.} Tradeoff: Duplicate effects are prevented at the cost of durable deduplication state. Pitfall: Expiring keys before upstream retry windows can reintroduce duplicates.
\par\textit{Tradeoff: Duplicate effects are prevented at the cost of durable deduplication state. Pitfall: Expiring keys before upstream retry windows can reintroduce duplicates.}
\paragraph{LA-1530 [hard]} Explain API idempotency from first principles, then show how you would evaluate and operate it in production.
\noindent\textbf{Answer.} Start with the mechanism: Idempotency keys let retried mutation requests resolve to the same business operation and stored outcome. Make the tradeoff explicit: Duplicate effects are prevented at the cost of durable deduplication state. Define workload-specific offline metrics and a latency/cost budget, test important slices, instrument the critical path, and create a rollback criterion. The production review must explicitly guard against this failure: Expiring keys before upstream retry windows can reintroduce duplicates.
\par\textit{A strong answer connects mechanism, measurable tradeoffs, failure isolation, observability, and rollback rather than listing product features.}
\paragraph{MCQ-1531 [medium]} Which statement best describes Rate limiting?
\begin{enumerate}[label=\Alph*.]
\item Token buckets, leaky buckets, and concurrency limits protect shared resources by identity and cost.
\item Boundaries should follow ownership, data, scaling, failure isolation, and change cadence rather than nouns alone.
\item Encapsulation, composition, interfaces, and explicit invariants organize stateful behavior.
\item Flask is a minimal WSGI web framework with a large extension ecosystem.
\end{enumerate}
\noindent\textbf{Answer.} Token buckets, leaky buckets, and concurrency limits protect shared resources by identity and cost.
\par\textit{Token buckets, leaky buckets, and concurrency limits protect shared resources by identity and cost. Tradeoff: Stability and fairness improve but strict limits can reject valuable bursts. Watch for: Limiting requests instead of tokens lets a few long prompts monopolize capacity.}
\paragraph{MCQ-1532 [hard]} What is the central engineering tradeoff of Rate limiting?
\begin{enumerate}[label=\Alph*.]
\item Simplicity suits synchronous services, while async-first streaming needs different architecture.
\item Abstraction improves changeability, while deep inheritance hides coupling.
\item Independent deployment helps teams but adds networks, consistency, and operations.
\item Stability and fairness improve but strict limits can reject valuable bursts.
\end{enumerate}
\noindent\textbf{Answer.} Stability and fairness improve but strict limits can reject valuable bursts.
\par\textit{Token buckets, leaky buckets, and concurrency limits protect shared resources by identity and cost. Tradeoff: Stability and fairness improve but strict limits can reject valuable bursts. Watch for: Limiting requests instead of tokens lets a few long prompts monopolize capacity.}
\paragraph{MCQ-1533 [hard]} Which failure mode should an interviewer expect you to identify for Rate limiting?
\begin{enumerate}[label=\Alph*.]
\item Assuming an async view converts a WSGI deployment into high-concurrency ASGI is wrong.
\item Splitting too early creates a distributed monolith.
\item Limiting requests instead of tokens lets a few long prompts monopolize capacity.
\item Modeling every data record as a behavior-heavy class adds ceremony without invariants.
\end{enumerate}
\noindent\textbf{Answer.} Limiting requests instead of tokens lets a few long prompts monopolize capacity.
\par\textit{Token buckets, leaky buckets, and concurrency limits protect shared resources by identity and cost. Tradeoff: Stability and fairness improve but strict limits can reject valuable bursts. Watch for: Limiting requests instead of tokens lets a few long prompts monopolize capacity.}
\paragraph{MCQ-1534 [medium]} A design review is specifically concerned with this risk: Limiting requests instead of tokens lets a few long prompts monopolize capacity. Which topic should be reviewed first?
\begin{enumerate}[label=\Alph*.]
\item Microservice boundaries
\item Flask
\item Object-oriented design
\item Rate limiting
\end{enumerate}
\noindent\textbf{Answer.} Rate limiting
\par\textit{Token buckets, leaky buckets, and concurrency limits protect shared resources by identity and cost. Tradeoff: Stability and fairness improve but strict limits can reject valuable bursts. Watch for: Limiting requests instead of tokens lets a few long prompts monopolize capacity.}
\paragraph{MCQ-1535 [hard]} Which design note most accurately balances the benefits and costs of Rate limiting?
\begin{enumerate}[label=\Alph*.]
\item Stability and fairness improve but strict limits can reject valuable bursts.
\item Abstraction improves changeability, while deep inheritance hides coupling.
\item Simplicity suits synchronous services, while async-first streaming needs different architecture.
\item Independent deployment helps teams but adds networks, consistency, and operations.
\end{enumerate}
\noindent\textbf{Answer.} Stability and fairness improve but strict limits can reject valuable bursts.
\par\textit{Token buckets, leaky buckets, and concurrency limits protect shared resources by identity and cost. Tradeoff: Stability and fairness improve but strict limits can reject valuable bursts. Watch for: Limiting requests instead of tokens lets a few long prompts monopolize capacity.}
\paragraph{MCQ-1536 [medium]} Which explanation would be strongest in a system-design interview about Rate limiting?
\begin{enumerate}[label=\Alph*.]
\item Token buckets, leaky buckets, and concurrency limits protect shared resources by identity and cost.
\item Boundaries should follow ownership, data, scaling, failure isolation, and change cadence rather than nouns alone.
\item Encapsulation, composition, interfaces, and explicit invariants organize stateful behavior.
\item Flask is a minimal WSGI web framework with a large extension ecosystem.
\end{enumerate}
\noindent\textbf{Answer.} Token buckets, leaky buckets, and concurrency limits protect shared resources by identity and cost.
\par\textit{Token buckets, leaky buckets, and concurrency limits protect shared resources by identity and cost. Tradeoff: Stability and fairness improve but strict limits can reject valuable bursts. Watch for: Limiting requests instead of tokens lets a few long prompts monopolize capacity.}
\paragraph{MCQ-1537 [hard]} Which production warning is most directly associated with Rate limiting?
\begin{enumerate}[label=\Alph*.]
\item Splitting too early creates a distributed monolith.
\item Limiting requests instead of tokens lets a few long prompts monopolize capacity.
\item Modeling every data record as a behavior-heavy class adds ceremony without invariants.
\item Assuming an async view converts a WSGI deployment into high-concurrency ASGI is wrong.
\end{enumerate}
\noindent\textbf{Answer.} Limiting requests instead of tokens lets a few long prompts monopolize capacity.
\par\textit{Token buckets, leaky buckets, and concurrency limits protect shared resources by identity and cost. Tradeoff: Stability and fairness improve but strict limits can reject valuable bursts. Watch for: Limiting requests instead of tokens lets a few long prompts monopolize capacity.}
\paragraph{FC-1538 [medium]} Define Rate limiting in interview-ready language.
\noindent\textbf{Answer.} Token buckets, leaky buckets, and concurrency limits protect shared resources by identity and cost.
\par\textit{Token buckets, leaky buckets, and concurrency limits protect shared resources by identity and cost.}
\paragraph{FC-1539 [hard]} State the main tradeoff and production pitfall for Rate limiting.
\noindent\textbf{Answer.} Tradeoff: Stability and fairness improve but strict limits can reject valuable bursts. Pitfall: Limiting requests instead of tokens lets a few long prompts monopolize capacity.
\par\textit{Tradeoff: Stability and fairness improve but strict limits can reject valuable bursts. Pitfall: Limiting requests instead of tokens lets a few long prompts monopolize capacity.}
\paragraph{LA-1540 [hard]} Explain Rate limiting from first principles, then show how you would evaluate and operate it in production.
\noindent\textbf{Answer.} Start with the mechanism: Token buckets, leaky buckets, and concurrency limits protect shared resources by identity and cost. Make the tradeoff explicit: Stability and fairness improve but strict limits can reject valuable bursts. Define workload-specific offline metrics and a latency/cost budget, test important slices, instrument the critical path, and create a rollback criterion. The production review must explicitly guard against this failure: Limiting requests instead of tokens lets a few long prompts monopolize capacity.
\par\textit{A strong answer connects mechanism, measurable tradeoffs, failure isolation, observability, and rollback rather than listing product features.}
\paragraph{MCQ-1541 [medium]} Which statement best describes Circuit breakers?
\begin{enumerate}[label=\Alph*.]
\item WebSockets provide bidirectional frames; server-sent events provide simpler one-way HTTP streams.
\item Circuit breakers stop calls to a failing dependency, probe recovery, and combine with timeouts and fallbacks.
\item Django supplies an ORM, admin, authentication, migrations, and conventions; DRF adds APIs.
\item Boundaries should follow ownership, data, scaling, failure isolation, and change cadence rather than nouns alone.
\end{enumerate}
\noindent\textbf{Answer.} Circuit breakers stop calls to a failing dependency, probe recovery, and combine with timeouts and fallbacks.
\par\textit{Circuit breakers stop calls to a failing dependency, probe recovery, and combine with timeouts and fallbacks. Tradeoff: Cascading failure is reduced, while bad thresholds can create synchronized outages. Watch for: A breaker without bounded timeouts reacts only after resources are already exhausted.}
\paragraph{MCQ-1542 [hard]} What is the central engineering tradeoff of Circuit breakers?
\begin{enumerate}[label=\Alph*.]
\item Independent deployment helps teams but adds networks, consistency, and operations.
\item SSE fits token streaming and reconnect semantics; WebSockets fit interactive duplex protocols.
\item Batteries accelerate business systems but can be heavy for a narrow inference gateway.
\item Cascading failure is reduced, while bad thresholds can create synchronized outages.
\end{enumerate}
\noindent\textbf{Answer.} Cascading failure is reduced, while bad thresholds can create synchronized outages.
\par\textit{Circuit breakers stop calls to a failing dependency, probe recovery, and combine with timeouts and fallbacks. Tradeoff: Cascading failure is reduced, while bad thresholds can create synchronized outages. Watch for: A breaker without bounded timeouts reacts only after resources are already exhausted.}
\paragraph{MCQ-1543 [hard]} Which failure mode should an interviewer expect you to identify for Circuit breakers?
\begin{enumerate}[label=\Alph*.]
\item Splitting too early creates a distributed monolith.
\item A breaker without bounded timeouts reacts only after resources are already exhausted.
\item Using WebSockets when one-way streaming is enough adds state and proxy complexity.
\item Performing N+1 ORM queries around an expensive model magnifies tail latency.
\end{enumerate}
\noindent\textbf{Answer.} A breaker without bounded timeouts reacts only after resources are already exhausted.
\par\textit{Circuit breakers stop calls to a failing dependency, probe recovery, and combine with timeouts and fallbacks. Tradeoff: Cascading failure is reduced, while bad thresholds can create synchronized outages. Watch for: A breaker without bounded timeouts reacts only after resources are already exhausted.}
\paragraph{MCQ-1544 [medium]} A design review is specifically concerned with this risk: A breaker without bounded timeouts reacts only after resources are already exhausted. Which topic should be reviewed first?
\begin{enumerate}[label=\Alph*.]
\item Django and DRF
\item WebSockets and SSE
\item Circuit breakers
\item Microservice boundaries
\end{enumerate}
\noindent\textbf{Answer.} Circuit breakers
\par\textit{Circuit breakers stop calls to a failing dependency, probe recovery, and combine with timeouts and fallbacks. Tradeoff: Cascading failure is reduced, while bad thresholds can create synchronized outages. Watch for: A breaker without bounded timeouts reacts only after resources are already exhausted.}
\paragraph{MCQ-1545 [hard]} Which design note most accurately balances the benefits and costs of Circuit breakers?
\begin{enumerate}[label=\Alph*.]
\item Cascading failure is reduced, while bad thresholds can create synchronized outages.
\item Independent deployment helps teams but adds networks, consistency, and operations.
\item SSE fits token streaming and reconnect semantics; WebSockets fit interactive duplex protocols.
\item Batteries accelerate business systems but can be heavy for a narrow inference gateway.
\end{enumerate}
\noindent\textbf{Answer.} Cascading failure is reduced, while bad thresholds can create synchronized outages.
\par\textit{Circuit breakers stop calls to a failing dependency, probe recovery, and combine with timeouts and fallbacks. Tradeoff: Cascading failure is reduced, while bad thresholds can create synchronized outages. Watch for: A breaker without bounded timeouts reacts only after resources are already exhausted.}
\paragraph{MCQ-1546 [medium]} Which explanation would be strongest in a system-design interview about Circuit breakers?
\begin{enumerate}[label=\Alph*.]
\item Circuit breakers stop calls to a failing dependency, probe recovery, and combine with timeouts and fallbacks.
\item WebSockets provide bidirectional frames; server-sent events provide simpler one-way HTTP streams.
\item Boundaries should follow ownership, data, scaling, failure isolation, and change cadence rather than nouns alone.
\item Django supplies an ORM, admin, authentication, migrations, and conventions; DRF adds APIs.
\end{enumerate}
\noindent\textbf{Answer.} Circuit breakers stop calls to a failing dependency, probe recovery, and combine with timeouts and fallbacks.
\par\textit{Circuit breakers stop calls to a failing dependency, probe recovery, and combine with timeouts and fallbacks. Tradeoff: Cascading failure is reduced, while bad thresholds can create synchronized outages. Watch for: A breaker without bounded timeouts reacts only after resources are already exhausted.}
\paragraph{MCQ-1547 [hard]} Which production warning is most directly associated with Circuit breakers?
\begin{enumerate}[label=\Alph*.]
\item Using WebSockets when one-way streaming is enough adds state and proxy complexity.
\item Performing N+1 ORM queries around an expensive model magnifies tail latency.
\item A breaker without bounded timeouts reacts only after resources are already exhausted.
\item Splitting too early creates a distributed monolith.
\end{enumerate}
\noindent\textbf{Answer.} A breaker without bounded timeouts reacts only after resources are already exhausted.
\par\textit{Circuit breakers stop calls to a failing dependency, probe recovery, and combine with timeouts and fallbacks. Tradeoff: Cascading failure is reduced, while bad thresholds can create synchronized outages. Watch for: A breaker without bounded timeouts reacts only after resources are already exhausted.}
\paragraph{FC-1548 [medium]} Define Circuit breakers in interview-ready language.
\noindent\textbf{Answer.} Circuit breakers stop calls to a failing dependency, probe recovery, and combine with timeouts and fallbacks.
\par\textit{Circuit breakers stop calls to a failing dependency, probe recovery, and combine with timeouts and fallbacks.}
\paragraph{FC-1549 [hard]} State the main tradeoff and production pitfall for Circuit breakers.
\noindent\textbf{Answer.} Tradeoff: Cascading failure is reduced, while bad thresholds can create synchronized outages. Pitfall: A breaker without bounded timeouts reacts only after resources are already exhausted.
\par\textit{Tradeoff: Cascading failure is reduced, while bad thresholds can create synchronized outages. Pitfall: A breaker without bounded timeouts reacts only after resources are already exhausted.}
\paragraph{LA-1550 [hard]} Explain Circuit breakers from first principles, then show how you would evaluate and operate it in production.
\noindent\textbf{Answer.} Start with the mechanism: Circuit breakers stop calls to a failing dependency, probe recovery, and combine with timeouts and fallbacks. Make the tradeoff explicit: Cascading failure is reduced, while bad thresholds can create synchronized outages. Define workload-specific offline metrics and a latency/cost budget, test important slices, instrument the critical path, and create a rollback criterion. The production review must explicitly guard against this failure: A breaker without bounded timeouts reacts only after resources are already exhausted.
\par\textit{A strong answer connects mechanism, measurable tradeoffs, failure isolation, observability, and rollback rather than listing product features.}
\paragraph{MCQ-1551 [medium]} Which statement best describes Microservice boundaries?
\begin{enumerate}[label=\Alph*.]
\item FastAPI builds typed ASGI APIs from Python hints with validation, dependency injection, and OpenAPI generation.
\item Boundaries should follow ownership, data, scaling, failure isolation, and change cadence rather than nouns alone.
\item gRPC uses protobuf contracts and HTTP/2 for efficient unary and streaming RPCs with generated clients.
\item Token buckets, leaky buckets, and concurrency limits protect shared resources by identity and cost.
\end{enumerate}
\noindent\textbf{Answer.} Boundaries should follow ownership, data, scaling, failure isolation, and change cadence rather than nouns alone.
\par\textit{Boundaries should follow ownership, data, scaling, failure isolation, and change cadence rather than nouns alone. Tradeoff: Independent deployment helps teams but adds networks, consistency, and operations. Watch for: Splitting too early creates a distributed monolith.}
\paragraph{MCQ-1552 [hard]} What is the central engineering tradeoff of Microservice boundaries?
\begin{enumerate}[label=\Alph*.]
\item Typed internal services perform well, while browser and human debugging are less direct.
\item Development is fast, while blocking work must be isolated from the event loop.
\item Stability and fairness improve but strict limits can reject valuable bursts.
\item Independent deployment helps teams but adds networks, consistency, and operations.
\end{enumerate}
\noindent\textbf{Answer.} Independent deployment helps teams but adds networks, consistency, and operations.
\par\textit{Boundaries should follow ownership, data, scaling, failure isolation, and change cadence rather than nouns alone. Tradeoff: Independent deployment helps teams but adds networks, consistency, and operations. Watch for: Splitting too early creates a distributed monolith.}
\paragraph{MCQ-1553 [hard]} Which failure mode should an interviewer expect you to identify for Microservice boundaries?
\begin{enumerate}[label=\Alph*.]
\item Splitting too early creates a distributed monolith.
\item Changing field numbers or incompatible meanings breaks wire compatibility.
\item Declaring a blocking database or model call inside async code stalls all requests.
\item Limiting requests instead of tokens lets a few long prompts monopolize capacity.
\end{enumerate}
\noindent\textbf{Answer.} Splitting too early creates a distributed monolith.
\par\textit{Boundaries should follow ownership, data, scaling, failure isolation, and change cadence rather than nouns alone. Tradeoff: Independent deployment helps teams but adds networks, consistency, and operations. Watch for: Splitting too early creates a distributed monolith.}
\paragraph{MCQ-1554 [medium]} A design review is specifically concerned with this risk: Splitting too early creates a distributed monolith. Which topic should be reviewed first?
\begin{enumerate}[label=\Alph*.]
\item Rate limiting
\item Microservice boundaries
\item FastAPI
\item gRPC
\end{enumerate}
\noindent\textbf{Answer.} Microservice boundaries
\par\textit{Boundaries should follow ownership, data, scaling, failure isolation, and change cadence rather than nouns alone. Tradeoff: Independent deployment helps teams but adds networks, consistency, and operations. Watch for: Splitting too early creates a distributed monolith.}
\paragraph{MCQ-1555 [hard]} Which design note most accurately balances the benefits and costs of Microservice boundaries?
\begin{enumerate}[label=\Alph*.]
\item Stability and fairness improve but strict limits can reject valuable bursts.
\item Independent deployment helps teams but adds networks, consistency, and operations.
\item Development is fast, while blocking work must be isolated from the event loop.
\item Typed internal services perform well, while browser and human debugging are less direct.
\end{enumerate}
\noindent\textbf{Answer.} Independent deployment helps teams but adds networks, consistency, and operations.
\par\textit{Boundaries should follow ownership, data, scaling, failure isolation, and change cadence rather than nouns alone. Tradeoff: Independent deployment helps teams but adds networks, consistency, and operations. Watch for: Splitting too early creates a distributed monolith.}
\paragraph{MCQ-1556 [medium]} Which explanation would be strongest in a system-design interview about Microservice boundaries?
\begin{enumerate}[label=\Alph*.]
\item Boundaries should follow ownership, data, scaling, failure isolation, and change cadence rather than nouns alone.
\item gRPC uses protobuf contracts and HTTP/2 for efficient unary and streaming RPCs with generated clients.
\item Token buckets, leaky buckets, and concurrency limits protect shared resources by identity and cost.
\item FastAPI builds typed ASGI APIs from Python hints with validation, dependency injection, and OpenAPI generation.
\end{enumerate}
\noindent\textbf{Answer.} Boundaries should follow ownership, data, scaling, failure isolation, and change cadence rather than nouns alone.
\par\textit{Boundaries should follow ownership, data, scaling, failure isolation, and change cadence rather than nouns alone. Tradeoff: Independent deployment helps teams but adds networks, consistency, and operations. Watch for: Splitting too early creates a distributed monolith.}
\paragraph{MCQ-1557 [hard]} Which production warning is most directly associated with Microservice boundaries?
\begin{enumerate}[label=\Alph*.]
\item Changing field numbers or incompatible meanings breaks wire compatibility.
\item Declaring a blocking database or model call inside async code stalls all requests.
\item Splitting too early creates a distributed monolith.
\item Limiting requests instead of tokens lets a few long prompts monopolize capacity.
\end{enumerate}
\noindent\textbf{Answer.} Splitting too early creates a distributed monolith.
\par\textit{Boundaries should follow ownership, data, scaling, failure isolation, and change cadence rather than nouns alone. Tradeoff: Independent deployment helps teams but adds networks, consistency, and operations. Watch for: Splitting too early creates a distributed monolith.}
\paragraph{FC-1558 [medium]} Define Microservice boundaries in interview-ready language.
\noindent\textbf{Answer.} Boundaries should follow ownership, data, scaling, failure isolation, and change cadence rather than nouns alone.
\par\textit{Boundaries should follow ownership, data, scaling, failure isolation, and change cadence rather than nouns alone.}
\paragraph{FC-1559 [hard]} State the main tradeoff and production pitfall for Microservice boundaries.
\noindent\textbf{Answer.} Tradeoff: Independent deployment helps teams but adds networks, consistency, and operations. Pitfall: Splitting too early creates a distributed monolith.
\par\textit{Tradeoff: Independent deployment helps teams but adds networks, consistency, and operations. Pitfall: Splitting too early creates a distributed monolith.}
\paragraph{LA-1560 [hard]} Explain Microservice boundaries from first principles, then show how you would evaluate and operate it in production.
\noindent\textbf{Answer.} Start with the mechanism: Boundaries should follow ownership, data, scaling, failure isolation, and change cadence rather than nouns alone. Make the tradeoff explicit: Independent deployment helps teams but adds networks, consistency, and operations. Define workload-specific offline metrics and a latency/cost budget, test important slices, instrument the critical path, and create a rollback criterion. The production review must explicitly guard against this failure: Splitting too early creates a distributed monolith.
\par\textit{A strong answer connects mechanism, measurable tradeoffs, failure isolation, observability, and rollback rather than listing product features.}
\paragraph{MCQ-1561 [medium]} Which statement best describes Event-driven architecture?
\begin{enumerate}[label=\Alph*.]
\item Producers publish immutable events to brokers while consumers process independently with delivery semantics.
\item WebSockets provide bidirectional frames; server-sent events provide simpler one-way HTTP streams.
\item Starlette is a lightweight ASGI toolkit providing routing, middleware, WebSockets, background tasks, and test support.
\item Django supplies an ORM, admin, authentication, migrations, and conventions; DRF adds APIs.
\end{enumerate}
\noindent\textbf{Answer.} Producers publish immutable events to brokers while consumers process independently with delivery semantics.
\par\textit{Producers publish immutable events to brokers while consumers process independently with delivery semantics. Tradeoff: Decoupling and replay improve, while ordering, duplication, schemas, and lag must be managed. Watch for: Assuming exactly-once business effects from at-least-once delivery causes duplicates.}
\paragraph{MCQ-1562 [hard]} What is the central engineering tradeoff of Event-driven architecture?
\begin{enumerate}[label=\Alph*.]
\item Decoupling and replay improve, while ordering, duplication, schemas, and lag must be managed.
\item SSE fits token streaming and reconnect semantics; WebSockets fit interactive duplex protocols.
\item It offers control with fewer batteries than a full API framework.
\item Batteries accelerate business systems but can be heavy for a narrow inference gateway.
\end{enumerate}
\noindent\textbf{Answer.} Decoupling and replay improve, while ordering, duplication, schemas, and lag must be managed.
\par\textit{Producers publish immutable events to brokers while consumers process independently with delivery semantics. Tradeoff: Decoupling and replay improve, while ordering, duplication, schemas, and lag must be managed. Watch for: Assuming exactly-once business effects from at-least-once delivery causes duplicates.}
\paragraph{MCQ-1563 [hard]} Which failure mode should an interviewer expect you to identify for Event-driven architecture?
\begin{enumerate}[label=\Alph*.]
\item Reimplementing validation and dependency patterns inconsistently increases defects.
\item Performing N+1 ORM queries around an expensive model magnifies tail latency.
\item Using WebSockets when one-way streaming is enough adds state and proxy complexity.
\item Assuming exactly-once business effects from at-least-once delivery causes duplicates.
\end{enumerate}
\noindent\textbf{Answer.} Assuming exactly-once business effects from at-least-once delivery causes duplicates.
\par\textit{Producers publish immutable events to brokers while consumers process independently with delivery semantics. Tradeoff: Decoupling and replay improve, while ordering, duplication, schemas, and lag must be managed. Watch for: Assuming exactly-once business effects from at-least-once delivery causes duplicates.}
\paragraph{MCQ-1564 [medium]} A design review is specifically concerned with this risk: Assuming exactly-once business effects from at-least-once delivery causes duplicates. Which topic should be reviewed first?
\begin{enumerate}[label=\Alph*.]
\item Event-driven architecture
\item WebSockets and SSE
\item Starlette
\item Django and DRF
\end{enumerate}
\noindent\textbf{Answer.} Event-driven architecture
\par\textit{Producers publish immutable events to brokers while consumers process independently with delivery semantics. Tradeoff: Decoupling and replay improve, while ordering, duplication, schemas, and lag must be managed. Watch for: Assuming exactly-once business effects from at-least-once delivery causes duplicates.}
\paragraph{MCQ-1565 [hard]} Which design note most accurately balances the benefits and costs of Event-driven architecture?
\begin{enumerate}[label=\Alph*.]
\item Decoupling and replay improve, while ordering, duplication, schemas, and lag must be managed.
\item It offers control with fewer batteries than a full API framework.
\item SSE fits token streaming and reconnect semantics; WebSockets fit interactive duplex protocols.
\item Batteries accelerate business systems but can be heavy for a narrow inference gateway.
\end{enumerate}
\noindent\textbf{Answer.} Decoupling and replay improve, while ordering, duplication, schemas, and lag must be managed.
\par\textit{Producers publish immutable events to brokers while consumers process independently with delivery semantics. Tradeoff: Decoupling and replay improve, while ordering, duplication, schemas, and lag must be managed. Watch for: Assuming exactly-once business effects from at-least-once delivery causes duplicates.}
\paragraph{MCQ-1566 [medium]} Which explanation would be strongest in a system-design interview about Event-driven architecture?
\begin{enumerate}[label=\Alph*.]
\item Producers publish immutable events to brokers while consumers process independently with delivery semantics.
\item WebSockets provide bidirectional frames; server-sent events provide simpler one-way HTTP streams.
\item Django supplies an ORM, admin, authentication, migrations, and conventions; DRF adds APIs.
\item Starlette is a lightweight ASGI toolkit providing routing, middleware, WebSockets, background tasks, and test support.
\end{enumerate}
\noindent\textbf{Answer.} Producers publish immutable events to brokers while consumers process independently with delivery semantics.
\par\textit{Producers publish immutable events to brokers while consumers process independently with delivery semantics. Tradeoff: Decoupling and replay improve, while ordering, duplication, schemas, and lag must be managed. Watch for: Assuming exactly-once business effects from at-least-once delivery causes duplicates.}
\paragraph{MCQ-1567 [hard]} Which production warning is most directly associated with Event-driven architecture?
\begin{enumerate}[label=\Alph*.]
\item Reimplementing validation and dependency patterns inconsistently increases defects.
\item Using WebSockets when one-way streaming is enough adds state and proxy complexity.
\item Assuming exactly-once business effects from at-least-once delivery causes duplicates.
\item Performing N+1 ORM queries around an expensive model magnifies tail latency.
\end{enumerate}
\noindent\textbf{Answer.} Assuming exactly-once business effects from at-least-once delivery causes duplicates.
\par\textit{Producers publish immutable events to brokers while consumers process independently with delivery semantics. Tradeoff: Decoupling and replay improve, while ordering, duplication, schemas, and lag must be managed. Watch for: Assuming exactly-once business effects from at-least-once delivery causes duplicates.}
\paragraph{FC-1568 [medium]} Define Event-driven architecture in interview-ready language.
\noindent\textbf{Answer.} Producers publish immutable events to brokers while consumers process independently with delivery semantics.
\par\textit{Producers publish immutable events to brokers while consumers process independently with delivery semantics.}
\paragraph{FC-1569 [hard]} State the main tradeoff and production pitfall for Event-driven architecture.
\noindent\textbf{Answer.} Tradeoff: Decoupling and replay improve, while ordering, duplication, schemas, and lag must be managed. Pitfall: Assuming exactly-once business effects from at-least-once delivery causes duplicates.
\par\textit{Tradeoff: Decoupling and replay improve, while ordering, duplication, schemas, and lag must be managed. Pitfall: Assuming exactly-once business effects from at-least-once delivery causes duplicates.}
\paragraph{LA-1570 [hard]} Explain Event-driven architecture from first principles, then show how you would evaluate and operate it in production.
\noindent\textbf{Answer.} Start with the mechanism: Producers publish immutable events to brokers while consumers process independently with delivery semantics. Make the tradeoff explicit: Decoupling and replay improve, while ordering, duplication, schemas, and lag must be managed. Define workload-specific offline metrics and a latency/cost budget, test important slices, instrument the critical path, and create a rollback criterion. The production review must explicitly guard against this failure: Assuming exactly-once business effects from at-least-once delivery causes duplicates.
\par\textit{A strong answer connects mechanism, measurable tradeoffs, failure isolation, observability, and rollback rather than listing product features.}
\paragraph{MCQ-1571 [medium]} Which statement best describes Object-oriented design?
\begin{enumerate}[label=\Alph*.]
\item Circuit breakers stop calls to a failing dependency, probe recovery, and combine with timeouts and fallbacks.
\item Token buckets, leaky buckets, and concurrency limits protect shared resources by identity and cost.
\item FastAPI builds typed ASGI APIs from Python hints with validation, dependency injection, and OpenAPI generation.
\item Encapsulation, composition, interfaces, and explicit invariants organize stateful behavior.
\end{enumerate}
\noindent\textbf{Answer.} Encapsulation, composition, interfaces, and explicit invariants organize stateful behavior.
\par\textit{Encapsulation, composition, interfaces, and explicit invariants organize stateful behavior. Tradeoff: Abstraction improves changeability, while deep inheritance hides coupling. Watch for: Modeling every data record as a behavior-heavy class adds ceremony without invariants.}
\paragraph{MCQ-1572 [hard]} What is the central engineering tradeoff of Object-oriented design?
\begin{enumerate}[label=\Alph*.]
\item Stability and fairness improve but strict limits can reject valuable bursts.
\item Development is fast, while blocking work must be isolated from the event loop.
\item Cascading failure is reduced, while bad thresholds can create synchronized outages.
\item Abstraction improves changeability, while deep inheritance hides coupling.
\end{enumerate}
\noindent\textbf{Answer.} Abstraction improves changeability, while deep inheritance hides coupling.
\par\textit{Encapsulation, composition, interfaces, and explicit invariants organize stateful behavior. Tradeoff: Abstraction improves changeability, while deep inheritance hides coupling. Watch for: Modeling every data record as a behavior-heavy class adds ceremony without invariants.}
\paragraph{MCQ-1573 [hard]} Which failure mode should an interviewer expect you to identify for Object-oriented design?
\begin{enumerate}[label=\Alph*.]
\item Limiting requests instead of tokens lets a few long prompts monopolize capacity.
\item Modeling every data record as a behavior-heavy class adds ceremony without invariants.
\item Declaring a blocking database or model call inside async code stalls all requests.
\item A breaker without bounded timeouts reacts only after resources are already exhausted.
\end{enumerate}
\noindent\textbf{Answer.} Modeling every data record as a behavior-heavy class adds ceremony without invariants.
\par\textit{Encapsulation, composition, interfaces, and explicit invariants organize stateful behavior. Tradeoff: Abstraction improves changeability, while deep inheritance hides coupling. Watch for: Modeling every data record as a behavior-heavy class adds ceremony without invariants.}
\paragraph{MCQ-1574 [medium]} A design review is specifically concerned with this risk: Modeling every data record as a behavior-heavy class adds ceremony without invariants. Which topic should be reviewed first?
\begin{enumerate}[label=\Alph*.]
\item Rate limiting
\item Circuit breakers
\item Object-oriented design
\item FastAPI
\end{enumerate}
\noindent\textbf{Answer.} Object-oriented design
\par\textit{Encapsulation, composition, interfaces, and explicit invariants organize stateful behavior. Tradeoff: Abstraction improves changeability, while deep inheritance hides coupling. Watch for: Modeling every data record as a behavior-heavy class adds ceremony without invariants.}
\paragraph{MCQ-1575 [hard]} Which design note most accurately balances the benefits and costs of Object-oriented design?
\begin{enumerate}[label=\Alph*.]
\item Abstraction improves changeability, while deep inheritance hides coupling.
\item Development is fast, while blocking work must be isolated from the event loop.
\item Stability and fairness improve but strict limits can reject valuable bursts.
\item Cascading failure is reduced, while bad thresholds can create synchronized outages.
\end{enumerate}
\noindent\textbf{Answer.} Abstraction improves changeability, while deep inheritance hides coupling.
\par\textit{Encapsulation, composition, interfaces, and explicit invariants organize stateful behavior. Tradeoff: Abstraction improves changeability, while deep inheritance hides coupling. Watch for: Modeling every data record as a behavior-heavy class adds ceremony without invariants.}
\paragraph{MCQ-1576 [medium]} Which explanation would be strongest in a system-design interview about Object-oriented design?
\begin{enumerate}[label=\Alph*.]
\item Circuit breakers stop calls to a failing dependency, probe recovery, and combine with timeouts and fallbacks.
\item Encapsulation, composition, interfaces, and explicit invariants organize stateful behavior.
\item FastAPI builds typed ASGI APIs from Python hints with validation, dependency injection, and OpenAPI generation.
\item Token buckets, leaky buckets, and concurrency limits protect shared resources by identity and cost.
\end{enumerate}
\noindent\textbf{Answer.} Encapsulation, composition, interfaces, and explicit invariants organize stateful behavior.
\par\textit{Encapsulation, composition, interfaces, and explicit invariants organize stateful behavior. Tradeoff: Abstraction improves changeability, while deep inheritance hides coupling. Watch for: Modeling every data record as a behavior-heavy class adds ceremony without invariants.}
\paragraph{MCQ-1577 [hard]} Which production warning is most directly associated with Object-oriented design?
\begin{enumerate}[label=\Alph*.]
\item Declaring a blocking database or model call inside async code stalls all requests.
\item A breaker without bounded timeouts reacts only after resources are already exhausted.
\item Limiting requests instead of tokens lets a few long prompts monopolize capacity.
\item Modeling every data record as a behavior-heavy class adds ceremony without invariants.
\end{enumerate}
\noindent\textbf{Answer.} Modeling every data record as a behavior-heavy class adds ceremony without invariants.
\par\textit{Encapsulation, composition, interfaces, and explicit invariants organize stateful behavior. Tradeoff: Abstraction improves changeability, while deep inheritance hides coupling. Watch for: Modeling every data record as a behavior-heavy class adds ceremony without invariants.}
\paragraph{FC-1578 [medium]} Define Object-oriented design in interview-ready language.
\noindent\textbf{Answer.} Encapsulation, composition, interfaces, and explicit invariants organize stateful behavior.
\par\textit{Encapsulation, composition, interfaces, and explicit invariants organize stateful behavior.}
\paragraph{FC-1579 [hard]} State the main tradeoff and production pitfall for Object-oriented design.
\noindent\textbf{Answer.} Tradeoff: Abstraction improves changeability, while deep inheritance hides coupling. Pitfall: Modeling every data record as a behavior-heavy class adds ceremony without invariants.
\par\textit{Tradeoff: Abstraction improves changeability, while deep inheritance hides coupling. Pitfall: Modeling every data record as a behavior-heavy class adds ceremony without invariants.}
\paragraph{LA-1580 [hard]} Explain Object-oriented design from first principles, then show how you would evaluate and operate it in production.
\noindent\textbf{Answer.} Start with the mechanism: Encapsulation, composition, interfaces, and explicit invariants organize stateful behavior. Make the tradeoff explicit: Abstraction improves changeability, while deep inheritance hides coupling. Define workload-specific offline metrics and a latency/cost budget, test important slices, instrument the critical path, and create a rollback criterion. The production review must explicitly guard against this failure: Modeling every data record as a behavior-heavy class adds ceremony without invariants.
\par\textit{A strong answer connects mechanism, measurable tradeoffs, failure isolation, observability, and rollback rather than listing product features.}
\subsection{Python Engineering}
\paragraph{MCQ-1581 [medium]} Which statement best describes Python threading and the GIL?
\begin{enumerate}[label=\Alph*.]
\item ThreadPoolExecutor and ProcessPoolExecutor expose a common future-based API for concurrent callables.
\item CPython combines reference counting with cyclic garbage collection; object graphs and native buffers affect real memory.
\item Generators suspend and resume execution around yield, enabling lazy iteration and streaming pipelines.
\item CPython threads share memory; the GIL limits simultaneous Python bytecode but releases around many blocking operations.
\end{enumerate}
\noindent\textbf{Answer.} CPython threads share memory; the GIL limits simultaneous Python bytecode but releases around many blocking operations.
\par\textit{CPython threads share memory; the GIL limits simultaneous Python bytecode but releases around many blocking operations. Tradeoff: Threads suit I/O-bound work and shared state, not pure Python CPU scaling. Watch for: Claiming threads are always useless ignores network and native-extension concurrency.}
\paragraph{MCQ-1582 [hard]} What is the central engineering tradeoff of Python threading and the GIL?
\begin{enumerate}[label=\Alph*.]
\item Automatic cleanup is convenient, while cycles and caches can retain large objects.
\item Memory use falls, while one-shot iteration and cleanup require care.
\item Threads suit I/O-bound work and shared state, not pure Python CPU scaling.
\item Integration is simple, but cancellation and process serialization have limits.
\end{enumerate}
\noindent\textbf{Answer.} Threads suit I/O-bound work and shared state, not pure Python CPU scaling.
\par\textit{CPython threads share memory; the GIL limits simultaneous Python bytecode but releases around many blocking operations. Tradeoff: Threads suit I/O-bound work and shared state, not pure Python CPU scaling. Watch for: Claiming threads are always useless ignores network and native-extension concurrency.}
\paragraph{MCQ-1583 [hard]} Which failure mode should an interviewer expect you to identify for Python threading and the GIL?
\begin{enumerate}[label=\Alph*.]
\item Submitting unbounded tasks creates memory pressure and queue latency.
\item Claiming threads are always useless ignores network and native-extension concurrency.
\item Reusing an exhausted generator returns no data without warning.
\item Watching only Python heap misses tensors allocated in native or GPU memory.
\end{enumerate}
\noindent\textbf{Answer.} Claiming threads are always useless ignores network and native-extension concurrency.
\par\textit{CPython threads share memory; the GIL limits simultaneous Python bytecode but releases around many blocking operations. Tradeoff: Threads suit I/O-bound work and shared state, not pure Python CPU scaling. Watch for: Claiming threads are always useless ignores network and native-extension concurrency.}
\paragraph{MCQ-1584 [medium]} A design review is specifically concerned with this risk: Claiming threads are always useless ignores network and native-extension concurrency. Which topic should be reviewed first?
\begin{enumerate}[label=\Alph*.]
\item Python threading and the GIL
\item Generators
\item Concurrent futures
\item Memory management
\end{enumerate}
\noindent\textbf{Answer.} Python threading and the GIL
\par\textit{CPython threads share memory; the GIL limits simultaneous Python bytecode but releases around many blocking operations. Tradeoff: Threads suit I/O-bound work and shared state, not pure Python CPU scaling. Watch for: Claiming threads are always useless ignores network and native-extension concurrency.}
\paragraph{MCQ-1585 [hard]} Which design note most accurately balances the benefits and costs of Python threading and the GIL?
\begin{enumerate}[label=\Alph*.]
\item Threads suit I/O-bound work and shared state, not pure Python CPU scaling.
\item Integration is simple, but cancellation and process serialization have limits.
\item Memory use falls, while one-shot iteration and cleanup require care.
\item Automatic cleanup is convenient, while cycles and caches can retain large objects.
\end{enumerate}
\noindent\textbf{Answer.} Threads suit I/O-bound work and shared state, not pure Python CPU scaling.
\par\textit{CPython threads share memory; the GIL limits simultaneous Python bytecode but releases around many blocking operations. Tradeoff: Threads suit I/O-bound work and shared state, not pure Python CPU scaling. Watch for: Claiming threads are always useless ignores network and native-extension concurrency.}
\paragraph{MCQ-1586 [medium]} Which explanation would be strongest in a system-design interview about Python threading and the GIL?
\begin{enumerate}[label=\Alph*.]
\item ThreadPoolExecutor and ProcessPoolExecutor expose a common future-based API for concurrent callables.
\item CPython threads share memory; the GIL limits simultaneous Python bytecode but releases around many blocking operations.
\item Generators suspend and resume execution around yield, enabling lazy iteration and streaming pipelines.
\item CPython combines reference counting with cyclic garbage collection; object graphs and native buffers affect real memory.
\end{enumerate}
\noindent\textbf{Answer.} CPython threads share memory; the GIL limits simultaneous Python bytecode but releases around many blocking operations.
\par\textit{CPython threads share memory; the GIL limits simultaneous Python bytecode but releases around many blocking operations. Tradeoff: Threads suit I/O-bound work and shared state, not pure Python CPU scaling. Watch for: Claiming threads are always useless ignores network and native-extension concurrency.}
\paragraph{MCQ-1587 [hard]} Which production warning is most directly associated with Python threading and the GIL?
\begin{enumerate}[label=\Alph*.]
\item Reusing an exhausted generator returns no data without warning.
\item Watching only Python heap misses tensors allocated in native or GPU memory.
\item Submitting unbounded tasks creates memory pressure and queue latency.
\item Claiming threads are always useless ignores network and native-extension concurrency.
\end{enumerate}
\noindent\textbf{Answer.} Claiming threads are always useless ignores network and native-extension concurrency.
\par\textit{CPython threads share memory; the GIL limits simultaneous Python bytecode but releases around many blocking operations. Tradeoff: Threads suit I/O-bound work and shared state, not pure Python CPU scaling. Watch for: Claiming threads are always useless ignores network and native-extension concurrency.}
\paragraph{FC-1588 [medium]} Define Python threading and the GIL in interview-ready language.
\noindent\textbf{Answer.} CPython threads share memory; the GIL limits simultaneous Python bytecode but releases around many blocking operations.
\par\textit{CPython threads share memory; the GIL limits simultaneous Python bytecode but releases around many blocking operations.}
\paragraph{FC-1589 [hard]} State the main tradeoff and production pitfall for Python threading and the GIL.
\noindent\textbf{Answer.} Tradeoff: Threads suit I/O-bound work and shared state, not pure Python CPU scaling. Pitfall: Claiming threads are always useless ignores network and native-extension concurrency.
\par\textit{Tradeoff: Threads suit I/O-bound work and shared state, not pure Python CPU scaling. Pitfall: Claiming threads are always useless ignores network and native-extension concurrency.}
\paragraph{LA-1590 [hard]} Explain Python threading and the GIL from first principles, then show how you would evaluate and operate it in production.
\noindent\textbf{Answer.} Start with the mechanism: CPython threads share memory; the GIL limits simultaneous Python bytecode but releases around many blocking operations. Make the tradeoff explicit: Threads suit I/O-bound work and shared state, not pure Python CPU scaling. Define workload-specific offline metrics and a latency/cost budget, test important slices, instrument the critical path, and create a rollback criterion. The production review must explicitly guard against this failure: Claiming threads are always useless ignores network and native-extension concurrency.
\par\textit{A strong answer connects mechanism, measurable tradeoffs, failure isolation, observability, and rollback rather than listing product features.}
\paragraph{MCQ-1591 [medium]} Which statement best describes Multiprocessing?
\begin{enumerate}[label=\Alph*.]
\item Unit, integration, property, contract, and load tests cover different risks; mocks isolate only well-understood boundaries.
\item Separate processes bypass the GIL and isolate memory while communicating through serialization or shared memory.
\item CPython threads share memory; the GIL limits simultaneous Python bytecode but releases around many blocking operations.
\item Generators suspend and resume execution around yield, enabling lazy iteration and streaming pipelines.
\end{enumerate}
\noindent\textbf{Answer.} Separate processes bypass the GIL and isolate memory while communicating through serialization or shared memory.
\par\textit{Separate processes bypass the GIL and isolate memory while communicating through serialization or shared memory. Tradeoff: CPU parallelism improves at startup, memory, and coordination cost. Watch for: Passing large models to spawned workers duplicates memory unexpectedly.}
\paragraph{MCQ-1592 [hard]} What is the central engineering tradeoff of Multiprocessing?
\begin{enumerate}[label=\Alph*.]
\item Fast tests aid iteration, while excessive mocking tests implementation rather than behavior.
\item Threads suit I/O-bound work and shared state, not pure Python CPU scaling.
\item Memory use falls, while one-shot iteration and cleanup require care.
\item CPU parallelism improves at startup, memory, and coordination cost.
\end{enumerate}
\noindent\textbf{Answer.} CPU parallelism improves at startup, memory, and coordination cost.
\par\textit{Separate processes bypass the GIL and isolate memory while communicating through serialization or shared memory. Tradeoff: CPU parallelism improves at startup, memory, and coordination cost. Watch for: Passing large models to spawned workers duplicates memory unexpectedly.}
\paragraph{MCQ-1593 [hard]} Which failure mode should an interviewer expect you to identify for Multiprocessing?
\begin{enumerate}[label=\Alph*.]
\item Reusing an exhausted generator returns no data without warning.
\item Claiming threads are always useless ignores network and native-extension concurrency.
\item Passing large models to spawned workers duplicates memory unexpectedly.
\item Mocking the model, database, and network in the same test proves little about integration.
\end{enumerate}
\noindent\textbf{Answer.} Passing large models to spawned workers duplicates memory unexpectedly.
\par\textit{Separate processes bypass the GIL and isolate memory while communicating through serialization or shared memory. Tradeoff: CPU parallelism improves at startup, memory, and coordination cost. Watch for: Passing large models to spawned workers duplicates memory unexpectedly.}
\paragraph{MCQ-1594 [medium]} A design review is specifically concerned with this risk: Passing large models to spawned workers duplicates memory unexpectedly. Which topic should be reviewed first?
\begin{enumerate}[label=\Alph*.]
\item Testing and mocking
\item Python threading and the GIL
\item Multiprocessing
\item Generators
\end{enumerate}
\noindent\textbf{Answer.} Multiprocessing
\par\textit{Separate processes bypass the GIL and isolate memory while communicating through serialization or shared memory. Tradeoff: CPU parallelism improves at startup, memory, and coordination cost. Watch for: Passing large models to spawned workers duplicates memory unexpectedly.}
\paragraph{MCQ-1595 [hard]} Which design note most accurately balances the benefits and costs of Multiprocessing?
\begin{enumerate}[label=\Alph*.]
\item CPU parallelism improves at startup, memory, and coordination cost.
\item Threads suit I/O-bound work and shared state, not pure Python CPU scaling.
\item Memory use falls, while one-shot iteration and cleanup require care.
\item Fast tests aid iteration, while excessive mocking tests implementation rather than behavior.
\end{enumerate}
\noindent\textbf{Answer.} CPU parallelism improves at startup, memory, and coordination cost.
\par\textit{Separate processes bypass the GIL and isolate memory while communicating through serialization or shared memory. Tradeoff: CPU parallelism improves at startup, memory, and coordination cost. Watch for: Passing large models to spawned workers duplicates memory unexpectedly.}
\paragraph{MCQ-1596 [medium]} Which explanation would be strongest in a system-design interview about Multiprocessing?
\begin{enumerate}[label=\Alph*.]
\item Unit, integration, property, contract, and load tests cover different risks; mocks isolate only well-understood boundaries.
\item Generators suspend and resume execution around yield, enabling lazy iteration and streaming pipelines.
\item Separate processes bypass the GIL and isolate memory while communicating through serialization or shared memory.
\item CPython threads share memory; the GIL limits simultaneous Python bytecode but releases around many blocking operations.
\end{enumerate}
\noindent\textbf{Answer.} Separate processes bypass the GIL and isolate memory while communicating through serialization or shared memory.
\par\textit{Separate processes bypass the GIL and isolate memory while communicating through serialization or shared memory. Tradeoff: CPU parallelism improves at startup, memory, and coordination cost. Watch for: Passing large models to spawned workers duplicates memory unexpectedly.}
\paragraph{MCQ-1597 [hard]} Which production warning is most directly associated with Multiprocessing?
\begin{enumerate}[label=\Alph*.]
\item Passing large models to spawned workers duplicates memory unexpectedly.
\item Mocking the model, database, and network in the same test proves little about integration.
\item Claiming threads are always useless ignores network and native-extension concurrency.
\item Reusing an exhausted generator returns no data without warning.
\end{enumerate}
\noindent\textbf{Answer.} Passing large models to spawned workers duplicates memory unexpectedly.
\par\textit{Separate processes bypass the GIL and isolate memory while communicating through serialization or shared memory. Tradeoff: CPU parallelism improves at startup, memory, and coordination cost. Watch for: Passing large models to spawned workers duplicates memory unexpectedly.}
\paragraph{FC-1598 [medium]} Define Multiprocessing in interview-ready language.
\noindent\textbf{Answer.} Separate processes bypass the GIL and isolate memory while communicating through serialization or shared memory.
\par\textit{Separate processes bypass the GIL and isolate memory while communicating through serialization or shared memory.}
\paragraph{FC-1599 [hard]} State the main tradeoff and production pitfall for Multiprocessing.
\noindent\textbf{Answer.} Tradeoff: CPU parallelism improves at startup, memory, and coordination cost. Pitfall: Passing large models to spawned workers duplicates memory unexpectedly.
\par\textit{Tradeoff: CPU parallelism improves at startup, memory, and coordination cost. Pitfall: Passing large models to spawned workers duplicates memory unexpectedly.}
\paragraph{LA-1600 [hard]} Explain Multiprocessing from first principles, then show how you would evaluate and operate it in production.
\noindent\textbf{Answer.} Start with the mechanism: Separate processes bypass the GIL and isolate memory while communicating through serialization or shared memory. Make the tradeoff explicit: CPU parallelism improves at startup, memory, and coordination cost. Define workload-specific offline metrics and a latency/cost budget, test important slices, instrument the critical path, and create a rollback criterion. The production review must explicitly guard against this failure: Passing large models to spawned workers duplicates memory unexpectedly.
\par\textit{A strong answer connects mechanism, measurable tradeoffs, failure isolation, observability, and rollback rather than listing product features.}
\paragraph{MCQ-1601 [medium]} Which statement best describes Asyncio?
\begin{enumerate}[label=\Alph*.]
\item Exceptions separate error propagation from normal return values and can preserve causal chains.
\item Asyncio cooperatively schedules coroutines on an event loop around awaitable I/O.
\item Unit, integration, property, contract, and load tests cover different risks; mocks isolate only well-understood boundaries.
\item Type hints support static analysis; Protocol defines structural interfaces without inheritance.
\end{enumerate}
\noindent\textbf{Answer.} Asyncio cooperatively schedules coroutines on an event loop around awaitable I/O.
\par\textit{Asyncio cooperatively schedules coroutines on an event loop around awaitable I/O. Tradeoff: High connection concurrency is efficient, while blocking functions freeze the loop. Watch for: Creating coroutines without awaiting them silently drops work.}
\paragraph{MCQ-1602 [hard]} What is the central engineering tradeoff of Asyncio?
\begin{enumerate}[label=\Alph*.]
\item High connection concurrency is efficient, while blocking functions freeze the loop.
\item Central handling simplifies APIs, while overly broad catches erase actionable context.
\item Fast tests aid iteration, while excessive mocking tests implementation rather than behavior.
\item Refactoring and contracts improve, while runtime behavior is unchanged unless validated.
\end{enumerate}
\noindent\textbf{Answer.} High connection concurrency is efficient, while blocking functions freeze the loop.
\par\textit{Asyncio cooperatively schedules coroutines on an event loop around awaitable I/O. Tradeoff: High connection concurrency is efficient, while blocking functions freeze the loop. Watch for: Creating coroutines without awaiting them silently drops work.}
\paragraph{MCQ-1603 [hard]} Which failure mode should an interviewer expect you to identify for Asyncio?
\begin{enumerate}[label=\Alph*.]
\item Assuming type annotations enforce production inputs creates vulnerabilities.
\item Catching Exception and returning None turns failures into corrupt state.
\item Creating coroutines without awaiting them silently drops work.
\item Mocking the model, database, and network in the same test proves little about integration.
\end{enumerate}
\noindent\textbf{Answer.} Creating coroutines without awaiting them silently drops work.
\par\textit{Asyncio cooperatively schedules coroutines on an event loop around awaitable I/O. Tradeoff: High connection concurrency is efficient, while blocking functions freeze the loop. Watch for: Creating coroutines without awaiting them silently drops work.}
\paragraph{MCQ-1604 [medium]} A design review is specifically concerned with this risk: Creating coroutines without awaiting them silently drops work. Which topic should be reviewed first?
\begin{enumerate}[label=\Alph*.]
\item Testing and mocking
\item Type hints and protocols
\item Exceptions
\item Asyncio
\end{enumerate}
\noindent\textbf{Answer.} Asyncio
\par\textit{Asyncio cooperatively schedules coroutines on an event loop around awaitable I/O. Tradeoff: High connection concurrency is efficient, while blocking functions freeze the loop. Watch for: Creating coroutines without awaiting them silently drops work.}
\paragraph{MCQ-1605 [hard]} Which design note most accurately balances the benefits and costs of Asyncio?
\begin{enumerate}[label=\Alph*.]
\item Fast tests aid iteration, while excessive mocking tests implementation rather than behavior.
\item Refactoring and contracts improve, while runtime behavior is unchanged unless validated.
\item High connection concurrency is efficient, while blocking functions freeze the loop.
\item Central handling simplifies APIs, while overly broad catches erase actionable context.
\end{enumerate}
\noindent\textbf{Answer.} High connection concurrency is efficient, while blocking functions freeze the loop.
\par\textit{Asyncio cooperatively schedules coroutines on an event loop around awaitable I/O. Tradeoff: High connection concurrency is efficient, while blocking functions freeze the loop. Watch for: Creating coroutines without awaiting them silently drops work.}
\paragraph{MCQ-1606 [medium]} Which explanation would be strongest in a system-design interview about Asyncio?
\begin{enumerate}[label=\Alph*.]
\item Type hints support static analysis; Protocol defines structural interfaces without inheritance.
\item Unit, integration, property, contract, and load tests cover different risks; mocks isolate only well-understood boundaries.
\item Asyncio cooperatively schedules coroutines on an event loop around awaitable I/O.
\item Exceptions separate error propagation from normal return values and can preserve causal chains.
\end{enumerate}
\noindent\textbf{Answer.} Asyncio cooperatively schedules coroutines on an event loop around awaitable I/O.
\par\textit{Asyncio cooperatively schedules coroutines on an event loop around awaitable I/O. Tradeoff: High connection concurrency is efficient, while blocking functions freeze the loop. Watch for: Creating coroutines without awaiting them silently drops work.}
\paragraph{MCQ-1607 [hard]} Which production warning is most directly associated with Asyncio?
\begin{enumerate}[label=\Alph*.]
\item Mocking the model, database, and network in the same test proves little about integration.
\item Assuming type annotations enforce production inputs creates vulnerabilities.
\item Catching Exception and returning None turns failures into corrupt state.
\item Creating coroutines without awaiting them silently drops work.
\end{enumerate}
\noindent\textbf{Answer.} Creating coroutines without awaiting them silently drops work.
\par\textit{Asyncio cooperatively schedules coroutines on an event loop around awaitable I/O. Tradeoff: High connection concurrency is efficient, while blocking functions freeze the loop. Watch for: Creating coroutines without awaiting them silently drops work.}
\paragraph{FC-1608 [medium]} Define Asyncio in interview-ready language.
\noindent\textbf{Answer.} Asyncio cooperatively schedules coroutines on an event loop around awaitable I/O.
\par\textit{Asyncio cooperatively schedules coroutines on an event loop around awaitable I/O.}
\paragraph{FC-1609 [hard]} State the main tradeoff and production pitfall for Asyncio.
\noindent\textbf{Answer.} Tradeoff: High connection concurrency is efficient, while blocking functions freeze the loop. Pitfall: Creating coroutines without awaiting them silently drops work.
\par\textit{Tradeoff: High connection concurrency is efficient, while blocking functions freeze the loop. Pitfall: Creating coroutines without awaiting them silently drops work.}
\paragraph{LA-1610 [hard]} Explain Asyncio from first principles, then show how you would evaluate and operate it in production.
\noindent\textbf{Answer.} Start with the mechanism: Asyncio cooperatively schedules coroutines on an event loop around awaitable I/O. Make the tradeoff explicit: High connection concurrency is efficient, while blocking functions freeze the loop. Define workload-specific offline metrics and a latency/cost budget, test important slices, instrument the critical path, and create a rollback criterion. The production review must explicitly guard against this failure: Creating coroutines without awaiting them silently drops work.
\par\textit{A strong answer connects mechanism, measurable tradeoffs, failure isolation, observability, and rollback rather than listing product features.}
\paragraph{MCQ-1611 [medium]} Which statement best describes Concurrent futures?
\begin{enumerate}[label=\Alph*.]
\item pyproject.toml, lockfiles, wheels, virtual environments, and reproducible builds define deliverable Python software.
\item Exceptions separate error propagation from normal return values and can preserve causal chains.
\item Type hints support static analysis; Protocol defines structural interfaces without inheritance.
\item ThreadPoolExecutor and ProcessPoolExecutor expose a common future-based API for concurrent callables.
\end{enumerate}
\noindent\textbf{Answer.} ThreadPoolExecutor and ProcessPoolExecutor expose a common future-based API for concurrent callables.
\par\textit{ThreadPoolExecutor and ProcessPoolExecutor expose a common future-based API for concurrent callables. Tradeoff: Integration is simple, but cancellation and process serialization have limits. Watch for: Submitting unbounded tasks creates memory pressure and queue latency.}
\paragraph{MCQ-1612 [hard]} What is the central engineering tradeoff of Concurrent futures?
\begin{enumerate}[label=\Alph*.]
\item Isolation prevents dependency drift but needs platform and supply-chain controls.
\item Central handling simplifies APIs, while overly broad catches erase actionable context.
\item Refactoring and contracts improve, while runtime behavior is unchanged unless validated.
\item Integration is simple, but cancellation and process serialization have limits.
\end{enumerate}
\noindent\textbf{Answer.} Integration is simple, but cancellation and process serialization have limits.
\par\textit{ThreadPoolExecutor and ProcessPoolExecutor expose a common future-based API for concurrent callables. Tradeoff: Integration is simple, but cancellation and process serialization have limits. Watch for: Submitting unbounded tasks creates memory pressure and queue latency.}
\paragraph{MCQ-1613 [hard]} Which failure mode should an interviewer expect you to identify for Concurrent futures?
\begin{enumerate}[label=\Alph*.]
\item Loose version ranges in production make rebuilds non-reproducible.
\item Catching Exception and returning None turns failures into corrupt state.
\item Assuming type annotations enforce production inputs creates vulnerabilities.
\item Submitting unbounded tasks creates memory pressure and queue latency.
\end{enumerate}
\noindent\textbf{Answer.} Submitting unbounded tasks creates memory pressure and queue latency.
\par\textit{ThreadPoolExecutor and ProcessPoolExecutor expose a common future-based API for concurrent callables. Tradeoff: Integration is simple, but cancellation and process serialization have limits. Watch for: Submitting unbounded tasks creates memory pressure and queue latency.}
\paragraph{MCQ-1614 [medium]} A design review is specifically concerned with this risk: Submitting unbounded tasks creates memory pressure and queue latency. Which topic should be reviewed first?
\begin{enumerate}[label=\Alph*.]
\item Concurrent futures
\item Type hints and protocols
\item Packaging and environments
\item Exceptions
\end{enumerate}
\noindent\textbf{Answer.} Concurrent futures
\par\textit{ThreadPoolExecutor and ProcessPoolExecutor expose a common future-based API for concurrent callables. Tradeoff: Integration is simple, but cancellation and process serialization have limits. Watch for: Submitting unbounded tasks creates memory pressure and queue latency.}
\paragraph{MCQ-1615 [hard]} Which design note most accurately balances the benefits and costs of Concurrent futures?
\begin{enumerate}[label=\Alph*.]
\item Isolation prevents dependency drift but needs platform and supply-chain controls.
\item Central handling simplifies APIs, while overly broad catches erase actionable context.
\item Refactoring and contracts improve, while runtime behavior is unchanged unless validated.
\item Integration is simple, but cancellation and process serialization have limits.
\end{enumerate}
\noindent\textbf{Answer.} Integration is simple, but cancellation and process serialization have limits.
\par\textit{ThreadPoolExecutor and ProcessPoolExecutor expose a common future-based API for concurrent callables. Tradeoff: Integration is simple, but cancellation and process serialization have limits. Watch for: Submitting unbounded tasks creates memory pressure and queue latency.}
\paragraph{MCQ-1616 [medium]} Which explanation would be strongest in a system-design interview about Concurrent futures?
\begin{enumerate}[label=\Alph*.]
\item Exceptions separate error propagation from normal return values and can preserve causal chains.
\item ThreadPoolExecutor and ProcessPoolExecutor expose a common future-based API for concurrent callables.
\item Type hints support static analysis; Protocol defines structural interfaces without inheritance.
\item pyproject.toml, lockfiles, wheels, virtual environments, and reproducible builds define deliverable Python software.
\end{enumerate}
\noindent\textbf{Answer.} ThreadPoolExecutor and ProcessPoolExecutor expose a common future-based API for concurrent callables.
\par\textit{ThreadPoolExecutor and ProcessPoolExecutor expose a common future-based API for concurrent callables. Tradeoff: Integration is simple, but cancellation and process serialization have limits. Watch for: Submitting unbounded tasks creates memory pressure and queue latency.}
\paragraph{MCQ-1617 [hard]} Which production warning is most directly associated with Concurrent futures?
\begin{enumerate}[label=\Alph*.]
\item Loose version ranges in production make rebuilds non-reproducible.
\item Submitting unbounded tasks creates memory pressure and queue latency.
\item Assuming type annotations enforce production inputs creates vulnerabilities.
\item Catching Exception and returning None turns failures into corrupt state.
\end{enumerate}
\noindent\textbf{Answer.} Submitting unbounded tasks creates memory pressure and queue latency.
\par\textit{ThreadPoolExecutor and ProcessPoolExecutor expose a common future-based API for concurrent callables. Tradeoff: Integration is simple, but cancellation and process serialization have limits. Watch for: Submitting unbounded tasks creates memory pressure and queue latency.}
\paragraph{FC-1618 [medium]} Define Concurrent futures in interview-ready language.
\noindent\textbf{Answer.} ThreadPoolExecutor and ProcessPoolExecutor expose a common future-based API for concurrent callables.
\par\textit{ThreadPoolExecutor and ProcessPoolExecutor expose a common future-based API for concurrent callables.}
\paragraph{FC-1619 [hard]} State the main tradeoff and production pitfall for Concurrent futures.
\noindent\textbf{Answer.} Tradeoff: Integration is simple, but cancellation and process serialization have limits. Pitfall: Submitting unbounded tasks creates memory pressure and queue latency.
\par\textit{Tradeoff: Integration is simple, but cancellation and process serialization have limits. Pitfall: Submitting unbounded tasks creates memory pressure and queue latency.}
\paragraph{LA-1620 [hard]} Explain Concurrent futures from first principles, then show how you would evaluate and operate it in production.
\noindent\textbf{Answer.} Start with the mechanism: ThreadPoolExecutor and ProcessPoolExecutor expose a common future-based API for concurrent callables. Make the tradeoff explicit: Integration is simple, but cancellation and process serialization have limits. Define workload-specific offline metrics and a latency/cost budget, test important slices, instrument the critical path, and create a rollback criterion. The production review must explicitly guard against this failure: Submitting unbounded tasks creates memory pressure and queue latency.
\par\textit{A strong answer connects mechanism, measurable tradeoffs, failure isolation, observability, and rollback rather than listing product features.}
\paragraph{MCQ-1621 [medium]} Which statement best describes Decorators?
\begin{enumerate}[label=\Alph*.]
\item CPython combines reference counting with cyclic garbage collection; object graphs and native buffers affect real memory.
\item Context managers guarantee paired setup and cleanup through with, \_\_enter\_\_/\_\_exit\_\_, or contextlib.
\item Descriptors define attribute access; properties provide managed getter, setter, and deleter behavior.
\item Decorators wrap or replace functions or classes to add behavior while preserving a callable interface.
\end{enumerate}
\noindent\textbf{Answer.} Decorators wrap or replace functions or classes to add behavior while preserving a callable interface.
\par\textit{Decorators wrap or replace functions or classes to add behavior while preserving a callable interface. Tradeoff: Cross-cutting concerns become reusable, but hidden control flow complicates debugging. Watch for: Omitting functools.wraps breaks metadata and framework introspection.}
\paragraph{MCQ-1622 [hard]} What is the central engineering tradeoff of Decorators?
\begin{enumerate}[label=\Alph*.]
\item Validation and lazy access fit attribute syntax but may hide expensive operations.
\item Cross-cutting concerns become reusable, but hidden control flow complicates debugging.
\item Automatic cleanup is convenient, while cycles and caches can retain large objects.
\item Resource safety improves and exceptions are localized.
\end{enumerate}
\noindent\textbf{Answer.} Cross-cutting concerns become reusable, but hidden control flow complicates debugging.
\par\textit{Decorators wrap or replace functions or classes to add behavior while preserving a callable interface. Tradeoff: Cross-cutting concerns become reusable, but hidden control flow complicates debugging. Watch for: Omitting functools.wraps breaks metadata and framework introspection.}
\paragraph{MCQ-1623 [hard]} Which failure mode should an interviewer expect you to identify for Decorators?
\begin{enumerate}[label=\Alph*.]
\item Performing network I/O in a property surprises callers and tooling.
\item Watching only Python heap misses tensors allocated in native or GPU memory.
\item Swallowing exceptions unintentionally in \_\_exit\_\_ hides failures.
\item Omitting functools.wraps breaks metadata and framework introspection.
\end{enumerate}
\noindent\textbf{Answer.} Omitting functools.wraps breaks metadata and framework introspection.
\par\textit{Decorators wrap or replace functions or classes to add behavior while preserving a callable interface. Tradeoff: Cross-cutting concerns become reusable, but hidden control flow complicates debugging. Watch for: Omitting functools.wraps breaks metadata and framework introspection.}
\paragraph{MCQ-1624 [medium]} A design review is specifically concerned with this risk: Omitting functools.wraps breaks metadata and framework introspection. Which topic should be reviewed first?
\begin{enumerate}[label=\Alph*.]
\item Descriptors and properties
\item Memory management
\item Context managers
\item Decorators
\end{enumerate}
\noindent\textbf{Answer.} Decorators
\par\textit{Decorators wrap or replace functions or classes to add behavior while preserving a callable interface. Tradeoff: Cross-cutting concerns become reusable, but hidden control flow complicates debugging. Watch for: Omitting functools.wraps breaks metadata and framework introspection.}
\paragraph{MCQ-1625 [hard]} Which design note most accurately balances the benefits and costs of Decorators?
\begin{enumerate}[label=\Alph*.]
\item Validation and lazy access fit attribute syntax but may hide expensive operations.
\item Cross-cutting concerns become reusable, but hidden control flow complicates debugging.
\item Resource safety improves and exceptions are localized.
\item Automatic cleanup is convenient, while cycles and caches can retain large objects.
\end{enumerate}
\noindent\textbf{Answer.} Cross-cutting concerns become reusable, but hidden control flow complicates debugging.
\par\textit{Decorators wrap or replace functions or classes to add behavior while preserving a callable interface. Tradeoff: Cross-cutting concerns become reusable, but hidden control flow complicates debugging. Watch for: Omitting functools.wraps breaks metadata and framework introspection.}
\paragraph{MCQ-1626 [medium]} Which explanation would be strongest in a system-design interview about Decorators?
\begin{enumerate}[label=\Alph*.]
\item Decorators wrap or replace functions or classes to add behavior while preserving a callable interface.
\item CPython combines reference counting with cyclic garbage collection; object graphs and native buffers affect real memory.
\item Context managers guarantee paired setup and cleanup through with, \_\_enter\_\_/\_\_exit\_\_, or contextlib.
\item Descriptors define attribute access; properties provide managed getter, setter, and deleter behavior.
\end{enumerate}
\noindent\textbf{Answer.} Decorators wrap or replace functions or classes to add behavior while preserving a callable interface.
\par\textit{Decorators wrap or replace functions or classes to add behavior while preserving a callable interface. Tradeoff: Cross-cutting concerns become reusable, but hidden control flow complicates debugging. Watch for: Omitting functools.wraps breaks metadata and framework introspection.}
\paragraph{MCQ-1627 [hard]} Which production warning is most directly associated with Decorators?
\begin{enumerate}[label=\Alph*.]
\item Watching only Python heap misses tensors allocated in native or GPU memory.
\item Swallowing exceptions unintentionally in \_\_exit\_\_ hides failures.
\item Omitting functools.wraps breaks metadata and framework introspection.
\item Performing network I/O in a property surprises callers and tooling.
\end{enumerate}
\noindent\textbf{Answer.} Omitting functools.wraps breaks metadata and framework introspection.
\par\textit{Decorators wrap or replace functions or classes to add behavior while preserving a callable interface. Tradeoff: Cross-cutting concerns become reusable, but hidden control flow complicates debugging. Watch for: Omitting functools.wraps breaks metadata and framework introspection.}
\paragraph{FC-1628 [medium]} Define Decorators in interview-ready language.
\noindent\textbf{Answer.} Decorators wrap or replace functions or classes to add behavior while preserving a callable interface.
\par\textit{Decorators wrap or replace functions or classes to add behavior while preserving a callable interface.}
\paragraph{FC-1629 [hard]} State the main tradeoff and production pitfall for Decorators.
\noindent\textbf{Answer.} Tradeoff: Cross-cutting concerns become reusable, but hidden control flow complicates debugging. Pitfall: Omitting functools.wraps breaks metadata and framework introspection.
\par\textit{Tradeoff: Cross-cutting concerns become reusable, but hidden control flow complicates debugging. Pitfall: Omitting functools.wraps breaks metadata and framework introspection.}
\paragraph{LA-1630 [hard]} Explain Decorators from first principles, then show how you would evaluate and operate it in production.
\noindent\textbf{Answer.} Start with the mechanism: Decorators wrap or replace functions or classes to add behavior while preserving a callable interface. Make the tradeoff explicit: Cross-cutting concerns become reusable, but hidden control flow complicates debugging. Define workload-specific offline metrics and a latency/cost budget, test important slices, instrument the critical path, and create a rollback criterion. The production review must explicitly guard against this failure: Omitting functools.wraps breaks metadata and framework introspection.
\par\textit{A strong answer connects mechanism, measurable tradeoffs, failure isolation, observability, and rollback rather than listing product features.}
\paragraph{MCQ-1631 [medium]} Which statement best describes Context managers?
\begin{enumerate}[label=\Alph*.]
\item Dataclasses reduce class boilerplate; Pydantic adds parsing, validation, serialization, and schema generation.
\item Context managers guarantee paired setup and cleanup through with, \_\_enter\_\_/\_\_exit\_\_, or contextlib.
\item Type hints support static analysis; Protocol defines structural interfaces without inheritance.
\item ThreadPoolExecutor and ProcessPoolExecutor expose a common future-based API for concurrent callables.
\end{enumerate}
\noindent\textbf{Answer.} Context managers guarantee paired setup and cleanup through with, \_\_enter\_\_/\_\_exit\_\_, or contextlib.
\par\textit{Context managers guarantee paired setup and cleanup through with, \_\_enter\_\_/\_\_exit\_\_, or contextlib. Tradeoff: Resource safety improves and exceptions are localized. Watch for: Swallowing exceptions unintentionally in \_\_exit\_\_ hides failures.}
\paragraph{MCQ-1632 [hard]} What is the central engineering tradeoff of Context managers?
\begin{enumerate}[label=\Alph*.]
\item Refactoring and contracts improve, while runtime behavior is unchanged unless validated.
\item Typed models improve boundaries but validation can add runtime cost.
\item Integration is simple, but cancellation and process serialization have limits.
\item Resource safety improves and exceptions are localized.
\end{enumerate}
\noindent\textbf{Answer.} Resource safety improves and exceptions are localized.
\par\textit{Context managers guarantee paired setup and cleanup through with, \_\_enter\_\_/\_\_exit\_\_, or contextlib. Tradeoff: Resource safety improves and exceptions are localized. Watch for: Swallowing exceptions unintentionally in \_\_exit\_\_ hides failures.}
\paragraph{MCQ-1633 [hard]} Which failure mode should an interviewer expect you to identify for Context managers?
\begin{enumerate}[label=\Alph*.]
\item Submitting unbounded tasks creates memory pressure and queue latency.
\item Assuming type annotations enforce production inputs creates vulnerabilities.
\item Treating parsed external input as semantically authorized is unsafe.
\item Swallowing exceptions unintentionally in \_\_exit\_\_ hides failures.
\end{enumerate}
\noindent\textbf{Answer.} Swallowing exceptions unintentionally in \_\_exit\_\_ hides failures.
\par\textit{Context managers guarantee paired setup and cleanup through with, \_\_enter\_\_/\_\_exit\_\_, or contextlib. Tradeoff: Resource safety improves and exceptions are localized. Watch for: Swallowing exceptions unintentionally in \_\_exit\_\_ hides failures.}
\paragraph{MCQ-1634 [medium]} A design review is specifically concerned with this risk: Swallowing exceptions unintentionally in \_\_exit\_\_ hides failures. Which topic should be reviewed first?
\begin{enumerate}[label=\Alph*.]
\item Concurrent futures
\item Dataclasses and Pydantic
\item Context managers
\item Type hints and protocols
\end{enumerate}
\noindent\textbf{Answer.} Context managers
\par\textit{Context managers guarantee paired setup and cleanup through with, \_\_enter\_\_/\_\_exit\_\_, or contextlib. Tradeoff: Resource safety improves and exceptions are localized. Watch for: Swallowing exceptions unintentionally in \_\_exit\_\_ hides failures.}
\paragraph{MCQ-1635 [hard]} Which design note most accurately balances the benefits and costs of Context managers?
\begin{enumerate}[label=\Alph*.]
\item Typed models improve boundaries but validation can add runtime cost.
\item Refactoring and contracts improve, while runtime behavior is unchanged unless validated.
\item Integration is simple, but cancellation and process serialization have limits.
\item Resource safety improves and exceptions are localized.
\end{enumerate}
\noindent\textbf{Answer.} Resource safety improves and exceptions are localized.
\par\textit{Context managers guarantee paired setup and cleanup through with, \_\_enter\_\_/\_\_exit\_\_, or contextlib. Tradeoff: Resource safety improves and exceptions are localized. Watch for: Swallowing exceptions unintentionally in \_\_exit\_\_ hides failures.}
\paragraph{MCQ-1636 [medium]} Which explanation would be strongest in a system-design interview about Context managers?
\begin{enumerate}[label=\Alph*.]
\item Context managers guarantee paired setup and cleanup through with, \_\_enter\_\_/\_\_exit\_\_, or contextlib.
\item Dataclasses reduce class boilerplate; Pydantic adds parsing, validation, serialization, and schema generation.
\item ThreadPoolExecutor and ProcessPoolExecutor expose a common future-based API for concurrent callables.
\item Type hints support static analysis; Protocol defines structural interfaces without inheritance.
\end{enumerate}
\noindent\textbf{Answer.} Context managers guarantee paired setup and cleanup through with, \_\_enter\_\_/\_\_exit\_\_, or contextlib.
\par\textit{Context managers guarantee paired setup and cleanup through with, \_\_enter\_\_/\_\_exit\_\_, or contextlib. Tradeoff: Resource safety improves and exceptions are localized. Watch for: Swallowing exceptions unintentionally in \_\_exit\_\_ hides failures.}
\paragraph{MCQ-1637 [hard]} Which production warning is most directly associated with Context managers?
\begin{enumerate}[label=\Alph*.]
\item Submitting unbounded tasks creates memory pressure and queue latency.
\item Treating parsed external input as semantically authorized is unsafe.
\item Swallowing exceptions unintentionally in \_\_exit\_\_ hides failures.
\item Assuming type annotations enforce production inputs creates vulnerabilities.
\end{enumerate}
\noindent\textbf{Answer.} Swallowing exceptions unintentionally in \_\_exit\_\_ hides failures.
\par\textit{Context managers guarantee paired setup and cleanup through with, \_\_enter\_\_/\_\_exit\_\_, or contextlib. Tradeoff: Resource safety improves and exceptions are localized. Watch for: Swallowing exceptions unintentionally in \_\_exit\_\_ hides failures.}
\paragraph{FC-1638 [medium]} Define Context managers in interview-ready language.
\noindent\textbf{Answer.} Context managers guarantee paired setup and cleanup through with, \_\_enter\_\_/\_\_exit\_\_, or contextlib.
\par\textit{Context managers guarantee paired setup and cleanup through with, \_\_enter\_\_/\_\_exit\_\_, or contextlib.}
\paragraph{FC-1639 [hard]} State the main tradeoff and production pitfall for Context managers.
\noindent\textbf{Answer.} Tradeoff: Resource safety improves and exceptions are localized. Pitfall: Swallowing exceptions unintentionally in \_\_exit\_\_ hides failures.
\par\textit{Tradeoff: Resource safety improves and exceptions are localized. Pitfall: Swallowing exceptions unintentionally in \_\_exit\_\_ hides failures.}
\paragraph{LA-1640 [hard]} Explain Context managers from first principles, then show how you would evaluate and operate it in production.
\noindent\textbf{Answer.} Start with the mechanism: Context managers guarantee paired setup and cleanup through with, \_\_enter\_\_/\_\_exit\_\_, or contextlib. Make the tradeoff explicit: Resource safety improves and exceptions are localized. Define workload-specific offline metrics and a latency/cost budget, test important slices, instrument the critical path, and create a rollback criterion. The production review must explicitly guard against this failure: Swallowing exceptions unintentionally in \_\_exit\_\_ hides failures.
\par\textit{A strong answer connects mechanism, measurable tradeoffs, failure isolation, observability, and rollback rather than listing product features.}
\paragraph{MCQ-1641 [medium]} Which statement best describes Generators?
\begin{enumerate}[label=\Alph*.]
\item Generators suspend and resume execution around yield, enabling lazy iteration and streaming pipelines.
\item Asyncio cooperatively schedules coroutines on an event loop around awaitable I/O.
\item Decorators wrap or replace functions or classes to add behavior while preserving a callable interface.
\item CPython threads share memory; the GIL limits simultaneous Python bytecode but releases around many blocking operations.
\end{enumerate}
\noindent\textbf{Answer.} Generators suspend and resume execution around yield, enabling lazy iteration and streaming pipelines.
\par\textit{Generators suspend and resume execution around yield, enabling lazy iteration and streaming pipelines. Tradeoff: Memory use falls, while one-shot iteration and cleanup require care. Watch for: Reusing an exhausted generator returns no data without warning.}
\paragraph{MCQ-1642 [hard]} What is the central engineering tradeoff of Generators?
\begin{enumerate}[label=\Alph*.]
\item High connection concurrency is efficient, while blocking functions freeze the loop.
\item Memory use falls, while one-shot iteration and cleanup require care.
\item Threads suit I/O-bound work and shared state, not pure Python CPU scaling.
\item Cross-cutting concerns become reusable, but hidden control flow complicates debugging.
\end{enumerate}
\noindent\textbf{Answer.} Memory use falls, while one-shot iteration and cleanup require care.
\par\textit{Generators suspend and resume execution around yield, enabling lazy iteration and streaming pipelines. Tradeoff: Memory use falls, while one-shot iteration and cleanup require care. Watch for: Reusing an exhausted generator returns no data without warning.}
\paragraph{MCQ-1643 [hard]} Which failure mode should an interviewer expect you to identify for Generators?
\begin{enumerate}[label=\Alph*.]
\item Creating coroutines without awaiting them silently drops work.
\item Claiming threads are always useless ignores network and native-extension concurrency.
\item Omitting functools.wraps breaks metadata and framework introspection.
\item Reusing an exhausted generator returns no data without warning.
\end{enumerate}
\noindent\textbf{Answer.} Reusing an exhausted generator returns no data without warning.
\par\textit{Generators suspend and resume execution around yield, enabling lazy iteration and streaming pipelines. Tradeoff: Memory use falls, while one-shot iteration and cleanup require care. Watch for: Reusing an exhausted generator returns no data without warning.}
\paragraph{MCQ-1644 [medium]} A design review is specifically concerned with this risk: Reusing an exhausted generator returns no data without warning. Which topic should be reviewed first?
\begin{enumerate}[label=\Alph*.]
\item Decorators
\item Generators
\item Python threading and the GIL
\item Asyncio
\end{enumerate}
\noindent\textbf{Answer.} Generators
\par\textit{Generators suspend and resume execution around yield, enabling lazy iteration and streaming pipelines. Tradeoff: Memory use falls, while one-shot iteration and cleanup require care. Watch for: Reusing an exhausted generator returns no data without warning.}
\paragraph{MCQ-1645 [hard]} Which design note most accurately balances the benefits and costs of Generators?
\begin{enumerate}[label=\Alph*.]
\item Memory use falls, while one-shot iteration and cleanup require care.
\item High connection concurrency is efficient, while blocking functions freeze the loop.
\item Cross-cutting concerns become reusable, but hidden control flow complicates debugging.
\item Threads suit I/O-bound work and shared state, not pure Python CPU scaling.
\end{enumerate}
\noindent\textbf{Answer.} Memory use falls, while one-shot iteration and cleanup require care.
\par\textit{Generators suspend and resume execution around yield, enabling lazy iteration and streaming pipelines. Tradeoff: Memory use falls, while one-shot iteration and cleanup require care. Watch for: Reusing an exhausted generator returns no data without warning.}
\paragraph{MCQ-1646 [medium]} Which explanation would be strongest in a system-design interview about Generators?
\begin{enumerate}[label=\Alph*.]
\item Decorators wrap or replace functions or classes to add behavior while preserving a callable interface.
\item Generators suspend and resume execution around yield, enabling lazy iteration and streaming pipelines.
\item Asyncio cooperatively schedules coroutines on an event loop around awaitable I/O.
\item CPython threads share memory; the GIL limits simultaneous Python bytecode but releases around many blocking operations.
\end{enumerate}
\noindent\textbf{Answer.} Generators suspend and resume execution around yield, enabling lazy iteration and streaming pipelines.
\par\textit{Generators suspend and resume execution around yield, enabling lazy iteration and streaming pipelines. Tradeoff: Memory use falls, while one-shot iteration and cleanup require care. Watch for: Reusing an exhausted generator returns no data without warning.}
\paragraph{MCQ-1647 [hard]} Which production warning is most directly associated with Generators?
\begin{enumerate}[label=\Alph*.]
\item Creating coroutines without awaiting them silently drops work.
\item Reusing an exhausted generator returns no data without warning.
\item Omitting functools.wraps breaks metadata and framework introspection.
\item Claiming threads are always useless ignores network and native-extension concurrency.
\end{enumerate}
\noindent\textbf{Answer.} Reusing an exhausted generator returns no data without warning.
\par\textit{Generators suspend and resume execution around yield, enabling lazy iteration and streaming pipelines. Tradeoff: Memory use falls, while one-shot iteration and cleanup require care. Watch for: Reusing an exhausted generator returns no data without warning.}
\paragraph{FC-1648 [medium]} Define Generators in interview-ready language.
\noindent\textbf{Answer.} Generators suspend and resume execution around yield, enabling lazy iteration and streaming pipelines.
\par\textit{Generators suspend and resume execution around yield, enabling lazy iteration and streaming pipelines.}
\paragraph{FC-1649 [hard]} State the main tradeoff and production pitfall for Generators.
\noindent\textbf{Answer.} Tradeoff: Memory use falls, while one-shot iteration and cleanup require care. Pitfall: Reusing an exhausted generator returns no data without warning.
\par\textit{Tradeoff: Memory use falls, while one-shot iteration and cleanup require care. Pitfall: Reusing an exhausted generator returns no data without warning.}
\paragraph{LA-1650 [hard]} Explain Generators from first principles, then show how you would evaluate and operate it in production.
\noindent\textbf{Answer.} Start with the mechanism: Generators suspend and resume execution around yield, enabling lazy iteration and streaming pipelines. Make the tradeoff explicit: Memory use falls, while one-shot iteration and cleanup require care. Define workload-specific offline metrics and a latency/cost budget, test important slices, instrument the critical path, and create a rollback criterion. The production review must explicitly guard against this failure: Reusing an exhausted generator returns no data without warning.
\par\textit{A strong answer connects mechanism, measurable tradeoffs, failure isolation, observability, and rollback rather than listing product features.}
\paragraph{MCQ-1651 [medium]} Which statement best describes Descriptors and properties?
\begin{enumerate}[label=\Alph*.]
\item Descriptors define attribute access; properties provide managed getter, setter, and deleter behavior.
\item Generators suspend and resume execution around yield, enabling lazy iteration and streaming pipelines.
\item Separate processes bypass the GIL and isolate memory while communicating through serialization or shared memory.
\item Unit, integration, property, contract, and load tests cover different risks; mocks isolate only well-understood boundaries.
\end{enumerate}
\noindent\textbf{Answer.} Descriptors define attribute access; properties provide managed getter, setter, and deleter behavior.
\par\textit{Descriptors define attribute access; properties provide managed getter, setter, and deleter behavior. Tradeoff: Validation and lazy access fit attribute syntax but may hide expensive operations. Watch for: Performing network I/O in a property surprises callers and tooling.}
\paragraph{MCQ-1652 [hard]} What is the central engineering tradeoff of Descriptors and properties?
\begin{enumerate}[label=\Alph*.]
\item Memory use falls, while one-shot iteration and cleanup require care.
\item Validation and lazy access fit attribute syntax but may hide expensive operations.
\item CPU parallelism improves at startup, memory, and coordination cost.
\item Fast tests aid iteration, while excessive mocking tests implementation rather than behavior.
\end{enumerate}
\noindent\textbf{Answer.} Validation and lazy access fit attribute syntax but may hide expensive operations.
\par\textit{Descriptors define attribute access; properties provide managed getter, setter, and deleter behavior. Tradeoff: Validation and lazy access fit attribute syntax but may hide expensive operations. Watch for: Performing network I/O in a property surprises callers and tooling.}
\paragraph{MCQ-1653 [hard]} Which failure mode should an interviewer expect you to identify for Descriptors and properties?
\begin{enumerate}[label=\Alph*.]
\item Passing large models to spawned workers duplicates memory unexpectedly.
\item Reusing an exhausted generator returns no data without warning.
\item Mocking the model, database, and network in the same test proves little about integration.
\item Performing network I/O in a property surprises callers and tooling.
\end{enumerate}
\noindent\textbf{Answer.} Performing network I/O in a property surprises callers and tooling.
\par\textit{Descriptors define attribute access; properties provide managed getter, setter, and deleter behavior. Tradeoff: Validation and lazy access fit attribute syntax but may hide expensive operations. Watch for: Performing network I/O in a property surprises callers and tooling.}
\paragraph{MCQ-1654 [medium]} A design review is specifically concerned with this risk: Performing network I/O in a property surprises callers and tooling. Which topic should be reviewed first?
\begin{enumerate}[label=\Alph*.]
\item Generators
\item Testing and mocking
\item Descriptors and properties
\item Multiprocessing
\end{enumerate}
\noindent\textbf{Answer.} Descriptors and properties
\par\textit{Descriptors define attribute access; properties provide managed getter, setter, and deleter behavior. Tradeoff: Validation and lazy access fit attribute syntax but may hide expensive operations. Watch for: Performing network I/O in a property surprises callers and tooling.}
\paragraph{MCQ-1655 [hard]} Which design note most accurately balances the benefits and costs of Descriptors and properties?
\begin{enumerate}[label=\Alph*.]
\item Validation and lazy access fit attribute syntax but may hide expensive operations.
\item Fast tests aid iteration, while excessive mocking tests implementation rather than behavior.
\item Memory use falls, while one-shot iteration and cleanup require care.
\item CPU parallelism improves at startup, memory, and coordination cost.
\end{enumerate}
\noindent\textbf{Answer.} Validation and lazy access fit attribute syntax but may hide expensive operations.
\par\textit{Descriptors define attribute access; properties provide managed getter, setter, and deleter behavior. Tradeoff: Validation and lazy access fit attribute syntax but may hide expensive operations. Watch for: Performing network I/O in a property surprises callers and tooling.}
\paragraph{MCQ-1656 [medium]} Which explanation would be strongest in a system-design interview about Descriptors and properties?
\begin{enumerate}[label=\Alph*.]
\item Descriptors define attribute access; properties provide managed getter, setter, and deleter behavior.
\item Separate processes bypass the GIL and isolate memory while communicating through serialization or shared memory.
\item Generators suspend and resume execution around yield, enabling lazy iteration and streaming pipelines.
\item Unit, integration, property, contract, and load tests cover different risks; mocks isolate only well-understood boundaries.
\end{enumerate}
\noindent\textbf{Answer.} Descriptors define attribute access; properties provide managed getter, setter, and deleter behavior.
\par\textit{Descriptors define attribute access; properties provide managed getter, setter, and deleter behavior. Tradeoff: Validation and lazy access fit attribute syntax but may hide expensive operations. Watch for: Performing network I/O in a property surprises callers and tooling.}
\paragraph{MCQ-1657 [hard]} Which production warning is most directly associated with Descriptors and properties?
\begin{enumerate}[label=\Alph*.]
\item Mocking the model, database, and network in the same test proves little about integration.
\item Performing network I/O in a property surprises callers and tooling.
\item Passing large models to spawned workers duplicates memory unexpectedly.
\item Reusing an exhausted generator returns no data without warning.
\end{enumerate}
\noindent\textbf{Answer.} Performing network I/O in a property surprises callers and tooling.
\par\textit{Descriptors define attribute access; properties provide managed getter, setter, and deleter behavior. Tradeoff: Validation and lazy access fit attribute syntax but may hide expensive operations. Watch for: Performing network I/O in a property surprises callers and tooling.}
\paragraph{FC-1658 [medium]} Define Descriptors and properties in interview-ready language.
\noindent\textbf{Answer.} Descriptors define attribute access; properties provide managed getter, setter, and deleter behavior.
\par\textit{Descriptors define attribute access; properties provide managed getter, setter, and deleter behavior.}
\paragraph{FC-1659 [hard]} State the main tradeoff and production pitfall for Descriptors and properties.
\noindent\textbf{Answer.} Tradeoff: Validation and lazy access fit attribute syntax but may hide expensive operations. Pitfall: Performing network I/O in a property surprises callers and tooling.
\par\textit{Tradeoff: Validation and lazy access fit attribute syntax but may hide expensive operations. Pitfall: Performing network I/O in a property surprises callers and tooling.}
\paragraph{LA-1660 [hard]} Explain Descriptors and properties from first principles, then show how you would evaluate and operate it in production.
\noindent\textbf{Answer.} Start with the mechanism: Descriptors define attribute access; properties provide managed getter, setter, and deleter behavior. Make the tradeoff explicit: Validation and lazy access fit attribute syntax but may hide expensive operations. Define workload-specific offline metrics and a latency/cost budget, test important slices, instrument the critical path, and create a rollback criterion. The production review must explicitly guard against this failure: Performing network I/O in a property surprises callers and tooling.
\par\textit{A strong answer connects mechanism, measurable tradeoffs, failure isolation, observability, and rollback rather than listing product features.}
\paragraph{MCQ-1661 [medium]} Which statement best describes Dataclasses and Pydantic?
\begin{enumerate}[label=\Alph*.]
\item Unit, integration, property, contract, and load tests cover different risks; mocks isolate only well-understood boundaries.
\item Dataclasses reduce class boilerplate; Pydantic adds parsing, validation, serialization, and schema generation.
\item CPython combines reference counting with cyclic garbage collection; object graphs and native buffers affect real memory.
\item Asyncio cooperatively schedules coroutines on an event loop around awaitable I/O.
\end{enumerate}
\noindent\textbf{Answer.} Dataclasses reduce class boilerplate; Pydantic adds parsing, validation, serialization, and schema generation.
\par\textit{Dataclasses reduce class boilerplate; Pydantic adds parsing, validation, serialization, and schema generation. Tradeoff: Typed models improve boundaries but validation can add runtime cost. Watch for: Treating parsed external input as semantically authorized is unsafe.}
\paragraph{MCQ-1662 [hard]} What is the central engineering tradeoff of Dataclasses and Pydantic?
\begin{enumerate}[label=\Alph*.]
\item Fast tests aid iteration, while excessive mocking tests implementation rather than behavior.
\item Typed models improve boundaries but validation can add runtime cost.
\item Automatic cleanup is convenient, while cycles and caches can retain large objects.
\item High connection concurrency is efficient, while blocking functions freeze the loop.
\end{enumerate}
\noindent\textbf{Answer.} Typed models improve boundaries but validation can add runtime cost.
\par\textit{Dataclasses reduce class boilerplate; Pydantic adds parsing, validation, serialization, and schema generation. Tradeoff: Typed models improve boundaries but validation can add runtime cost. Watch for: Treating parsed external input as semantically authorized is unsafe.}
\paragraph{MCQ-1663 [hard]} Which failure mode should an interviewer expect you to identify for Dataclasses and Pydantic?
\begin{enumerate}[label=\Alph*.]
\item Creating coroutines without awaiting them silently drops work.
\item Mocking the model, database, and network in the same test proves little about integration.
\item Watching only Python heap misses tensors allocated in native or GPU memory.
\item Treating parsed external input as semantically authorized is unsafe.
\end{enumerate}
\noindent\textbf{Answer.} Treating parsed external input as semantically authorized is unsafe.
\par\textit{Dataclasses reduce class boilerplate; Pydantic adds parsing, validation, serialization, and schema generation. Tradeoff: Typed models improve boundaries but validation can add runtime cost. Watch for: Treating parsed external input as semantically authorized is unsafe.}
\paragraph{MCQ-1664 [medium]} A design review is specifically concerned with this risk: Treating parsed external input as semantically authorized is unsafe. Which topic should be reviewed first?
\begin{enumerate}[label=\Alph*.]
\item Memory management
\item Asyncio
\item Testing and mocking
\item Dataclasses and Pydantic
\end{enumerate}
\noindent\textbf{Answer.} Dataclasses and Pydantic
\par\textit{Dataclasses reduce class boilerplate; Pydantic adds parsing, validation, serialization, and schema generation. Tradeoff: Typed models improve boundaries but validation can add runtime cost. Watch for: Treating parsed external input as semantically authorized is unsafe.}
\paragraph{MCQ-1665 [hard]} Which design note most accurately balances the benefits and costs of Dataclasses and Pydantic?
\begin{enumerate}[label=\Alph*.]
\item High connection concurrency is efficient, while blocking functions freeze the loop.
\item Typed models improve boundaries but validation can add runtime cost.
\item Automatic cleanup is convenient, while cycles and caches can retain large objects.
\item Fast tests aid iteration, while excessive mocking tests implementation rather than behavior.
\end{enumerate}
\noindent\textbf{Answer.} Typed models improve boundaries but validation can add runtime cost.
\par\textit{Dataclasses reduce class boilerplate; Pydantic adds parsing, validation, serialization, and schema generation. Tradeoff: Typed models improve boundaries but validation can add runtime cost. Watch for: Treating parsed external input as semantically authorized is unsafe.}
\paragraph{MCQ-1666 [medium]} Which explanation would be strongest in a system-design interview about Dataclasses and Pydantic?
\begin{enumerate}[label=\Alph*.]
\item Asyncio cooperatively schedules coroutines on an event loop around awaitable I/O.
\item Dataclasses reduce class boilerplate; Pydantic adds parsing, validation, serialization, and schema generation.
\item Unit, integration, property, contract, and load tests cover different risks; mocks isolate only well-understood boundaries.
\item CPython combines reference counting with cyclic garbage collection; object graphs and native buffers affect real memory.
\end{enumerate}
\noindent\textbf{Answer.} Dataclasses reduce class boilerplate; Pydantic adds parsing, validation, serialization, and schema generation.
\par\textit{Dataclasses reduce class boilerplate; Pydantic adds parsing, validation, serialization, and schema generation. Tradeoff: Typed models improve boundaries but validation can add runtime cost. Watch for: Treating parsed external input as semantically authorized is unsafe.}
\paragraph{MCQ-1667 [hard]} Which production warning is most directly associated with Dataclasses and Pydantic?
\begin{enumerate}[label=\Alph*.]
\item Treating parsed external input as semantically authorized is unsafe.
\item Mocking the model, database, and network in the same test proves little about integration.
\item Creating coroutines without awaiting them silently drops work.
\item Watching only Python heap misses tensors allocated in native or GPU memory.
\end{enumerate}
\noindent\textbf{Answer.} Treating parsed external input as semantically authorized is unsafe.
\par\textit{Dataclasses reduce class boilerplate; Pydantic adds parsing, validation, serialization, and schema generation. Tradeoff: Typed models improve boundaries but validation can add runtime cost. Watch for: Treating parsed external input as semantically authorized is unsafe.}
\paragraph{FC-1668 [medium]} Define Dataclasses and Pydantic in interview-ready language.
\noindent\textbf{Answer.} Dataclasses reduce class boilerplate; Pydantic adds parsing, validation, serialization, and schema generation.
\par\textit{Dataclasses reduce class boilerplate; Pydantic adds parsing, validation, serialization, and schema generation.}
\paragraph{FC-1669 [hard]} State the main tradeoff and production pitfall for Dataclasses and Pydantic.
\noindent\textbf{Answer.} Tradeoff: Typed models improve boundaries but validation can add runtime cost. Pitfall: Treating parsed external input as semantically authorized is unsafe.
\par\textit{Tradeoff: Typed models improve boundaries but validation can add runtime cost. Pitfall: Treating parsed external input as semantically authorized is unsafe.}
\paragraph{LA-1670 [hard]} Explain Dataclasses and Pydantic from first principles, then show how you would evaluate and operate it in production.
\noindent\textbf{Answer.} Start with the mechanism: Dataclasses reduce class boilerplate; Pydantic adds parsing, validation, serialization, and schema generation. Make the tradeoff explicit: Typed models improve boundaries but validation can add runtime cost. Define workload-specific offline metrics and a latency/cost budget, test important slices, instrument the critical path, and create a rollback criterion. The production review must explicitly guard against this failure: Treating parsed external input as semantically authorized is unsafe.
\par\textit{A strong answer connects mechanism, measurable tradeoffs, failure isolation, observability, and rollback rather than listing product features.}
\paragraph{MCQ-1671 [medium]} Which statement best describes Type hints and protocols?
\begin{enumerate}[label=\Alph*.]
\item Asyncio cooperatively schedules coroutines on an event loop around awaitable I/O.
\item Dataclasses reduce class boilerplate; Pydantic adds parsing, validation, serialization, and schema generation.
\item Exceptions separate error propagation from normal return values and can preserve causal chains.
\item Type hints support static analysis; Protocol defines structural interfaces without inheritance.
\end{enumerate}
\noindent\textbf{Answer.} Type hints support static analysis; Protocol defines structural interfaces without inheritance.
\par\textit{Type hints support static analysis; Protocol defines structural interfaces without inheritance. Tradeoff: Refactoring and contracts improve, while runtime behavior is unchanged unless validated. Watch for: Assuming type annotations enforce production inputs creates vulnerabilities.}
\paragraph{MCQ-1672 [hard]} What is the central engineering tradeoff of Type hints and protocols?
\begin{enumerate}[label=\Alph*.]
